% Generated mechanically from the frozen exercises.tex target.
% Do not edit this generated file; edit the bound locale source instead.

% OK, start here.
%

\setcounter{chapter}{110}
\chapter{练习}



\phantomsection
\label{exercises-section-phantom}



\section{代数}
\label{exercises-section-algebra}

\noindent
本节首先收录若干杂题。

\begin{exercise}
\label{exercises-exercise-isomorphism-localizations}
设 $A$ 为环，${\mathfrak m}$ 为极大理想。在 $A[X]$ 中令
$\tilde {\mathfrak m}_1 = ({\mathfrak m}, X)$ 且
$\tilde {\mathfrak m}_2 = ({\mathfrak m}, X-1)$。证明
$$
A[X]_{\tilde {\mathfrak m}_1} \cong A[X]_{\tilde {\mathfrak m}_2}.
$$
\end{exercise}

\begin{exercise}
\label{exercises-exercise-coherent}
% P12-ERR-0415
求一个非 Noether 环 $R$ 的例子，使得 $R$ 的每个有限生成理想作为
$R$-模都是有限表示的。（若一个环具有后一性质，则称它是{\it 凝聚环}。）
\end{exercise}

\begin{exercise}
\label{exercises-exercise-flat-ideals-pid}
设 $(A, {\mathfrak m}, k)$ 为 Noether 局部环。对任意有限 $A$-模 $M$，
定义 $r(M)$ 为 $M$ 作为 $A$-模的最少生成元个数。由 NAK，此数等于
$\dim_k M/{\mathfrak m} M = \dim_k M \otimes_A k$。
\begin{enumerate}
\item 证明 $r(M \otimes_A N) = r(M)r(N)$。
\item 设 $I\subset A $ 为满足 $r(I) > 1$ 的理想。证明
$r(I^2) < r(I)^2$。
\item 推出：若 $A$ 中每个理想都是平坦模，则 $A$ 是主理想整环
（或域）。
\end{enumerate}
\end{exercise}

\begin{exercise}
\label{exercises-exercise-not-isomorphic}
设 $k$ 为域。证明下列每一对 $k$-代数不同构：
\begin{enumerate}
\item $k[x_1, \ldots, x_n]$ 与 $k[x_1, \ldots, x_{n + 1}]$，
其中 $n\geq 1$ 任意；
\item $k[a, b, c, d, e, f]/(ab + cd + ef)$ 与
$k[x_1, \ldots, x_n]$，其中 $n = 5$；
\item $k[a, b, c, d, e, f]/(ab + cd + ef)$ 与
$k[x_1, \ldots, x_n]$，其中 $n = 6$。
\end{enumerate}
\end{exercise}

\begin{remark}
\label{exercises-remark-simple-geometric}
当然，本练习的意图是在每种情形中找到简单论证，而不是应用某个“大”定理。
不过，以一般原理为指引仍然很有益。
\end{remark}

\begin{exercise}
\label{exercises-exercise-silly}
代数。（这是个很初等且应当容易的问题。）
\begin{enumerate}
\item 给出一个环 $A$ 及一个不分裂的 $A$-模短正合列的例子
$$
0 \to M_1 \to M_2 \to M_3 \to 0.
$$
\item 给出一个如上的不分裂 $A$-模正合列，以及一个忠实平坦映射
$A \to B$，使得
$$
0 \to M_1\otimes_A B \to M_2\otimes_A B \to M_3\otimes_A B \to 0.
$$
作为 $B$-模正合列是分裂的。
\end{enumerate}
\end{exercise}

\begin{exercise}
\label{exercises-exercise-field-kummer}
设域 $k$ 含有一个本原 $n$ 次单位根 $\zeta$。这意味着
$\zeta^n = 1$，但 $\zeta^m\not = 1$ 对 $0 < m < n$ 成立。
\begin{enumerate}
\item 证明 $k$ 的特征与 $n$ 互素。
\item 设 $a \in k$ 是 $k$ 中一个元素。
% P12-ERR-0416
若 $d$ 是 $n$ 的任意满足 $n \geq d > 1$ 的因子，则该元素在 $k$ 中
不是 $d$ 次幂。证明 $k[x]/(x^n-a)$ 是域。（提示：考虑
$x^n-a$ 的一个分裂域，并使用 Galois 理论。）
\end{enumerate}
\end{exercise}

\begin{exercise}
\label{exercises-exercise-valuation}
设 $\nu : k[x]\setminus \{0\}  \to {\mathbf Z}$ 为满足下列性质的映射：
只要 $f$、$g$ 非零，就有 $\nu(fg) = \nu(f) + \nu(g)$；只要
$f$、$g$、$f + g$ 非零，就有
% P12-ERR-0417
$\nu(f + g) \geq min(\nu(f), \nu(g))$；并且
$\nu(c) = 0$ 对所有 $c\in k^*$ 都成立。
\begin{enumerate}
\item 证明：若 $f$、$g$ 和 $f + g$ 均非零，且
$\nu(f) \not = \nu(g)$，则
$\nu(f + g) = min(\nu(f), \nu(g))$。
\item 证明：若 $f = \sum a_i x^i$ 且 $f\not = 0$，则
$\nu(f) \geq min(\{i\nu(x)\}_{a_i\not = 0})$。何时等号成立？
\item 证明：若 $\nu$ 取到负值，则
$\nu(f) = -n \deg(f)$ 对某个 $n\in {\mathbf N}$ 成立。
\item 设 $\nu(x) \geq 0$。证明
$\{f \mid f = 0, \ or\ \nu(f) > 0\}$ 是 $k[x]$ 的素理想。
\item 描述所有可能的 $\nu$。
\end{enumerate}
\end{exercise}


\noindent
设 $A$ 为环。若元素 $e \in A$ 满足 $e^2 = e$，则称它为{\it 幂等元}。
元素 $1$ 和 $0$ 总是幂等元。
% P12-ERR-0418
{\it 非平凡幂等元}是一个不等于零的幂等元。若两个幂等元
$e, e' \in A$ 满足 $ee' = 0$，则称它们{\it 正交}。

\begin{exercise}
\label{exercises-exercise-product}
设 $A$ 为环。证明：$A$ 是两个非零环的直积，当且仅当 $A$
具有非平凡幂等元。
\end{exercise}

\begin{exercise}
\label{exercises-exercise-lift-idempotents}
设 $A$ 为环，$I \subset A$ 为局部幂零理想。证明映射
$A \to A/I$ 在幂等元上诱导双射。（提示：先证明 $I$ 为幂零理想时的结论，
这可能更容易。然后用“绝对 Noether 约化”化归到幂零情形。）
\end{exercise}







\section{余极限}
\label{exercises-section-colimits}


\begin{definition}
\label{exercises-definition-directed-poset}
{\it 有向集}是一个非空集合 $I$，其上赋有预序 $\leq$，并且对任意一对
$i, j \in I$，存在 $k \in I$ 使得 $i \leq k$ 且 $j \leq k$。
$I$ 上的一个{\it 环系}由每个 $i \in I$ 所对应的环 $A_i$，以及每当
$i \leq j$ 时的环映射 $\varphi_{ij} : A_i \to A_j$ 组成，并要求：
只要 $i \leq j \leq k$，复合 $A_i \to A_j \to A_k$ 就等于
$A_i \to A_k$。
\end{definition}

\noindent
群系、固定环上的模系、域上的向量空间系等可类似定义。

\begin{exercise}
\label{exercises-exercise-directed-colimit}
设 $I$ 为有向集，$(A_i, \varphi_{ij})$ 为 $I$ 上的环系。
证明存在环 $A$ 及映射 $\varphi_i : A_i \to A$，使得对所有
$i \leq j$ 都有 $\varphi_j \circ \varphi_{ij} = \varphi_i$，并具有
下列泛性质：给定任意环 $B$ 及映射 $\psi_i : A_i \to B$，若对所有
$i \leq j$ 都有 $\psi_j \circ \varphi_{ij} = \psi_i$，则存在唯一的
环映射 $\psi : A \to B$ 使得 $\psi_i = \psi \circ \varphi_i$。
\end{exercise}

\begin{definition}
\label{exercises-definition-colimit}
练习 \ref{exercises-exercise-directed-colimit} 中构造的环 $A$ 称为该系统的
{\it 余极限}，记作 $\colim A_i$。
\end{definition}

\begin{exercise}
\label{exercises-exercise-prime-in-colimit}
% P12-ERR-0419
设 $(I, \geq)$ 为有向集，$(A_i, \varphi_{ij})$ 为 $I$ 上的环系，
其余极限为 $A$。证明存在双射
$$
\Spec(A) = \{(\mathfrak p_i)_{i \in I} \mid
\mathfrak p_i \subset A_i \text{ 且 }
\mathfrak p_i = \varphi_{ij}^{-1}(\mathfrak p_j)\ \forall i \leq j\}
\subset \prod\nolimits_{i \in I} \Spec(A_i)
$$
等式右侧的集合是各集合 $\Spec(A_i)$ 的极限，记作
$\lim \Spec(A_i)$。
\end{exercise}

\begin{exercise}
\label{exercises-exercise-colimit-surjective}
设 $(I, \geq)$ 为有向集，$(A_i, \varphi_{ij})$ 为 $I$ 上的环系，
其余极限为 $A$。设对所有 $i \leq j$，映射
$\Spec(A_j) \to \Spec(A_i)$ 都是满射。证明对所有 $i$，
$\Spec(A) \to \Spec(A_i)$ 都是满射。（提示：可以尝试使用
Tychonoff 定理；不过，由《代数》引理
\ref{algebra-lemma-in-image} 也可给出一个基本上平凡的直接代数证明。）
\end{exercise}

\begin{exercise}
\label{exercises-exercise-integral-colimit-finite}
设 $A \subset B$ 为整环扩张。证明
$\Spec(B) \to \Spec(A)$ 是满射。
% P12-ERR-0420
使用上面的练习、课堂上已证明的有限环扩张情形，并证明
$B = \colim B_i$ 是有限扩张 $A \subset B_i$ 的有向余极限。
\end{exercise}

\begin{exercise}
\label{exercises-exercise-colimit-tensor}
设 $(I, \geq)$ 为有向集。设 $A$ 为环，
$(N_i, \varphi_{i, i'})$ 为由 $I$ 标号的 $A$-模有向系，并设
$M$ 为另一个 $A$-模。证明
$$
\colim_{i\in I} M \otimes_A N_i\cong
M \otimes_A \Big( \colim_{i\in I} N_i\Big).
$$
\end{exercise}

\begin{definition}
\label{exercises-definition-finite-presentation}
若 $R$-模 $M$ 同构于有限秩自由模之间某个映射
$ R^{\oplus n} \to R^{\oplus m}$ 的余核，则称它在 $R$ 上
{\it 有限表示}。
\end{definition}

\begin{exercise}
\label{exercises-exercise-colimit-modules}
证明任意环上的任意模：
\begin{enumerate}
\item 是其有限生成子模的余极限；
\item 可按某种方式写成有限表示模的余极限。
\end{enumerate}
\end{exercise}





\section{加性范畴与 Abel 范畴}
\label{exercises-section-additive}

\begin{exercise}
\label{exercises-exercise-filtered-vector-spaces}
设 $k$ 为域，$\mathcal{C}$ 为 $k$ 上滤过向量空间所成的范畴。
任意范畴中滤过对象的定义见《同调代数》定义
\ref{homology-definition-filtered}。
\begin{enumerate}
\item 证明这是一个加性范畴（并仔细说明两个对象的直和是什么）。
\item 设 $f : (V, F) \to (W, F)$ 为 $\mathcal{C}$ 的一个态射。
证明 $f$ 有核与余核（并准确说明 $f$ 的核与余核是什么）。
\item 给出 $\mathcal{C}$ 中一个映射的例子，使得典范映射
$\Coim(f) \to \Im(f)$ 不是同构。
\end{enumerate}
\end{exercise}

\begin{exercise}
\label{exercises-exercise-torsion-free}
设 $R$ 为 Noether 整环，$\mathcal{C}$ 为有限生成无挠 $R$-模
所成的范畴。
\begin{enumerate}
\item 证明这是一个加性范畴。
\item 设 $f : N \to M$ 为 $\mathcal{C}$ 的一个态射。证明 $f$
有核与余核（务必准确地定义 $f$ 的核与余核）。
\item 给出一个 Noether 整环 $R$ 以及 $\mathcal{C}$ 中的一个映射，
使得典范映射 $\Coim(f) \to \Im(f)$ 不是同构。
\end{enumerate}
\end{exercise}

\begin{exercise}
\label{exercises-exercise-other}
给出一个加性且具有核与余核的范畴的例子，但它不同于练习
\ref{exercises-exercise-filtered-vector-spaces} 和
\ref{exercises-exercise-torsion-free} 中的范畴。
\end{exercise}





\section{张量积}
\label{exercises-section-tensor-product}

\noindent
张量积在《代数》第 \ref{algebra-section-tensor-product} 节中引入。
设 $R$ 为环，$\text{Mod}_R$ 为 $R$-模范畴。我们称函子
$F : \text{Mod}_R \to \text{Mod}_R$
\begin{enumerate}
\item 是加性的，如果对任意 $R$-模 $M, N$，映射
$F : \Hom_R(M, N) \to \Hom_R(F(M), F(N))$
都是 Abel 群同态；见《同调代数》定义
\ref{homology-definition-preadditive}；
% P12-ERR-0421
\item 是 $R$-线性的，如果对任意 $R$-模 $M, N$，映射
$F : \Hom_R(M, N) \to \Hom_R(F(M), F(N))$ 都是 $R$-线性的；
\item 是右正合的，如果对任意短正合列
$0 \to M_1 \to M_2 \to M_3 \to 0$，序列
$F(M_1) \to F(M_2) \to F(M_3) \to 0$ 都正合；
\item 是左正合的，如果对任意短正合列
$0 \to M_1 \to M_2 \to M_3 \to 0$，序列
$0 \to F(M_1) \to F(M_2) \to F(M_3)$ 都正合；
\item 与直和交换，如果给定集合 $I$ 及 $R$-模 $M_i$，各映射
$F(M_i) \to F(\bigoplus M_i)$ 诱导同构
$\bigoplus F(M_i) = F(\bigoplus M_i)$。
\end{enumerate}

\begin{exercise}
\label{exercises-exercise-characterize-tensor-functor}
设 $R$ 为环，沿用上述记号。
\begin{enumerate}
\item 给出环 $R$ 及加性函子
$F : \text{Mod}_R \to \text{Mod}_R$ 的例子，使得该函子不是
$R$-线性的。
\item 设 $N$ 为 $R$-模。证明函子 $F(M) = M \otimes_R N$
是 $R$-线性的、右正合的，并且与直和交换。
\item 反过来，证明任意函子 $F : \text{Mod}_R \to \text{Mod}_R$，
若它是 $R$-线性的、右正合的且与直和交换，则它具有形式
$F(M) = M \otimes_R N$，其中 $N$ 是某个 $R$-模。
\item 证明：若在 (3) 中去掉 $F$ 与直和交换的假设，则结论不再成立。
\end{enumerate}
\end{exercise}






\section{平坦环映射}
\label{exercises-section-flat}

\begin{exercise}
\label{exercises-exercise-localization-flat}
设 $S$ 为环 $A$ 的乘法子集。
\begin{enumerate}
\item 对 $A$-模 $M$，证明 $S^{-1}M = S^{-1}A \otimes_A M$。
\item 证明 $S^{-1}A$ 在 $A$ 上平坦。
\end{enumerate}
\end{exercise}

\begin{exercise}
\label{exercises-exercise-examples-not-flat}
在下列各情形中，求一个 $A$-模单射 $M_1 \to M_2$，使得
$M_1\otimes N \to M_2 \otimes N$ 不是单射：
\begin{enumerate}
\item $A = k[x, y]$ 且 $N = (x, y) \subset A$。（这里及以下，
$k$ 都是域。）
\item $A = k[x, y]$ 且 $N = A/(x, y)$。
\end{enumerate}
\end{exercise}

\begin{exercise}
\label{exercises-exercise-flat-not-projective}
% P12-ERR-0422
给出环 $A$ 及有限 $A$-模 $M$ 的例子，使得它是平坦的但不是
投射 $A$-模。
\end{exercise}

\begin{remark}
\label{exercises-remark-flat-not-projective}
若 $M$ 在 $A$ 上有限表示且平坦，则 $M$ 在 $A$ 上投射。
因此，你的例子必须涉及一个非 Noether 环 $A$。我知道一个例子，
其中 $A$ 是 ${\mathcal C}^\infty$-函数在 ${\mathbf R}$ 上的环。
\end{remark}

\begin{exercise}
\label{exercises-exercise-flat-not-free-dvr}
求 ${\mathbf Z}_{(2)}$ 上一个平坦但不自由的模。
\end{exercise}

\begin{exercise}
\label{exercises-exercise-flat-deformations}
平坦形变。
\begin{enumerate}
\item 设 $k$ 为域，$k[\epsilon]$ 为对偶数环
$k[\epsilon] = k[x]/(x^2)$，且 $\epsilon = \bar x$。证明：对任意
$k$-代数 $A$，存在平坦 $k[\epsilon]$-代数 $B$，使得
$A$ 同构于 $B/\epsilon B$。
\item 设 $k = {\mathbf F}_p = {\mathbf Z}/p{\mathbf Z}$ 且
$$
A = k[x_1, x_2, x_3, x_4, x_5, x_6]/
(x_1^p, x_2^p, x_3^p, x_4^p, x_5^p, x_6^p).
$$
证明存在平坦 ${\mathbf Z}/p^2{\mathbf Z}$-代数 $B$，使得
$B/pB$ 同构于 $A$。（因此，这里 $p$ 扮演 $\epsilon$ 的角色。）
\item 现在令 $p = 2$，并对
$k = {\mathbf F}_2 = {\mathbf Z}/2{\mathbf Z}$ 及
$$
A = k[x_1, x_2, x_3, x_4, x_5, x_6]/
(x_1^2, x_2^2, x_3^2, x_4^2, x_5^2, x_6^2, x_1x_2 + x_3x_4 + x_5x_6).
$$
考虑同一问题。不过，在这种情形下证明：{\it 不}存在平坦
${\mathbf Z}/4{\mathbf Z}$-代数 $B$ 使得 $B/2B$ 同构于 $A$。
（找出其中的窍门！同一例子对任意特征 $p > 0$ 都成立，只是计算更困难。）
\end{enumerate}
\end{exercise}

\begin{exercise}
\label{exercises-exercise-flat-given-residue-field-extension}
设 $(A, {\mathfrak m}, k)$ 为局部环，$k'/k$ 为有限域扩张。
证明存在局部环之间的平坦局部映射 $A \to B$，使得
${\mathfrak m}_B = {\mathfrak m} B$，并且
$B/{\mathfrak m} B$ 作为 $k$-代数同构于 $k'$。（提示：先处理
$k \subset k'$ 由单个元素生成的情形。）
\end{exercise}

\begin{remark}
\label{exercises-remark-flat-given-residue-field-extension-general}
同一结果对任意域扩张 $K/k$ 都成立。
\end{remark}

\section{环的谱}
\label{exercises-section-spectrum-ring}

\begin{exercise}
\label{exercises-exercise-spec-Z}
计算集合 $\Spec(\mathbf{Z})$ 并描述其拓扑。
\end{exercise}

\begin{exercise}
\label{exercises-exercise-basis-opens-standard}
设 $A$ 为任意环。对 $f\in A$，定义
$D(f):= \{\mathfrak p \subset A \mid f \not \in \mathfrak p\}$。
证明开子集 $D(f)$ 构成 $\Spec(A)$ 拓扑的一组基。
\end{exercise}

\begin{exercise}
\label{exercises-exercise-radical-ideals-closed}
证明映射 $I\mapsto V(I)$ 定义一个自然双射
$$
\{I\subset A\text{ 满足 }I = \sqrt{I}\}
\longrightarrow
\{T\subset \Spec(A)\text{ 为闭集}\}
$$
\end{exercise}

\begin{definition}
\label{exercises-definition-quasi-compact}
若对任意开覆盖 $X = \bigcup_{i\in I} U_i$，都存在有限子集
$\{i_1, \ldots, i_n\}\subset I$ 使得
$X = U_{i_1}\cup\ldots U_{i_n}$，则称拓扑空间 $X$ {\it 拟紧}。
\end{definition}

\begin{exercise}
\label{exercises-exercise-spec-quasi-compact}
证明对任意环 $A$，$\Spec(A)$ 都拟紧。
\end{exercise}

\begin{definition}
\label{exercises-definition-Hausdorff}
若拓扑空间 $X$ 中任意两个不同的点 $x, y\in X$，$x\not = y$，
都有一个开子集仅含其中一点而不含另一点，则称 $X$ 满足分离公理
$T_0$。若对任意一对 $x, y\in X$，$x\not = y$，都存在不交的
开子集 $U, V$，使得 $x\in U$ 且 $y\in V$，则称 $X$
{\it 豪斯多夫}。
\end{definition}

\begin{exercise}
\label{exercises-exercise-not-hausdorff}
证明一般而言 $\Spec(A)$ {\bf 不}是豪斯多夫的。
证明 $\Spec(A)$ 满足 $T_0$。给出一个不满足 $T_0$ 的拓扑空间
$X$ 的例子。
\end{exercise}

\begin{remark}
\label{exercises-remark-not-hausdorff}
“紧”一词通常保留给既拟紧又豪斯多夫的空间。
\end{remark}

\begin{definition}
\label{exercises-definition-irreducible}
若拓扑空间 $X$ 非空，并且每当 $X = Z_1\cup Z_2$、其中
$Z_1, Z_2\subset X$ 均为闭集时，都有 $Z_1 = X$ 或
$Z_2 = X$，则称 $X$ {\it 不可约}。拓扑空间的子集
$T\subset X$ 若在 $X$ 诱导的拓扑下是不可约拓扑空间，则称
$T$ {\it 不可约}。由此定义可知，该子集不可约，当且仅当
$T$ 在 $X$ 中的闭包 $\bar T$ 不可约。
\end{definition}

\begin{exercise}
\label{exercises-exercise-irreducible-spec}
证明 $\Spec(A)$ 不可约，当且仅当 $Nil(A)$ 是素理想；并证明在
这种情形下，它是 $A$ 的唯一极小素理想。
\end{exercise}

\begin{exercise}
\label{exercises-exercise-irreducible-prime}
证明闭子集 $T\subset \Spec(A)$ 不可约，当且仅当存在素理想
${\mathfrak p}\subset A$ 使它具有形式 $T = V({\mathfrak p})$。
\end{exercise}

\begin{definition}
\label{exercises-definition-generic-point}
设 $X$ 为不可约拓扑空间。若点 $x$ 满足 $X$ 等于单点子集
$\{x\}$ 的闭包，则称其为 $X$ 的{\it 泛点}。
\end{definition}

\begin{exercise}
\label{exercises-exercise-irreducible-T0-at-most-one-generic}
证明在 $T_0$ 空间 $X$ 中，每个不可约闭子集至多有一个泛点。
\end{exercise}

\begin{exercise}
\label{exercises-exercise-spec-sober}
证明在 $\Spec(A)$ 中，每个不可约闭子集{\it 确实}有泛点。
事实上，证明映射 ${\mathfrak p} \mapsto
\overline{\{{\mathfrak p}\}}$ 是从 $\Spec(A)$ 到 $X$ 的不可约
闭子集所成集合的双射。
% P12-ERR-0423
\end{exercise}

\begin{exercise}
\label{exercises-exercise-irreducible-subset-not-generic}
给出一个例子，说明 $\Spec(\mathbf{Z})$ 的不可约子集不一定有泛点。
\end{exercise}

\begin{definition}
\label{exercises-definition-Noetherian-space}
若 $X$ 的任意闭子集降链 $Z_1\supset Z_2 \supset Z_3\supset \ldots$
都稳定，则称拓扑空间 $X$ 是{\it Noether} 的。（若任意闭子集升链
都稳定，则称它是{\it Artin} 的。）
\end{definition}

\begin{exercise}
\label{exercises-exercise-Noetherian-spec}
证明：若环 $A$ 是 Noether 环，则拓扑空间 $\Spec(A)$ 是
Noether 空间。给出一个例子，说明逆命题不成立。（如果愿意，也可对
Artin 情形做同样的事。）
\end{exercise}

\begin{definition}
\label{exercises-definition-irreducible-component}
极大不可约子集 $T\subset X$ 称为 $X$ 的一个
{\it 不可约分支}。$X$ 的任何不可约分支自动是 $X$ 的闭子集。
\end{definition}

\begin{exercise}
\label{exercises-exercise-irreducible-in-irreducible}
证明 $X$ 的任意不可约子集都包含于 $X$ 的某个不可约分支中。
\end{exercise}

\begin{exercise}
\label{exercises-exercise-Noetherian-finite-nr-irreducible}
证明 Noether 拓扑空间 $X$ 只有有限多个不可约分支，记为
$X_1, \ldots, X_n$，且 $X = X_1\cup X_2\cup\ldots\cup X_n$。
（注意，任意 $X$ 总是其不可约分支的并；不过，例如若
$X = {\mathbf R}$ 取通常拓扑，则 $X$ 的不可约分支都是单点子集，
这并没有多大意思。）
\end{exercise}

\begin{exercise}
\label{exercises-exercise-irreducible-components-minimal-primes}
证明 $\Spec(A)$ 的不可约分支对应于 $A$ 的极小素理想。
\end{exercise}

\begin{definition}
\label{exercises-definition-closed}
若 $x\in X$ 满足 $\overline{\{x\}} = \{ x\}$，则称该点
为{\it 闭点}。设 $x, y$ 是 $X$ 的点。若
$x\in \overline{\{y\}}$，则称 $x$ 是 $y$ 的一个{\it 专化}，
或称 $y$ 是 $x$ 的一个{\it 一般化}。
\end{definition}

\begin{exercise}
\label{exercises-exercise-closed-maximal}
证明 $\Spec(A)$ 的闭点对应于 $A$ 的极大理想。
\end{exercise}

\begin{exercise}
\label{exercises-exercise-generalization}
证明 ${\mathfrak p}$ 是 $\Spec(A)$ 中 ${\mathfrak q}$ 的一般化，
当且仅当 ${\mathfrak p}\subset {\mathfrak q}$。用一般化与专化
刻画闭点、极大理想、泛点和极小素理想。（这里采用如下术语：
可约拓扑空间 $X$ 的一个点若是 $X$ 某个不可约分支的泛点，就称它为
% P12-ERR-0424
一个泛点。）
\end{exercise}

\begin{exercise}
\label{exercises-exercise-disjoint-closed-spec}
设 $I$ 和 $J$ 为 $A$ 的理想。$V(I)$ 与 $V(J)$ 不交的条件是什么？
\end{exercise}

\begin{definition}
\label{exercises-definition-connected-component}
若拓扑空间 $X$ 非空，且不能写成两个非空不交开子集的并，则称
它{\it 连通}。它的极大连通子集称为 $X$ 的一个
{\it 连通分支}。$X$ 的每个点都包含于 $X$ 的某个连通分支中，
而 $X$ 的每个连通分支在 $X$ 中都是闭的。（但一般而言，连通分支
在 $X$ 中未必是开的。）
\end{definition}

\begin{exercise}
\label{exercises-exercise-disconnected-spec}
设 $A$ 为非零环。证明 $\Spec(A)$ 不连通，当且仅当对某些非零环
$B, C$ 有 $A\cong B \times C$。
\end{exercise}

\begin{exercise}
\label{exercises-exercise-connected-component-stable-generalization}
设 $T$ 为 $\Spec(A)$ 的一个连通分支。证明 $T$ 在一般化下稳定。
若 $A$ 是 Noether 环，证明 $T$ 是 $\Spec(A)$ 的开子集。
（注：例如，当 $A$ 是无穷多个 ${\mathbf F}_2$ 的直积时，该结论错误。
此环的谱由无穷多个闭点组成。）
\end{exercise}

\begin{exercise}
\label{exercises-exercise-primes-kx}
计算 $\Spec(k[x])$，即描述该环的素理想、可能的专化以及拓扑。
（在 $k$ 为代数闭域时完成计算，也在 $k$ 不是代数闭域时完成计算。）
\end{exercise}

\begin{exercise}
\label{exercises-exercise-primes-kxy}
计算 $\Spec(k[x, y])$，其中 $k$ 是代数闭域。
[提示：使用态射 $\varphi : \Spec(k[x, y]) \to \Spec(k[x])$；若
$\varphi({\mathfrak p}) = (0)$，则关于
$S = \{f\in k[x] \mid f \not = 0\}$ 局部化，并使用课堂上关于
局部化与 $\Spec$ 的结果。]
（你认为代数几何学家为什么称它为仿射 2-空间？）
\end{exercise}

\begin{exercise}
\label{exercises-exercise-primes-Zy}
计算 $\Spec(\mathbf{Z}[y])$。
[提示：同上。]（$\mathbf{Z}$ 上的仿射 1-空间。）
\end{exercise}






\section{局部化}
\label{exercises-section-localization}


\begin{exercise}
\label{exercises-exercise-submodule-localization}
设 $A$ 为环，$S \subset A$ 为乘法子集，$M$ 为 $A$-模，
$N \subset S^{-1}M$ 为 $S^{-1}A$-子模。证明存在 $A$-子模
$N' \subset M$ 使得 $N = S^{-1}N'$。（这一有用结果特别适用于
$S^{-1}A$ 的理想。）
\end{exercise}

\begin{exercise}
\label{exercises-exercise-localize-zero}
设 $A$ 为环，$M$ 为 $A$-模，$m \in M$。
\begin{enumerate}
\item 证明 $I = \{a \in A \mid am = 0\}$ 是 $A$ 的理想。
\item 对 $A$ 的素理想 $\mathfrak p$，证明 $m$ 在
$M_\mathfrak p$ 中的像为零，当且仅当 $I \not \subset \mathfrak p$。
\item 证明 $m$ 为零，当且仅当 $m$ 在所有 $A$ 的素理想
$\mathfrak p$ 所对应的 $M_\mathfrak p$ 中的像都为零。
\item 证明 $m$ 为零，当且仅当 $m$ 在所有 $A$ 的极大理想
$\mathfrak m$ 所对应的 $M_\mathfrak m$ 中的像都为零。
\item 证明 $M = 0$，当且仅当对所有极大理想 $\mathfrak m$，
$M_{\mathfrak m}$ 都为零。
\end{enumerate}
\end{exercise}

\begin{exercise}
\label{exercises-exercise-localization-is-field}
求一对 $(A, f)$，其中 $A$ 是具有三个或更多两两不同素理想的整环，
$f \in A$ 是一个元素，并且主局部化
$A_f = \{1, f, f^2, \ldots \}^{-1}A$ 是域。
\end{exercise}

\begin{exercise}
\label{exercises-exercise-localize-finite-module-zero}
设 $A$ 为环，$M$ 为有限 $A$-模，$S \subset A$ 为乘法集。
假设 $S^{-1}M = 0$。证明存在 $f \in S$，使得主局部化
$M_f = \{1, f, f^2, \ldots \}^{-1}M$ 为零。
\end{exercise}

\begin{exercise}
\label{exercises-exercise-localization-is-quotient}
给出三元组 $(A, I, S)$ 的例子，其中 $A$ 是环，
$0 \not = I \not = A$ 是真非零理想，$S \subset A$ 是乘法子集，
并且作为 $A$-代数有 $A/I \cong S^{-1}A$。
\end{exercise}




\section{Nakayama 引理}
\label{exercises-section-nakayama}

\begin{exercise}
\label{exercises-exercise-nakayama}
设 $A$ 为环，$I$ 为 $A$ 的理想，$M$ 为 $A$-模，并设
$x_1, \ldots, x_n \in M$。假设：
\begin{enumerate}
\item $M/IM$ 由 $x_1, \ldots, x_n$ 生成；
\item $M$ 是有限 $A$-模；
\item $I$ 包含于 $A$ 的每个极大理想中。
\end{enumerate}
证明 $x_1, \ldots, x_n$ 生成 $M$。（建议解法：使用练习
\ref{exercises-exercise-localize-zero} 以及局部化的正合性，化归到 $A$
的一个极大理想处的局部化。然后考察 $M$ 模去由
$x_1, \ldots, x_n$ 生成的子模所得的商，化归到课堂上的
Nakayama 引理。）
\end{exercise}







\section{长度}
\label{exercises-section-length}

\begin{definition}
\label{exercises-definition-length}
设 $A$ 为环，$M$ 为 $A$-模。
% P12-ERR-0425
$M$ 作为 $R$-模的{\it 长度}为
$$
\text{length}_A(M)
=
\sup
\{
n
\mid
\exists\ 0 = M_0 \subset M_1 \subset \ldots \subset M_n = M,
\text{ }M_i \not = M_{i + 1}
\}.
$$
换言之，它是子模链长度的上确界。
\end{definition}

\begin{exercise}
\label{exercises-exercise-length-is-one}
证明环 $A$ 上的模 $M$ 长度为 $1$，当且仅当它同构于
$A/\mathfrak m$，其中 $\mathfrak m$ 是 $A$ 的某个极大理想。
\end{exercise}

\begin{exercise}
\label{exercises-exercise-length-easy}
计算下列各模在相应环上的长度。简要（！）说明你的答案。
（可自由使用短正合列中长度函数的可加性；见《代数》引理
\ref{algebra-lemma-length-additive}。）
\begin{enumerate}
\item $\mathbf{Z}/120\mathbf{Z}$ 在 $\mathbf{Z}$ 上的长度。
\item $\mathbf{C}[x]/(x^{100} + x + 1)$ 在 $\mathbf{C}[x]$ 上的长度。
\item $\mathbf{R}[x]/(x^4 + 2x^2 + 1)$ 在 $\mathbf{R}[x]$ 上的长度。
\end{enumerate}
\end{exercise}

\begin{exercise}
\label{exercises-exercise-compute-length}
设 $A = k[x, y]_{(x, y)}$ 为仿射平面在原点处的局部环。
可对域 $k$ 作任意你喜欢的假设。设
$f = x^3 + x^2y^2 + y^{100}$ 且 $g = y^3 - x^{999}$。
$A/(f, g)$ 作为 $A$-模的长度是多少？（一种可能的方法是：
考察 $f$ 和 $g$ 在形如
% P12-ERR-0426
$A/{\mathfrak m}_A^n= k[x, y]/(x, y)^n$ 的商中生成的理想，其中
$n$ 变化。尝试寻找 $n$，使得
$A/(f, g)+{\mathfrak m}_A^n \cong A/(f, g)+{\mathfrak m}_A^{n + 1}$，
并使用 NAK。）
\end{exercise}






\section{伴随素理想}
\label{exercises-section-ass}

\noindent
伴随素理想在《代数》第 \ref{algebra-section-ass} 节中讨论。

\begin{exercise}
\label{exercises-exercise-compute-ass}
计算下列各模的伴随素理想集合。
\begin{enumerate}
\item $R = k[x, y]$ 且 $M = R/(xy(x + y))$；
\item $R = \mathbf{Z}[x]$ 且 $M = R/(300x + 75)$；
\item $R = k[x, y, z]$ 且 $M = R/(x^3, x^2y, xz)$。
\end{enumerate}
这里照例 $k$ 是域。
\end{exercise}

\begin{exercise}
\label{exercises-exercise-prime-power-not-primary}
给出 Noether 环 $R$ 及素理想 $\mathfrak p$ 的例子，使得
$\mathfrak p$ 不是 $R/\mathfrak p^2$ 的唯一伴随素理想。
\end{exercise}

\begin{exercise}
\label{exercises-exercise-product-primes-not-only-associated}
设 $R$ 为 Noether 环，具有不可比的素理想
$\mathfrak p$、$\mathfrak q$；即
$\mathfrak p \not \subset \mathfrak q$ 且
$\mathfrak q \not \subset \mathfrak p$。
\begin{enumerate}
\item 证明对 $N = R/(\mathfrak p \cap \mathfrak q)$ 有
$\text{Ass}(N) = \{\mathfrak p, \mathfrak q\}$。
\item 用例子说明，模 $M = R/\mathfrak p \mathfrak q$ 可以有一个
既不等于 $\mathfrak p$ 也不等于 $\mathfrak q$ 的伴随素理想。
\end{enumerate}
\end{exercise}





\section{Ext 群}
\label{exercises-section-ext}

\noindent
Ext 群在《代数》第 \ref{algebra-section-ext} 节中定义。

\begin{exercise}
\label{exercises-exercise-compute-ext-abelian-groups}
计算所给各模的全部 Ext 群 $\Ext^i(M, N)$；这里取
$\mathbf{Z}$-模范畴（亦称 Abel 群范畴）。
\begin{enumerate}
\item $M = \mathbf{Z}$ 且 $N = \mathbf{Z}$；
\item $M = \mathbf{Z}/4\mathbf{Z}$ 且 $N = \mathbf{Z}/8\mathbf{Z}$；
\item $M = \mathbf{Q}$ 且 $N = \mathbf{Z}/2\mathbf{Z}$；
\item $M = \mathbf{Z}/2\mathbf{Z}$ 且 $N = \mathbf{Q}/\mathbf{Z}$。
\end{enumerate}
\end{exercise}

\begin{exercise}
\label{exercises-exercise-compute-in-regular}
设 $R = k[x, y]$，其中 $k$ 是域。
\begin{enumerate}
\item 直接证明 Koszul 复形
$$
0 \to R
\xrightarrow{
\left(
\begin{matrix}
y \\
-x
\end{matrix}
\right)
}
R^{\oplus 2} \xrightarrow{(x, y)} R \xrightarrow{f \mapsto f(0, 0)} k \to 0
$$
是正合的。
\item 计算 $\Ext^i_R(k, k)$，其中把 $k = R/(x, y)$ 视为 $R$-模。
\end{enumerate}
\end{exercise}

\begin{exercise}
\label{exercises-exercise-infinitely-many-nonzero-ext}
给出 Noether 环 $R$ 以及有限模 $M$、$N$ 的例子，使得对所有
$i \geq 0$，$\Ext^i_R(M, N)$ 都非零。
\end{exercise}

\begin{exercise}
\label{exercises-exercise-infinite-ext}
给出环 $R$ 和理想 $I$ 的例子，使得
$\Ext^1_R(R/I, R/I)$ 不是有限 $R$-模。
（由《代数》引理 \ref{algebra-lemma-ext-noetherian} 可知，当
$R$ 为 Noether 环时不可能发生这种情形。）
\end{exercise}







\section{深度}
\label{exercises-section-depth}

\noindent
深度在《代数》第 \ref{algebra-section-depth} 节中定义，并在
《对偶化复形》第 \ref{dualizing-section-depth} 节中进一步研究。

\begin{exercise}
\label{exercises-exercise-compute-depth}
设 $R$ 为环，$I \subset R$ 为理想，$M$ 为 $R$-模。
在下列各情形中计算 $\text{depth}_I(M)$。
\begin{enumerate}
\item $R = \mathbf{Z}$，$I = (30)$，$M = \mathbf{Z}$；
\item $R = \mathbf{Z}$，$I = (30)$，$M = \mathbf{Z}/(300)$；
\item $R = \mathbf{Z}$，$I = (30)$，$M = \mathbf{Z}/(7)$；
\item $R = k[x, y, z]/(x^2 + y^2 + z^2)$，$I = (x, y, z)$，$M = R$；
\item $R = k[x, y, z, w]/(xz, xw, yz, yw)$，$I = (x, y, z, w)$，
$M = R$。
\end{enumerate}
这里 $k$ 是域。在最后两种情形中，可自由地在极大理想 $I$ 处局部化。
\end{exercise}

\begin{exercise}
\label{exercises-exercise-depth-not-inherited-localization}
给出深度 $\geq 1$ 的 Noether 局部环 $(R, \mathfrak m, \kappa)$
以及素理想 $\mathfrak p$ 的例子，使得
\begin{enumerate}
\item $\text{depth}_\mathfrak m(R) \geq 1$；
\item $\text{depth}_\mathfrak p(R_\mathfrak p) = 0$；
\item $\dim(R_\mathfrak p) \geq 1$。
\end{enumerate}
若不要求（3），这个练习就太容易了。为什么？
\end{exercise}

\begin{exercise}
\label{exercises-exercise-depth-torsion-free}
设 $(R, \mathfrak m)$ 为 Noether 局部整环，$M$ 为有限 $R$-模。
\begin{enumerate}
\item 若 $M$ 无挠，证明 $M$ 在 $R$ 上的深度至少为 $1$。
\item 给出深度恰为 $1$ 的例子。
\end{enumerate}
\end{exercise}

\begin{exercise}
\label{exercises-exercise-depth-examples}
对每一对 $m \geq n \geq 0$，给出 Noether 局部环 $R$ 的例子，使得
$\dim(R) = m$ 且 $\text{depth}(R) = n$。
\end{exercise}

\begin{exercise}
\label{exercises-exercise-make-depth-1}
设 $(R, \mathfrak m)$ 为 Noether 局部环，$M$ 为有限 $R$-模。
证明存在典范短正合列
$$
0 \to K \to M \to Q \to 0
$$
满足下列条件：
\begin{enumerate}
\item $\text{depth}(Q) \geq 1$；
\item $K$ 为零，或 $\text{Supp}(K) = \{\mathfrak m\}$；
\item $\text{length}_R(K) < \infty$。
\end{enumerate}
提示：利用 Noether 性证明存在满足（2）的极大子模 $K$，再证明
$Q = M/K$ 满足（1），且 $K$ 满足（3）。
\end{exercise}

\begin{exercise}
\label{exercises-exercise-make-depth-2}
设 $(R, \mathfrak m)$ 为 Noether 局部环，$M$ 为深度
$\geq 2$ 的有限 $R$-模。设 $N \subset M$ 为非零子模。
\begin{enumerate}
\item 证明 $\text{depth}(N) \geq 1$。
\item 证明 $\text{depth}(N) = 1$，当且仅当商模 $M/N$ 满足
$\text{depth}(M/N) = 0$。
\item 证明存在子模 $N' \subset M$，其中 $N \subset N'$ 且余长度有限，
即 $\text{length}_R(N'/N) < \infty$，使得 $N'$ 的深度
$\geq 2$。提示：把练习 \ref{exercises-exercise-make-depth-1} 应用于 $M/N$，
并取 $N'$ 为 $K$ 的逆像。
\end{enumerate}
\end{exercise}

\begin{exercise}
\label{exercises-exercise-Hartshorne-reduced}
设 $(R, \mathfrak m)$ 为 Noether 局部环。假设 $R$ 既约，即
$R$ 没有非零幂零元。再假设 $R$ 有两个不同的极小素理想
$\mathfrak p$ 和 $\mathfrak q$。
\begin{enumerate}
\item 证明 $R$-模序列
% P12-ERR-0427
$$
0 \to R \to R/\mathfrak p \oplus R/\mathfrak q \to
R/\mathfrak p + \mathfrak q \to 0
$$
是正合的（逐处检查）。这些映射分别为
$x \mapsto (x \bmod \mathfrak p, x \bmod \mathfrak q)$
以及 $(y \bmod \mathfrak p, z \bmod \mathfrak q) \mapsto
(y - z \bmod \mathfrak p + \mathfrak q)$。
\item 证明若 $\text{depth}(R) \geq 2$，则
$\dim(R/\mathfrak p + \mathfrak q) \geq 1$。
\item 证明若 $\text{depth}(R) \geq 2$，则
$U = \Spec(R) \setminus \{\mathfrak m\}$ 是连通拓扑空间。
\end{enumerate}
这证明了 Hartshorne 连通性定理的一个非常特殊的情形。该定理断言：
深度 $\text{depth} \geq 2$ 的 Noether 局部环之去心谱 $U$ 是连通的。
\end{exercise}

\begin{exercise}
\label{exercises-exercise-depth-2}
设 $(R, \mathfrak m)$ 为 Noether 局部环。设
$x, y \in \mathfrak m$ 为长度 $2$ 的正则序列。证明对任意
$n \geq 2$，不存在 $a, b \in R$ 使得
$$
x^{n - 1}y^{n - 1} = a x^n + b y^n
$$
建议：先尝试 $n = 2$，以了解论证方法。
注：此结果有一个极为广泛的推广，称为单项式猜想。
\end{exercise}







\section{Cohen--Macaulay 模与环}
\label{exercises-section-CM}

\noindent
Cohen--Macaulay 模在《代数》第 \ref{algebra-section-CM} 节中研究，
Cohen--Macaulay 环则在《代数》第
\ref{algebra-section-CM-ring} 节中研究。

\begin{exercise}
\label{exercises-exercise-examples-CM}
请在下列各情形中回答“是”或“否”，无需解释或证明。
\begin{enumerate}
\item 设 $p$ 为素数。局部环 $\mathbf{Z}_{(p)}$ 是
Cohen--Macaulay 局部环吗？
\item 设 $p$ 为素数。局部环 $\mathbf{Z}_{(p)}$ 是正则局部环吗？
\item 设 $k$ 为域。局部环 $k[x]_{(x)}$ 是
Cohen--Macaulay 局部环吗？
\item 设 $k$ 为域。局部环 $k[x]_{(x)}$ 是正则局部环吗？
\item 设 $k$ 为域。局部环 $(k[x, y]/(y^2 - x^3))_{(x, y)} =
k[x, y]_{(x, y)}/(y^2 - x^3)$ 是 Cohen--Macaulay 局部环吗？
\item 设 $k$ 为域。局部环 $(k[x, y]/(y^2, xy))_{(x, y)} =
k[x, y]_{(x, y)}/(y^2, xy)$ 是 Cohen--Macaulay 局部环吗？
\end{enumerate}
\end{exercise}













\section{奇点}
\label{exercises-section-singularities}

\begin{exercise}
\label{exercises-exercise-singularities}
任取域 $k$。设 $A = k[[x, y]]/(f)$ 且
$B = k[[u, v]]/(g)$，其中 $f = xy$，$g = uv + \delta$，并且
$\delta \in (u, v)^3$。证明 $A$ 与 $B$ 是同构的环。
\end{exercise}

\begin{remark}
\label{exercises-remark-singularities}
设曲线定义在域 $k$ 上。若曲线在某点的完备局部环形如
$k'[[x, y]]/(f)$，其中（a）$k'$ 是 $k$ 的有限可分扩张；
（b）$f$ 的初始项次数为二，即形如
$q = ax^2 + bxy + cy^2$，其中 $a, b, c\in k'$ 不全为零；
（c）$q$ 是 $k'$ 上的非退化二次型（在特征 2 时，这意味着
$b$ 非零），则称该点为普通二重点。一般而言，每个同构类的序对
$(k', q)$ 都对应这类环的一个同构类。
\end{remark}

\begin{exercise}
\label{exercises-exercise-periodic-resolution}
设 $R$ 为环，$n \geq 1$。设 $A$、$B$ 为系数在 $R$ 中的
$n \times n$ 矩阵，并且对 $R$ 中某个非零因子 $f$ 有
$AB = f 1_{n \times n}$。令 $S = R/(f)$。证明
$$
\ldots \to
S^{\oplus n} \xrightarrow{B}
S^{\oplus n} \xrightarrow{A}
S^{\oplus n} \xrightarrow{B}
S^{\oplus n} \to \ldots
$$
是正合的。
\end{exercise}







\section{可构造集}
\label{exercises-section-constructible}

\noindent
设 $k$ 为代数闭域，例如复数域 $\mathbf{C}$。设 $n \geq 0$。
多项式 $f \in k[x_1, \ldots, x_n]$ 通过取值给出函数
$f : k^n \to k$。若子集 $Z \subset k^n$ 是一族多项式的公共零点集，
则称该子集为{\it 代数集}。

\begin{exercise}
\label{exercises-exercise-finite-nr-equations}
证明代数集总能写成有限多个多项式的零点集。
\end{exercise}

\noindent
沿用上述记号。若子集 $E \subset k^n$ 是有限多个形如
$Z \cap \{f \not = 0\}$ 的集合之并，其中 $f$ 是多项式，
则称该子集{\it 可构造}。

\begin{exercise}
\label{exercises-exercise-constructible-classical}
证明下列各点：
\begin{enumerate}
\item 可构造集的补集是可构造集；
\item 有限多个可构造集的并是可构造集；
\item 有限多个可构造集的交是可构造集；
\item 任意可构造集 $E$ 都可写成有限不交并
$E = \coprod E_i$，其中每个 $E_i$ 都形如
$Z \cap \{f \not = 0\}$，$Z$ 为代数集，$f$ 为多项式。
\end{enumerate}
\end{exercise}

\begin{exercise}
\label{exercises-exercise-division-with-remainder}
设 $R$ 为环，并设
$f = a_d x^d + a_{d - 1} x^{d - 1} + \ldots + a_0 \in R[x]$。
（照例，这一记号意味着 $a_0, \ldots, a_d \in R$。）设 $g \in R[x]$。
证明可以找到 $N \geq 0$ 以及 $r, q \in R[x]$，使得
$$
a_d^N g = q f + r
$$
且 $\deg(r) < d$；即对某些 $c_i \in R$ 有
$r = c_0 + c_1 x + \ldots + c_{d - 1}x^{d - 1}$。
\end{exercise}







\section{Hilbert 零点定理}
\label{exercises-section-Hilbert-Nullstellensatz}


\begin{exercise}
\label{exercises-exercise-uncountable}
{\it 一个仅用复数的有趣论证！}
设 ${\mathbf C}$ 为复数域，$V$ 为 ${\mathbf C}$ 上的向量空间。
线性算子 $T : V \to V$ 的谱，是满足算子
$T - \lambda \text{id}_V$ 不可逆的复数 $\lambda \in {\mathbf C}$
所成的集合。
\begin{enumerate}
\item 证明 $\mathbf{C}(X)$ 在 ${\mathbf C}$ 上具有不可数维数。
\item 证明：若 $V$ 的维数有限或可数，则 $V$ 上的任意线性算子
都有非空谱。
\item 证明：若有限生成 ${\mathbf C}$-代数 $R$ 是域，则映射
${\mathbf C}\to R$ 是同构。
\item 证明 ${\mathbf C}[x_1, \ldots, x_n]$ 的任意极大理想
${\mathfrak m}$ 都形如
$(x_1-\alpha_1, \ldots, x_n-\alpha_n)$，其中某些
$\alpha_i \in {\mathbf C}$。
\end{enumerate}
\end{exercise}

\begin{remark}
\label{exercises-remark-HNSS}
设 $k$ 为域。那么，对每个整数 $n\in {\mathbf N}$ 以及每个极大理想
${\mathfrak m} \subset k[x_1, \ldots, x_n]$，商
$k[x_1, \ldots, x_n]/{\mathfrak m}$ 都是 $k$ 的有限域扩张。
这将在课程后面证明。当然（请检查），它蕴含有限生成
$k$-代数之极大理想的类似结论。上面的练习在
$k = {\mathbf C}$ 的情形证明了这一点。
\end{remark}

\begin{exercise}
\label{exercises-exercise-Hilbert-Nullstellensatz}
设 $k$ 为域。请使用注 \ref{exercises-remark-HNSS}。
\begin{enumerate}
\item 设 $R$ 为 $k$-代数。假设 $\dim_k R < \infty$ 且
$R$ 是整环。证明 $R$ 是域。
\item 假设 $R$ 是有限生成 $k$-代数，且 $f\in R$ 不是幂零元。
证明存在极大理想 ${\mathfrak m} \subset R$，使得
$f\not\in {\mathfrak m}$。
\item 用例子说明，当 $R$ 不是域上的有限型代数时，该命题不成立。
\item 证明任意根理想 $I \subset {\mathbf C}[x_1, \ldots, x_n]$
都是所有包含它的极大理想之交。
\end{enumerate}
\end{exercise}

\begin{remark}
\label{exercises-remark-Hilbert-Nullstellensatz}
这就是 Hilbert 零点定理。它断言
$\Spec(k[x_1, \ldots, x_n])$ 的闭子集
（由前面的练习，它们对应于根理想）由其中包含的闭点决定。
\end{remark}

\begin{exercise}
\label{exercises-exercise-product-matrices-ring}
设
$A =
{\mathbf C}[x_{11}, x_{12}, x_{21}, x_{22}, y_{11}, y_{12}, y_{21}, y_{22}]$。
设 $I$ 为 $A$ 中由矩阵 $XY$ 的各个元素生成的理想，其中
$$
X = \left(
\begin{matrix}
x_{11} & x_{12}\\
x_{21} & x_{22}
\end{matrix}
\right)
\quad\text{且}\quad
Y = \left(
\begin{matrix}
y_{11} & y_{12}\\
y_{21} & y_{22}
\end{matrix}
\right).
$$
求 $\Spec(A)$ 的闭子集 $V(I)$ 的不可约分支。
（这里是要描述这些分支并给出每个分支的方程；不必证明所写方程定义
素理想。）提示：
\begin{enumerate}
\item 可以使用 Hilbert 零点定理；只需找出覆盖 $V(I)$
之闭点集的不可约局部闭子集。
\item 有两个容易找到的分支。
\item 不可约集在连续映射下的像不可约。
\end{enumerate}
\end{exercise}




\section{维数}
\label{exercises-section-dimension}

\begin{exercise}
\label{exercises-exercise-dimension-bigger-one-finite-nr-primes}
构造仅有有限多个素理想、但维数 $> 1$ 的环 $A$。
\end{exercise}

\begin{exercise}
\label{exercises-exercise-hypersurface-in-A2-dimension-one}
设 $f \in \mathbf{C}[x, y]$ 为非常数多项式。
证明 $\mathbf{C}[x, y]/(f)$ 的维数为 $1$。
\end{exercise}

\begin{exercise}
\label{exercises-exercise-dimension-polynomial-ring}
设 $(R, \mathfrak m)$ 为 Noether 局部环。设 $n \geq 1$。
在多项式环 $R[x_1, \ldots, x_n]$ 中令
$\mathfrak m' = (\mathfrak m, x_1, \ldots, x_n)$。证明
$$
\dim(R[x_1, \ldots, x_n]_{\mathfrak m'}) = \dim(R) + n.
$$
\end{exercise}





\section{链环}
\label{exercises-section-catenary}

\begin{definition}
\label{exercises-definition-catenary}
若对任意三个素理想
${\mathfrak p}_1 \subset {\mathfrak p}_2 \subset {\mathfrak p}_3$ 均有
$$
ht({\mathfrak p}_3 / {\mathfrak p}_1) = ht({\mathfrak p}_3/{\mathfrak p}_2) +
ht({\mathfrak p}_2/{\mathfrak p}_1).
$$
则称 Noether 环 $A$ 为{\it 链环}。这里
$ht(\mathfrak p/\mathfrak q)$ 表示 $\mathfrak p/\mathfrak q$
在环 $A/\mathfrak q$ 中的高度。写成公式即
$$
ht(\mathfrak p/\mathfrak q) =
\dim(A_\mathfrak p/\mathfrak qA_\mathfrak p) =
\dim((A/\mathfrak q)_\mathfrak p) =
\dim((A/\mathfrak q)_{\mathfrak p/\mathfrak q})
$$
若 $T \subset T' \subset X$，其中 $T$ 与 $T'$ 为闭不可约子集，
则存在一条不可约闭子集的极大链
$$
T = T_0 \subset T_1 \subset \ldots \subset T_n = T'
$$
并且每条这样的链都有相同的有限长度，则称拓扑空间 $X$ {\it 链状}。
\end{definition}

\begin{exercise}
\label{exercises-exercise-catenary-the-same}
证明：对 Noether 环而言，《代数》定义
\ref{algebra-definition-catenary} 中定义的链状性概念，与定义
\ref{exercises-definition-catenary} 中的概念一致。
\end{exercise}

\begin{exercise}
\label{exercises-exercise-Noetherian-local-domain-dim-2-catenary}
证明维数为 $2$ 的 Noether 局部整环是链环。
\end{exercise}

\begin{exercise}
\label{exercises-exercise-finite-type-over-field-catenary}
设 $k$ 为域。证明有限型 $k$-代数是链环。
\end{exercise}

\begin{exercise}
\label{exercises-exercise-example-no-dim-function}
给出一个有限、清醒且链状的拓扑空间 $X$，使它不存在维数函数
$\delta : X \to \mathbf{Z}$。这里，映射
$\delta : X \to \mathbf{Z}$ 称为维数函数，是指对
$x, y \in X$ 有：
\begin{enumerate}
\item 若 $x \leadsto y$ 且 $x \not = y$，则
$\delta(x) > \delta(y)$；
\item 若 $x \leadsto y$ 且 $\delta(x) \geq \delta(y) + 2$，
则存在 $z \in X$，使得 $x \leadsto z \leadsto y$ 且
$\delta(x) > \delta(z) > \delta(y)$。
\end{enumerate}
清楚地描述你的空间，并简洁说明为什么它不可能有维数函数。
\end{exercise}




\section{分式域}
\label{exercises-section-fraction-fields}

\begin{exercise}
\label{exercises-exercise-find-fraction-field}
考察整环
$$
{\mathbf Q}[r, s, t]/(s^2-(r-1)(r-2)(r-3), t^2-(r + 1)(r + 2)(r + 3)).
$$
找出一个形如 ${\mathbf Q}[x, y]/(f)$ 的整环，使其分式域与之同构。
\end{exercise}



\section{超越次数}
\label{exercises-section-transcendence}

\begin{exercise}
\label{exercises-exercise-algebraic-extension}
设 $K'/K/k$ 为域扩张，并且 $K'$ 在 $K$ 上代数。
证明 $\text{trdeg}_k(K) = \text{trdeg}_k(K')$。
（提示：证明若 $x_1, \ldots, x_d \in K$ 在 $k$ 上代数无关，且
$d < \text{trdeg}_k(K')$，则 $k(x_1, \ldots, x_d) \subset K$
不可能是代数扩张。）
\end{exercise}

\begin{exercise}
\label{exercises-exercise-growth-powers-subvector-space}
设 $k$ 为域，$K/k$ 为超越次数 $d$ 的有限生成扩张。
若 $V, W \subset K$ 是有限维 $k$-向量子空间，记
$$
VW = \{f \in K \mid f = \sum\nolimits_{i = 1, \ldots, n} v_i w_i
\text{，其中某个 }n\text{ 且 }v_i \in V, w_i \in W\}
$$
这是有限维 $k$-向量子空间。令 $V^2 = VV$、$V^3 = V V^2$，依此类推。
\begin{enumerate}
\item 证明可找到 $V \subset K$ 和 $\epsilon > 0$，使得对所有
$n \geq 1$ 都有 $\dim V^n \geq \epsilon n^d$。
\item 反过来，证明对每个有限维 $V \subset K$，存在 $C > 0$，
使得对所有 $n \geq 1$ 都有 $\dim V^n \leq C n^d$。
（一种可能的做法是：先对 $k[x_1, \ldots, x_d]$ 的向量子空间证明；
再对 $k(x_1, \ldots, x_d)$ 的向量子空间证明；最后，若
$K/k(x_1, \ldots, x_d)$ 是有限扩张，就选取 $K$ 在
$k(x_1, \ldots, x_d)$ 上的一组基，并按此基展开论证。）
\item 推出：可以用 $K$ 的有限维向量子空间各次幂的增长来重新定义
超越次数。
\end{enumerate}
这与 $k$ 上（非交换）代数的 Gelfand--Kirillov 维数有关。
\end{exercise}








\section{纤维维数}
\label{exercises-section-dimension-fibres}

\noindent
下面是一些与维数公式有关的问题；见《代数》第
\ref{algebra-section-dimension-formula} 节。

\begin{exercise}
\label{exercises-exercise-nr-components-fibre}
设 $k$ 为你偏爱的代数闭域。下文中 $k[x]$ 和 $k[x, y]$
表示多项式环。
\begin{enumerate}
\item 对每个整数 $n \geq 0$，找出整环的有限型扩张
$k[x] \subset A$，使得 $A/xA$ 的谱恰有 $n$ 个不可约分支。
\item 构造整环的有限型扩张 $k[x] \subset A$，使得对每个
$\alpha \in k$，$A/(x - \alpha)A$ 的谱都非空且可约。
\item 构造整环的有限型扩张 $k[x, y] \subset A$，使得对所有
$(\alpha, \beta) \in k^2 \setminus \{(0, 0)\}$，
$A/(x - \alpha, y - \beta)A$ 的谱都不可约
\footnote{回忆：不可约蕴含非空。}，而 $A/(x, y)A$ 的谱非空且可约。
\end{enumerate}
\end{exercise}

\begin{exercise}
\label{exercises-exercise-codim-1}
设 $k$ 为你偏爱的代数闭域。设 $n \geq 1$。
令 $k[x_1, \ldots, x_n]$ 为多项式环，并令
$\mathfrak m = (x_1, \ldots, x_n)$。设
$k[x_1, \ldots, x_n] \subset A$ 为整环的有限型扩张。
令 $d = \dim(A)$。
\begin{enumerate}
\item 若 $A/\mathfrak mA \not = 0$，证明
$d - 1 \geq \dim(A/\mathfrak m A) \geq d - n$。
\item 用例子说明每个数值都能出现。
\item 用例子说明 $\Spec(A/\mathfrak m A)$ 可以有维数不同的
不可约分支。
\end{enumerate}
\end{exercise}




\section{有限局部自由模}
\label{exercises-section-finite-locally-free}

\begin{definition}
\label{exercises-definition-finite-locally-free}
设 $A$ 为环。回忆：{\it 有限局部自由} $A$-模 $M$，是指满足如下条件的模：
对每个 ${\mathfrak p} \in \Spec(A)$，都存在
$f\in A$、$f \not \in {\mathfrak p}$，使 $M_f$ 是有限自由
$A_f$-模。若 $M$ 是秩为 $1$ 的有限局部自由模，即对每个
${\mathfrak p} \in \Spec(A)$，都存在
$f\in A$、$f \not \in \mathfrak p$，使得作为 $A_f$-模有
$M_f \cong A_f$，则称 $M$ 为{\it 可逆模}。
\end{definition}

\begin{exercise}
\label{exercises-exercise-tensor-finite-locally-free}
证明有限局部自由模的张量积仍有限局部自由。证明两个可逆模的张量积
仍可逆。
\end{exercise}

\begin{definition}
\label{exercises-definition-class-group}
设 $A$ 为环。$A$ 的{\it 类群}，有时也称 $A$ 的{\it Picard 群}，
是可逆 $A$-模的同构类集合 $\Pic(A)$，其群运算由张量积定义
（见练习 \ref{exercises-exercise-tensor-finite-locally-free}）。
\end{definition}

\noindent
注意，$A$ 的类群平凡，当且仅当每个可逆模都同构于秩为 1 的自由模。

\begin{exercise}
\label{exercises-exercise-class-group-trivial}
证明下列各环的类群平凡：
\begin{enumerate}
\item 多项式环 $A = k[x]$，其中 $k$ 为域；
\item 整数环 $A = \mathbf{Z}$；
\item 多项式环 $A = k[x, y]$，其中 $k$ 为域；
\item 商环 $k[x, y]/(xy)$，其中 $k$ 为域。
\end{enumerate}
\end{exercise}

\begin{exercise}
\label{exercises-exercise-class-group-not-trivial}
设 $k$ 为特征不为 $2$ 的域，并设
$f(x) = (x - t_1) \ldots (x - t_n)$，其中
$t_1, \ldots, t_n \in k$ 两两不同，$n \geq 3$ 为奇数。
证明环 $A = k[x, y]/(y^2 - f(x))$ 的类群非平凡。
（提示：证明理想 $(y, x - t_1)$ 定义 $\Pic(A)$ 的一个非平凡元素。）
\end{exercise}

\begin{exercise}
\label{exercises-exercise-trace-det}
设 $A$ 为环。
\begin{enumerate}
\item 假设 $M$ 是有限局部自由 $A$-模，且
$\varphi : M \to M$ 是自同态。定义或构造 $\varphi$ 的{\it 迹}
与{\it 行列式}，并证明你的构造“关于三元组 $(A, M, \varphi)$
具有函子性”。
\item 证明：若 $M, N$ 是有限局部自由 $A$-模，并且
$\varphi : M \to N$、$\psi : N \to M$，则
% P12-ERR-0428
$\text{Trace}(\varphi \circ \psi) = \text{Trace}(\psi \circ \varphi)$ 且
$\det(\varphi \circ \psi) = \det(\psi \circ \varphi)$。
\item 当 $M$ 有限局部自由时，证明 $\text{Trace}$ 定义
$A$-线性映射 $\text{End}_A(M) \to A$，而 $\det$ 定义乘法映射
$\text{End}_A(M) \to A$。
\end{enumerate}
\end{exercise}

\begin{exercise}
\label{exercises-exercise-trace-det-rings}
现在假设 $B$ 是一个 $A$-代数，且作为 $A$-模有限局部自由；
换言之，$B$ 是有限局部自由 $A$-代数。
\begin{enumerate}
\item 使用练习 \ref{exercises-exercise-trace-det} 中的 $\text{Trace}$ 和
$\det$ 定义 $\text{Trace}_{B/A}$ 与 $\text{Norm}_{B/A}$。
\item 设 $b\in B$，并设 $\pi : \Spec(B) \to \Spec(A)$ 为诱导态射。
证明 $\pi(V(b)) = V(\text{Norm}_{B/A}(b))$。
（回忆 $V(f) = \{ {\mathfrak p} \mid f \in {\mathfrak p}\}$。）
\item （基变换。）假设 $i : A \to A'$ 是环映射。令
$B' = B \otimes_A A'$。说明为什么 $i(\text{Norm}_{B/A}(b))$ 等于
$\text{Norm}_{B'/A'}(b \otimes 1)$。
\item 当 $B = A \times A \times A \times \ldots \times A$ 且
$b = (a_1, \ldots, a_n)$ 时，计算 $\text{Norm}_{B/A}(b)$。
\item 对有限平坦映射
${\mathbf Q}[x] \to {\mathbf Q}[y]$、$x \to y^n$，计算
$y-y^3$ 的范数。（提示：使用“基变换”
$A = {\mathbf Q}[x] \subset A' = {\mathbf Q}(\zeta_n)(x^{1/n})$。）
\end{enumerate}
\end{exercise}



\section{粘合}
\label{exercises-section-glueing}

\begin{exercise}
\label{exercises-exercise-cover}
假设 $A$ 是环，$M$ 是 $A$-模。设 $f_i$、$i \in I$
为 $A$ 的一族元素，并且
$$
\Spec(A) = \bigcup D(f_i).
$$
\begin{enumerate}
\item 证明若 $M_{f_i}$ 是有限 $A_{f_i}$-模，则 $M$ 是有限
$A$-模。
\item 证明若 $M_{f_i}$ 是平坦 $A_{f_i}$-模，则 $M$ 是平坦
$A$-模。（仔细想想，这一点其实有些显然。）
\end{enumerate}
\end{exercise}

\begin{remark}
\label{exercises-remark-cover}
用代数几何语言来说，这意味着“有限生成”或“平坦”这一性质
（在适当意义下）对 Zariski 拓扑是局部的。还可以对“有限表示”这一
性质证明相同结论。
\end{remark}

\begin{exercise}
\label{exercises-exercise-cover-ring-map}
假设 $A \to B$ 是环映射。设 $f_i \in A$、$i \in I$ 以及
$g_j \in B$、$j \in J$ 为两族元素，并且
$$
\Spec(A) = \bigcup D(f_i)
\quad\text{且}\quad
\Spec(B) = \bigcup D(g_j).
$$
证明：若对所有 $i, j$，$A_{f_i} \to B_{f_ig_j}$ 都是有限型的，
则 $A \to B$ 是有限型的。
\end{exercise}




\section{上升与下降}
\label{exercises-section-going-up}


\begin{definition}
\label{exercises-definition-GU-GD}
设 $\phi : A \to B$ 为环同态。若满足下列条件，则称 $\phi$
满足{\it 上升定理}：
\begin{itemize}
\item[(GU)] 对任意
${\mathfrak p}, {\mathfrak p}' \in \Spec(A)$，若
${\mathfrak p} \subset {\mathfrak p}'$，并且任取位于
${\mathfrak p}$ 上方的 $P \in \Spec(B)$，都存在位于
${\mathfrak p}'$ 上方的 $P'\in \Spec(B)$，使得 $P \subset P'$。
\end{itemize}
类似地，若满足下列条件，则称 $\phi$ 满足{\it 下降定理}：
\begin{itemize}
\item[(GD)] 对任意
${\mathfrak p}, {\mathfrak p}' \in \Spec(A)$，若
${\mathfrak p} \subset {\mathfrak p}'$，并且任取位于
${\mathfrak p}'$ 上方的 $P' \in \Spec(B)$，都存在位于
${\mathfrak p}$ 上方的 $P\in \Spec(B)$，使得 $P \subset P'$。
\end{itemize}
\end{definition}

\begin{exercise}
\label{exercises-exercise-GU-GD}
在下列每一种情形中，判断（GU）与（GD）是否成立，并说明理由。
（可以使用能找到的任何命题、定理或引理，但要逐一检查假设。）
\begin{enumerate}
\item $k$ 是域，$A = k$，$B = k[x]$。
\item $k$ 是域，$A = k[x]$，$B = k[x, y]$。
\item $A = {\mathbf Z}$，$B = {\mathbf Z}[1/11]$。
\item $k$ 是代数闭域，$A = k[x, y]$，
$B = k[x, y, z]/(x^2-y, z^2-x)$。
\item $A = {\mathbf Z}$，$B = {\mathbf Z}[i, 1/(2 + i)]$。
\item $A = {\mathbf Z}$，$B = {\mathbf Z}[i, 1/(14 + 7i)]$。
\item $k$ 是代数闭域，$A = k[x]$，
$B = k[x, y, 1/(xy-1)]/(y^2-y)$。
\end{enumerate}
\end{exercise}

\begin{exercise}
\label{exercises-exercise-image}
设 $A$ 为环，$B = A[x]$ 为 $A$ 上的单变量多项式代数。设
$f = a_0 + a_1 x + \ldots + a_r x^r \in B = A[x]$。
仔细证明 $D(f)$ 在 $\Spec(A)$ 中的像等于
$D(a_0) \cup \ldots \cup D(a_r)$。
\end{exercise}

\begin{exercise}
\label{exercises-exercise-images}
设 $k$ 为代数闭域。计算下列映射在 $\Spec(k[x, y])$ 中的像：
\begin{enumerate}
\item $\Spec(k[x, yx^{-1}]) \to \Spec(k[x, y])$，其中
$k[x, y] \subset k[x, yx^{-1}] \subset k[x, y, x^{-1}]$。
（提示：为避免混淆，给元素 $yx^{-1}$ 另取一个名字。）
\item $\Spec(k[x, y, a, b]/(ax-by-1))\to \Spec(k[x, y])$。
\item $\Spec(k[t, 1/(t-1)]) \to \Spec(k[x, y])$，它由
$x \mapsto t^2$ 和 $y \mapsto t^3$ 诱导。
\item $k = {\mathbf C}$（复数），
$\Spec(k[s, t]/(s^3 + t^3-1)) \to \Spec(k[x, y])$，其中
$x\mapsto s^2$，$y \mapsto t^2$。
\end{enumerate}
\end{exercise}

\begin{remark}
\label{exercises-remark-elimination-theory}
求上述像通常要使用消元理论。
\end{remark}




\section{Fitting 理想}
\label{exercises-section-fitting-ideals}

\begin{exercise}
\label{exercises-exercise-fitting}
设 $R$ 为环，$M$ 为有限 $R$-模。选取 $M$ 的一个表示
$$
\bigoplus\nolimits_{j \in J} R \longrightarrow R^{\oplus n}
\longrightarrow M \longrightarrow 0.
$$
注意，映射 $R^{\oplus n} \to M$ 由 $M$ 的一列元素
$x_1, \ldots, x_n$ 给出；这些元素 $x_i$ 是 $M$ 的{\it 生成元}。
映射 $\bigoplus_{j \in J} R \to R^{\oplus n}$ 由系数在 $R$ 中的
$n \times J$ 矩阵 $A$ 给出。换言之，
$A = (a_{ij})_{i = 1, \ldots, n, j \in J}$。$A$ 的各列
$(a_{1j}, \ldots, a_{nj})$（$j \in J$）称为{\it 关系}。
任何满足 $\sum r_i x_i = 0$ 的向量 $(r_i) \in R^{\oplus n}$
都是 $A$ 的各列的线性组合。当然，每个有限 $R$-模都有许多不同的表示。
\begin{enumerate}
\item 证明由 $A$ 的 $(n - k) \times (n - k)$ 子式生成的理想
与表示的选择无关。该理想称为 $M$ 的{\it 第 $k$ 个 Fitting 理想}，
记作 $Fit_k(M)$。
\item 证明
$Fit_0(M) \subset Fit_1(M) \subset Fit_2(M) \subset \ldots$。
（提示：利用按一列展开来计算行列式。）
\item 证明下列条件等价：
\begin{enumerate}
\item $Fit_{r - 1}(M) = (0)$ 且 $Fit_r(M) = R$；
\item $M$ 是秩为 $r$ 的局部自由模。
\end{enumerate}
\end{enumerate}
\end{exercise}




\section{Hilbert 函数}
\label{exercises-section-hilbert}

\begin{definition}
\label{exercises-definition-numerical-polynomial}
{\it 数值多项式}是满足如下条件的多项式
$f(x) \in {\mathbf Q}[x]$：对每个整数 $n$，都有
$f(n) \in {\mathbf Z}$。
\end{definition}

\begin{definition}
\label{exercises-definition-graded-module}
环 $A$ 上的{\it 分次模} $M$，是带有 $A$-子模直和分解
$
\bigoplus\nolimits_{n \in {\mathbf Z}} M_n
$
的 $A$-模 $M$。若所有 $M_n$ 都是有限 $A$-模，则称 $M$
{\it 局部有限}。假设 $A$ 是 Noether 环，并且 $\varphi$
是有限 $A$-模上的一个{\it Euler--Poincar\'e 函数}。
这意味着：对每个有限生成 $A$-模 $M$，给定一个整数
$\varphi(M) \in {\mathbf Z}$；并且对每个短正合列
$$
0
\longrightarrow
M'
\longrightarrow
M
\longrightarrow
M''
\longrightarrow
0
$$
都有 $\varphi(M) = \varphi(M') + \varphi(M'')$。
局部有限分次模 $M$（关于 $\varphi$）的{\it Hilbert 函数}是函数
$\chi_\varphi(M, n) = \varphi(M_n)$。若存在某个数值多项式
$P_\varphi$，使得对所有充分大的整数 $n$ 都有
$\chi_\varphi(M, n) = P_\varphi(n)$，则称 $M$ 有
{\it Hilbert 多项式}。
\end{definition}

\begin{definition}
\label{exercises-definition-graded-algebra}
{\it 分次 $A$-代数}是分次 $A$-模
$B = \bigoplus_{n \geq 0} B_n$，连同一个 $A$-双线性映射
$$
B \times B \longrightarrow B, \ (b, b') \longmapsto bb'
$$
使 $B$ 成为 $A$-代数，且 $B_n \cdot B_m \subset B_{n + m}$。
最后，{\it 分次 $A$-代数 $B$ 上的分次模 $M$}，由一个分次
$A$-模 $M$ 连同一个相容的 $B$-模结构给出，并满足
$B_n \cdot M_d \subset M_{n + d}$。现在可以定义
{\it 分次模或分次环的同态}、{\it 分次子模}、{\it 分次理想}、
{\it 分次模的正合列}等等。
\end{definition}

\begin{exercise}
\label{exercises-exercise-Euler-Poincare-field}
设 $A = k$ 为域。在此情形中，有限 $A$-模上所有可能的
Euler--Poincar\'e 函数是什么？
\end{exercise}

\begin{exercise}
\label{exercises-exercise-Euler-Poincare-Z}
设 $A ={\mathbf Z}$。在此情形中，有限 $A$-模上所有可能的
Euler--Poincar\'e 函数是什么？
\end{exercise}

\begin{exercise}
\label{exercises-exercise-Euler-Poincare-node}
设 $A = k[x, y]/(xy)$，其中 $k$ 代数闭。在此情形中，有限
$A$-模上所有可能的 Euler--Poincar\'e 函数是什么？
\end{exercise}

\begin{exercise}
\label{exercises-exercise-kernel-locally-finite}
假设 $A$ 是 Noether 环。证明局部有限分次 $A$-模之间的映射之核
仍局部有限。
\end{exercise}

\begin{exercise}
\label{exercises-exercise-no-hilbert}
设 $k$ 为域，$A = k$，并设 $B = k[x, y]$ 的分次由
$\deg(x) = 2$ 与 $\deg(y) = 3$ 决定。令
$\varphi(M) = \dim_k(M)$。计算 $B$ 作为分次 $k$-模的
Hilbert 函数。在此情形中有 Hilbert 多项式吗？
\end{exercise}

\begin{exercise}
\label{exercises-exercise-no-hilbert-or-is-there}
设 $k$ 为域，$A = k$，并设 $B = k[x, y]/(x^2, xy)$ 的分次由
$\deg(x) = 2$ 与 $\deg(y) = 3$ 决定。令
$\varphi(M) = \dim_k(M)$。计算 $B$ 作为分次 $k$-模的
Hilbert 函数。在此情形中有 Hilbert 多项式吗？
\end{exercise}

\begin{exercise}
\label{exercises-exercise-hilbert-to-compute}
设 $k$ 为域，$A = k$。令 $\varphi(M) = \dim_k(M)$。
固定 $d\in {\mathbf N}$。考察分次 $A$-代数
$B = k[x, y, z]/(x^d + y^d + z^d)$，其中 $x, y, z$
各自的次数都是 $1$。计算 $B$ 的 Hilbert 函数。在此情形中
有 Hilbert 多项式吗？
\end{exercise}



\section{环的 Proj}
\label{exercises-section-proj-ring}

\begin{definition}
\label{exercises-definition-homogeneous-ideal}
设 $R$ 为分次环。{\it 齐次}理想就是同时也是 $R$ 的分次子模的理想
$I \subset R$。等价地，它是由齐次元素生成的理想。再等价地，若
$f \in I$，并且
$$
f = f_0 + f_1 + \ldots + f_n
$$
是 $f$ 在 $R$ 中分解为齐次部分的分解，则对每个 $i$ 都有
$f_i \in I$。
\end{definition}

\begin{definition}
\label{exercises-definition-Proj-R}
定义分次环 $R$ 的{\it 齐次谱 $\text{Proj}(R)$}为 $R$ 的
齐次素理想 ${\mathfrak p}$ 所成的集合，其中要求
$R_{+} \not \subset {\mathfrak p}$。注意 $\text{Proj}(R)$
是 $\Spec(R)$ 的子集，因而带有自然的诱导拓扑。
\end{definition}

\begin{definition}
\label{exercises-definition-Dplus-Vplus}
设 $R = \oplus_{d \geq 0} R_d$ 为分次环，设 $f\in R_d$，并假设
$d \geq 1$。定义 {\it $R_{(f)}$} 为 $R_f$ 的子环，由所有形如
$r/f^n$ 的元素组成，其中 $r$ 齐次且 $\deg(r) = nd$。此外，定义
$$
D_{+}(f) = \{ {\mathfrak p} \in \text{Proj}(R) | f \not\in {\mathfrak p} \}.
$$
最后，对齐次理想 $I \subset R$，定义
$V_{+}(I) = V(I) \cap \text{Proj}(R)$。
\end{definition}

\begin{exercise}
\label{exercises-exercise-topology-proj}
关于 $\text{Proj}(R)$ 上的拓扑。使用上述定义与记号，证明下列命题。
\begin{enumerate}
\item 证明 $D_{+}(f)$ 在 $\text{Proj}(R)$ 中是开的。
\item 证明 $D_{+}(ff') = D_{+}(f) \cap D_{+}(f')$。
\item 设 $g = g_0 + \ldots + g_m$ 是 $R$ 的元素，其中
$g_i \in R_i$。用 $D_{+}(g_i)$（$i \geq 1$）以及
$D(g_0) \cap \text{Proj}(R)$ 表示
$D(g) \cap \text{Proj}(R)$。无需证明。
\item 设 $g\in R_0$ 为次数 $0$ 的齐次元素。用一族适当的
正次数齐次元素 $f_\alpha \in R$ 的 $D_{+}(f_\alpha)$
表示 $D(g) \cap \text{Proj}(R)$。
\item 证明开集族 $\{D_{+}(f)\}$ 构成 $\text{Proj}(R)$
拓扑的一组基。
\item
\label{exercises-item-bijection}
证明存在典范双射 $D_{+}(f) \to \Spec(R_{(f)})$。
% P12-ERR-0429
（提示：模仿 $\Spec$ 情形的证明，但在某一步加入一个理想的根。）
\item 证明（\ref{exercises-item-bijection}）中的映射是同胚。
\item 给出 $R$ 的例子，使得 $\text{Proj}(R)$ 不拟紧。无需证明。
\item 证明任意闭子集 $T \subset \text{Proj}(R)$ 都形如
$V_{+}(I)$，其中 $I \subset R$ 是某个齐次理想。
\end{enumerate}
\end{exercise}

\begin{remark}
\label{exercises-remark-continuous-proj-spec}
存在连续映射 $ \text{Proj}(R) \longrightarrow \Spec(R_0) $。
\end{remark}

\begin{exercise}
\label{exercises-exercise-iso-polynomial-ring-one-variable}
若 $R = A[X]$ 且 $\deg(X) = 1$，证明自然映射
$\text{Proj}(R) \to \Spec(A)$ 是双射，事实上是同胚。
\end{exercise}

\begin{exercise}
\label{exercises-exercise-blowing-up-I}
爆破：第一部分。
在本练习中，$R = Bl_I(A) = A \oplus I \oplus I^2 \oplus \ldots$。
考察自然映射 $b : \text{Proj}(R) \to \Spec(A)$。
令 $U = \Spec(A) - V(I)$。证明
$$
b : b^{-1}(U) \longrightarrow U
$$
是同胚。因此可以把 $U$ 看作 $\text{Proj}(R)$ 的开子集。
设 $Z \subset \Spec(A)$ 为不可约闭子概形，泛点为
$\xi \in Z$。假设 $\xi \not\in V(I)$；换言之，
$Z \not\subset V(I)$；换言之，$\xi \in U$；换言之，
$Z\cap U \not = \emptyset$。定义 $Z$ 的{\it 严格变换}
$Z'$ 为位于 $\xi$ 上方的唯一一点 $\xi'$ 的闭包。另一种说法是，
$Z'$ 是局部闭子集 $Z\cap U \subset U \subset \text{Proj}(R)$
在 $\text{Proj}(R)$ 中的闭包。
\end{exercise}

\begin{exercise}
\label{exercises-exercise-blowing-up-II}
爆破：第二部分。
设 $A = k[x, y]$，其中 $k$ 是域，并设 $I = (x, y)$。
设 $R$ 为 $A$ 关于 $I$ 的爆破代数。
\begin{enumerate}
\item 证明 $Z_1 = V(\{x\})$ 与 $Z_2 = V(\{y\})$ 的严格变换不交。
\item 证明 $Z_1 = V(\{x\})$ 与 $Z_2 = V(\{x-y^2\})$
的严格变换相交。
\item 找出理想 $J \subset A$，使得 $V(J) = V(I)$，并且
$Z_1 = V(\{x\})$ 与 $Z_2 = V(\{x-y^2\})$ 在沿 $J$
的爆破中的严格变换不交。
\end{enumerate}
\end{exercise}

\begin{exercise}
\label{exercises-exercise-proj-when-empty}
设 $R$ 为分次环。
\begin{enumerate}
% P12-ERR-0430
\item 证明若对所有 $n >> 0$ 都有 $R_n = (0)$，则
$\text{Proj}(R)$ 为空。
\item 证明若 $R$ 是整环且 $R_{+}$ 非零，则
$\text{Proj}(R)$ 是不可约拓扑空间。（回忆：空拓扑空间不可约。）
\end{enumerate}
\end{exercise}

\begin{exercise}
\label{exercises-exercise-blowing-up-III}
爆破：第三部分。
考察严格变换定义中的 $A$、$I$、$U$ 与 $Z$。设
$Z = V({\mathfrak p})$，其中 ${\mathfrak p}$ 为素理想。令
$\bar A = A/{\mathfrak p}$，并令 $\bar I$ 为 $I$ 在 $\bar A$
中的像。
\begin{enumerate}
\item 证明存在满环映射
% P12-ERR-0431
$R: = Bl_I(A) \to \bar R: = Bl_{\bar I}(\bar A)$。
\item 证明上述环映射诱导从 $\text{Proj}(\bar R)$ 到 $Z$
的严格变换 $Z'$ 上的双射。（这并不容易。
% P12-ERR-0432
提示：使用上面的 5(b)。）
\item 推出严格变换 $Z' = V_{+}(P)$，其中 $P \subset R$
是由 $P_d = I^d \cap {\mathfrak p}$ 定义的齐次理想。
\item 假设 $Z_1 = V({\mathfrak p})$ 与
$Z_2 = V({\mathfrak q})$ 是由素理想定义的不可约闭子集，并且
$Z_1 \not \subset Z_2$、$Z_2 \not \subset Z_1$。
证明爆破理想 $I = {\mathfrak p} + {\mathfrak q}$ 会分离
$Z_1$ 与 $Z_2$ 的严格变换；即
$Z_1' \cap Z_2' = \emptyset$。（提示：考察（c）中的齐次理想
$P$ 与 $Q$，并考察 $V(P + Q)$。）
\end{enumerate}
\end{exercise}


\section{一维 Cohen--Macaulay 环}
\label{exercises-section-CM-dim-1}

\begin{definition}
\label{exercises-definition-CM}
若 Noether 局部环 $A$ 的维数为 $d$，并且 $A$ 存在参数系
$x_1, \ldots, x_d$，使得对 $i = 1, \ldots, d$，$x_i$
在 $A/(x_1, \ldots, x_{i-1})$ 中都是非零因子，则称该环
为维数 $d$ 的{\it Cohen--Macaulay} 环。
\end{definition}

\begin{exercise}
\label{exercises-exercise-CM-dim-1-I}
一维 Cohen--Macaulay 环。第一部分：理论。
\begin{enumerate}
% P12-ERR-0433
\item 设 $(A, {\mathfrak m})$ 为局部 Noether 环，且
$\dim A = 1$。证明若 $x\in {\mathfrak m}$ 不是零因子，则
\begin{enumerate}
\item $\dim A/xA = 0$；换言之，$A/xA$ 是 Artin 环；再换言之，
$\{x\}$ 是 $A$ 的参数系。
% P12-ERR-0434
\item $A$ 没有嵌入素理想。
\end{enumerate}
\item 反过来，设 $(A, {\mathfrak m})$ 为维数为 $1$ 的局部 Noether 环。
证明若 $A$ 没有嵌入素理想，则 ${\mathfrak m}$ 中存在非零因子。
\end{enumerate}
\end{exercise}

\begin{exercise}
\label{exercises-exercise-CM-dim-1-II}
一维 Cohen--Macaulay 环。第二部分：例子。
\begin{enumerate}
\item 设 $A$ 为 $k[x, y]/(x^2, xy)$ 在 $(x, y)$ 处的局部环。
\begin{enumerate}
\item 证明 $A$ 的维数为 1。
\item 证明 ${\mathfrak m}\subset A$ 的每个元素都是零因子。
\item 找出 $z\in {\mathfrak m}$，使得 $\dim A/zA = 0$
（无需证明）。
\end{enumerate}
\item 设 $A$ 为 $k[x, y]/(x^2)$ 在 $(x, y)$ 处的局部环。
在 ${\mathfrak m}$ 中找出一个非零因子（无需证明）。
\end{enumerate}
\end{exercise}

\begin{exercise}
\label{exercises-exercise-embedding-dim-1}
嵌入维数为 $1$ 的局部环。
假设 $(A, {\mathfrak m}, k)$ 是嵌入维数为 $1$ 的 Noether 局部环，
即
$$
\dim_k {\mathfrak m}/{\mathfrak m}^2 = 1.
$$
证明函数 $f(n) = \dim_k {\mathfrak m}^n/{\mathfrak m}^{n + 1}$
或者恒等于 $1$，或者其值依次为
$$
1, 1, \ldots, 1, 0, 0, 0, 0, 0, \ldots
$$
\end{exercise}

\begin{exercise}
\label{exercises-exercise-regular-local-dim-1}
维数为 $1$ 的正则局部环。
假设 $(A, {\mathfrak m}, k)$ 是维数为 $1$ 的正则 Noether 局部环。
回忆：这意味着 $A$ 的维数为 $1$，嵌入维数也为 $1$，即
$$
\dim_k {\mathfrak m}/{\mathfrak m}^2 = 1.
$$
任取 $x\in{\mathfrak m}$，使其在
${\mathfrak m}/{\mathfrak m}^2$ 中的类非零。
\begin{enumerate}
\item 证明对 ${\mathfrak m}$ 的每个元素 $y$，都存在整数 $n$，
使得 $y$ 可写成 $y = ux^n$，其中 $u\in A^\ast$ 是单位。
\item 证明 $x$ 在 $A$ 中是非零因子。
\item 推出 $A$ 是整环。
\end{enumerate}
\end{exercise}

\begin{exercise}
\label{exercises-exercise-nonzerodivisor-graded}
设 $(A, {\mathfrak m}, k)$ 为 Noether 局部环，其伴随分次环为
$Gr_{\mathfrak m}(A)$。
\begin{enumerate}
\item 假设 $x\in {\mathfrak m}^d$ 映到
$Gr_{\mathfrak m}(A)$ 的次数 $d$ 部分中的非零因子
$\bar x \in {\mathfrak m}^d/{\mathfrak m}^{d + 1}$。
证明 $x$ 是非零因子。
\item 假设 $A$ 的深度至少为 $1$；也就是说，假设存在非零因子
$y \in {\mathfrak m}$。在此情形中可以加强结论：只假设
$x\in {\mathfrak m}^d$ 映到 $Gr_{\mathfrak m}(A)$ 的次数 $d$
部分中的元素 $\bar x \in {\mathfrak m}^d/{\mathfrak m}^{d + 1}$，
且它在充分高的次数上是非零因子：即 $\exists N$，使得对所有
$n \geq N$，乘以 $\bar x$ 的映射
$$
{\mathfrak m}^n/{\mathfrak m}^{n + 1} \longrightarrow
{\mathfrak m}^{n + d}/{\mathfrak m}^{n + d + 1}
$$
都是单射。证明此时 $x$ 是非零因子。
\end{enumerate}
\end{exercise}

\begin{exercise}
\label{exercises-exercise-embedding-2-dim-1}
假设 $(A, {\mathfrak m}, k)$ 是维数为 $1$ 的 Noether 局部环。
再假设 $A$ 的嵌入维数为 $2$，即
$$
\dim_k {\mathfrak m}/{\mathfrak m}^2 = 2.
$$
记号：$f(n) = \dim_k {\mathfrak m}^n/{\mathfrak m}^{n + 1}$。
选取生成元 $x, y \in {\mathfrak m}$，并对某个齐次理想 $I$ 写成
$Gr_{\mathfrak m}(A) = k[\bar x, \bar y]/I$。
\begin{enumerate}
\item 证明存在齐次元素 $F\in k[\bar x, \bar y]$，使得
$I \subset (F)$，并且在所有充分高的次数上取等号。
\item 证明 $f(n) \leq n + 1$。
\item 证明若 $f(n) < n + 1$，则 $n \geq \deg(F)$。
\item 证明若 $f(n) < n + 1$，则 $f(n + 1) \leq f(n)$。
% P12-ERR-0443
\item 证明对所有 $n >> 0$，都有 $f(n) = \deg(F)$。
\end{enumerate}
\end{exercise}

\begin{exercise}
\label{exercises-exercise-CM-dim-1-embedding-dim-2}
维数为 1、嵌入维数为 2 的 Cohen--Macaulay 环。
假设 $(A, {\mathfrak m}, k)$ 是维数为 $1$ 的 Cohen--Macaulay
Noether 局部环。再假设 $A$ 的嵌入维数为 $2$，即
$$
\dim_k {\mathfrak m}/{\mathfrak m}^2 = 2.
$$
记号 $f$、$F$、$x, y\in {\mathfrak m}$、$I$ 同练习
\ref{exercises-exercise-embedding-2-dim-1}。请自由使用上面各题的结果。
\begin{enumerate}
\item 假设 $z\in {\mathfrak m}$ 的类在
${\mathfrak m}/{\mathfrak m}^2$ 中是线性形式
$\alpha \bar x + \beta \bar y \in k[\bar x, \bar y]$，并且它与
$F$ 互素。
\begin{enumerate}
\item 证明 $z$ 在 $A$ 上是非零因子。
\item 令 $d = \deg(F)$。证明对所有充分大的 $n$ 都有
${\mathfrak m}^n = z^{n + 1-d}{\mathfrak m}^{d-1}$。
（提示：先利用关于 $Gr_{\mathfrak m}(A)$ 的已知结果，证明
$z^{n + 1-d}{\mathfrak m}^{d-1} \to
{\mathfrak m}^n/{\mathfrak m}^{n + 1}$ 是满射，再使用 NAK。）
\end{enumerate}
\item $k$ 满足什么条件才能保证存在这样的 $z$？
（无需证明；这太容易了。）

\noindent
现在假设存在上述 $z$。事实证明这是无害的假设（其含义是：
为了得到下文（d）和（e）的结果，可以化归到它成立的情形）。
\item 证明对所有 $\ell \geq d$ 都有
${\mathfrak m}^\ell = z^{\ell - d + 1} {\mathfrak m}^{d-1}$。
\item 推出 $I = (F)$。
% P12-ERR-0435
\item 推出函数 $f$ 的值依次为
$$
2, 3, 4, \ldots, d-1, d, d, d, d, d, d, d, \ldots
$$
\end{enumerate}
\end{exercise}

\begin{remark}
\label{exercises-remark-CM-dim-1-embedding-dim-2}
这表明，维数为 1、嵌入维数为 2 的 Cohen--Macaulay 局部 Noether 环
应当形如 $B/FB$，其中 $B$ 是二维正则局部环。在环具有适当“良好性”
的条件下，这基本上是正确的。
\end{remark}






\section{无穷多个素数}
\label{exercises-section-many-primes}

\noindent
本节收集了一些关于无穷多个素数不可逆之环的奇特问题。

\begin{exercise}
\label{exercises-exercise-not-in-Q}
给出具有下列两个性质的有限型 ${\mathbf Z}$-代数 $R$：
\begin{enumerate}
\item 不存在环映射 $R \to {\mathbf Q}$。
\item 对每个素数 $p$，都存在极大理想 ${\mathfrak m} \subset R$，
使得 $R/{\mathfrak m} \cong {\mathbf F}_p$。
\end{enumerate}
\end{exercise}

\begin{exercise}
\label{exercises-exercise-strange-fp-1}
对 $f \in {\mathbf Z}[x, u]$，定义
$f_p(x) = f(x, x^p) \bmod p \in {\mathbf F}_p[x]$。
给出满足下列两个性质的 $f \in {\mathbf Z}[x, u]$：
\begin{enumerate}
\item 有无穷多个 $p$，使 $f_p$ 在 ${\mathbf F}_p$ 中没有零点。
% P12-ERR-0443
\item 对所有 $p >> 0$，多项式 $f_p$ 都有一次因子或二次因子。
\end{enumerate}
\end{exercise}

\begin{exercise}
\label{exercises-exercise-strange-fp-2}
对 $f \in {\mathbf Z}[x, y, u, v]$，定义
$f_p(x, y) = f(x, y, x^p, y^p) \bmod p \in {\mathbf F}_p[x, y]$。
给出一个“有趣”的 $f$，使得对所有
% P12-ERR-0443
$p >> 0$，$f_p$ 都可约。例如，$f = xv-yu$ 对应
$f_p = xy^p-x^py = xy(x^{p-1}-y^{p-1})$，这是“无趣”的；
任何仅依赖 $x, u$ 的 $f$ 也是“无趣”的；等等。
\end{exercise}

\begin{remark}
\label{exercises-remark-strange-fp}
设 $h \in {\mathbf Z}[y]$ 为次数 $d$ 的首一多项式。那么：
\begin{enumerate}
\item 映射 $A = {\mathbf Z}[x] \to B ={\mathbf Z}[y]$、
$x \mapsto h$ 是秩为 $d$ 的有限局部自由映射。
% P12-ERR-0436
\item 对所有素数 $p$，映射
$A_p = {\mathbf F}_p[x]\to B_p = {\mathbf F}_p[y]$、
$y \mapsto h(y) \bmod p$ 都是秩为 $d$ 的有限局部自由映射。
\end{enumerate}
\end{remark}

\begin{exercise}
\label{exercises-exercise-strange-fp-3}
设 $h, A, B, A_p, B_p$ 同上注。对 $f \in {\mathbf Z}[x, u]$，
定义 $f_p(x) = f(x, x^p) \bmod p \in {\mathbf F}_p[x]$。
对 $g \in {\mathbf Z}[y, v]$，定义
$g_p(y) = g(y, y^p) \bmod p \in {\mathbf F}_p[y]$。
\begin{enumerate}
\item 给出 $h$ 和 $g$ 的例子，使得不存在满足下述性质的 $f$：
% P12-ERR-0437
$$
f_p  =  Norm_{B_p/A_p}(g_p).
$$
\item 证明对上述任意 $h$ 和 $g$，都存在非零 $f$，使得对所有
$p$ 都有
$$
Norm_{B_p/A_p}(g_p)\quad\text{整除}\quad f_p .
$$
如果愿意，可以只处理 $h = y^n$ 甚至 $n = 2$ 的情形，
不过一般情形也成立。
% P12-ERR-0438
\item 讨论这与前一题组的练习 6 和 7 有何关系。
\end{enumerate}
\end{exercise}

\begin{exercise}
\label{exercises-exercise-strange-fp-unsolved}
未解决的问题。它们可能非常困难，也可能很容易。我不知道。
\begin{enumerate}
\item 是否存在 $f \in {\mathbf Z}[x, u]$，使得 $f_p$ 对无穷多个
$p$ 都不可约？（提示：是的，$f(x, u) = u - x - 1$ 会发生这种情形，
$f(x, u) = u^2 - x^2 + 1$ 也会。）
\item 设 $f \in {\mathbf Z}[x, u]$ 非零，并假设对所有充分大的
$p$，都有 $\deg_x(f_p) = dp + d'$。（换言之，
$\deg_u(f) = d$，且 $f$ 中 $u^d$ 的系数 $c$ 满足
$\deg_x(c) = d'$。）假设可以写成 $d = d_1 + d_2$ 和
$d' = d'_1 + d'_2$，其中 $d_1, d_2 > 0$ 且
$d'_1, d'_2 \geq 0$；并且对所有充分大的 $p$，存在分解
$$
f_p = f_{1, p} f_{2, p}
$$
满足 $\deg_x(f_{1, p}) = d_1p + d'_1$。% P12-ERR-0439
$f$ 是否必定来自练习 4
中的范数构造？（更准确地说，是否存在 $h$ 和 $g$，使得
% P12-ERR-0437
$Norm_{B_p/A_p}(g_p)$ 对所有
% P12-ERR-0443
$p >> 0$ 都整除 $f_p$？）
\item 与（b）类似的问题，但现在假设
$f \in {\mathbf Z}[x_1, x_2, u_1, u_2]$ 不可约，并且只假设
$f_p(x_1, x_2) = f(x_1, x_2, x_1^p, x_2^p) \bmod p$ 对所有
% P12-ERR-0443
$p >> 0$ 都可分解。
\end{enumerate}
\end{exercise}


















\section{滤过导出范畴}
\label{exercises-section-filtered-derived}

\noindent
为完成本节练习，请阅读《同调代数》第
\ref{homology-section-filtrations} 节的材料。我们说 $A$ 是
$\mathcal{A}$ 的滤过对象，是指 $A$ 带有一个滤过 $F$，
但在记号中略去它。

\begin{exercise}
\label{exercises-exercise-split-injective}
设 $\mathcal{A}$ 为 Abel 范畴，$I$ 为 $\mathcal{A}$ 的滤过对象。
假设 $I$ 上的滤过有限，并且每个 $\text{gr}^p(I)$ 都是
$\mathcal{A}$ 的内射对象。证明存在同构
$I \cong \bigoplus \text{gr}^p(I)$，使得滤过 $F^p(I)$ 对应于
% P12-ERR-0440
$\bigoplus_{p' \geq p} \text{gr}^p(I)$。
\end{exercise}

\begin{exercise}
\label{exercises-exercise-filtered-injective}
设 $\mathcal{A}$ 为 Abel 范畴，$I$ 为 $\mathcal{A}$ 的滤过对象。
假设 $I$ 上的滤过有限。证明下列条件等价：
\begin{enumerate}
\item 对任意滤过对象的实线图
$$
\xymatrix{
A \ar[r]_\alpha \ar[d] & B \ar@{-->}[ld] \\
I &
}
$$
若（\romannumeral1）$A$ 与 $B$ 上的滤过有限，并且
% P12-ERR-0441
（\romannumeral2）$\text{gr}(\alpha)$ 为单射，则存在使图交换的虚线箭头。
\item 每个 $\text{gr}^p I$ 都是内射对象。
\end{enumerate}
\end{exercise}

\noindent
注意，给定有限滤过对象间的态射 $\alpha : A \to B$，
% P12-ERR-0441
$\text{gr}(\alpha)$ 为单射，等价于 $\alpha$ 在范畴
$\text{Fil}(\mathcal{A})$ 中是{\it 严格单态射}。事实上，成为单态射
意味着 $\Ker(\alpha) = 0$，而严格意味着这还蕴含
$\Ker(\text{gr}(\alpha)) = 0$。见《同调代数》引理
\ref{homology-lemma-characterize-strict}。（我们只把“单射”一词用于
Abel 范畴中的态射，尽管它在任何具有核的加性范畴中都有意义。）
% P12-ERR-0442
上面的练习说明了下述定义的合理性。

\begin{definition}
\label{exercises-definition-injective-filtered}
设 $\mathcal{A}$ 为 Abel 范畴，$I$ 为 $\mathcal{A}$ 的滤过对象。
假设 $I$ 上的滤过有限。若每个 $\text{gr}^p(I)$ 都是
$\mathcal{A}$ 的内射对象，则称 $I$ {\it 滤过内射}。
\end{definition}

\noindent
作如下定义，以免处处重复“具有有限滤过”。

\begin{definition}
\label{exercises-definition-finite-filtration-category}
设 $\mathcal{A}$ 为 Abel 范畴。记 $\text{Fil}(\mathcal{A})$
的全子范畴 {\it $\text{Fil}^f(\mathcal{A})$} 为这样的对象
$A \in \Ob(\text{Fil}(\mathcal{A}))$ 所成的范畴：其滤过有限。
\end{definition}

\begin{exercise}
\label{exercises-exercise-inject-into-injective}
设 $\mathcal{A}$ 为 Abel 范畴，并假设 $\mathcal{A}$ 有足够多内射对象。
设 $A$ 为 $\text{Fil}^f(\mathcal{A})$ 的对象。证明存在严格单态射
$\alpha : A \to I$，把 $A$ 嵌入
$\text{Fil}^f(\mathcal{A})$ 的滤过内射对象 $I$。
\end{exercise}

\begin{definition}
\label{exercises-definition-filtered-quasi-isomorphism}
设 $\mathcal{A}$ 为 Abel 范畴，并设
$\alpha : K^\bullet \to L^\bullet$ 为
$\text{Fil}(\mathcal{A})$ 中复形的态射。若对每个
$p \in \mathbf{Z}$，态射
$\text{gr}^p(K^\bullet) \to \text{gr}^p(L^\bullet)$
都是拟同构，则称 $\alpha$ 为{\it 滤过拟同构}。
\end{definition}

\begin{definition}
\label{exercises-definition-filtered-acyclic}
设 $\mathcal{A}$ 为 Abel 范畴，并设 $K^\bullet$ 为
$\text{Fil}^f(\mathcal{A})$ 中的复形。若对每个
$p \in \mathbf{Z}$，复形 $\text{gr}^p(K^\bullet)$ 都无环，
则称 $K^\bullet$ {\it 滤过无环}。
\end{definition}

\begin{exercise}
\label{exercises-exercise-filtered-quasi-isomorphism}
设 $\mathcal{A}$ 为 Abel 范畴，并设
$\alpha : K^\bullet \to L^\bullet$ 为
$\text{Fil}^f(\mathcal{A})$ 中下有界复形的态射。
（注意上标 $f$。）证明下列条件等价：
\begin{enumerate}
\item $\alpha$ 是滤过拟同构；
\item 对每个 $p \in \mathbf{Z}$，映射
$\alpha : F^pK^\bullet \to F^pL^\bullet$ 是拟同构；
\item 对每个 $p \in \mathbf{Z}$，映射
$\alpha : K^\bullet/F^pK^\bullet \to L^\bullet/F^pL^\bullet$
是拟同构；
\item $\alpha$ 的锥（见《导出范畴》定义
\ref{derived-definition-cone}）是滤过无环复形。
\end{enumerate}
此外，证明若 $\alpha$ 是滤过拟同构，则 $\alpha$ 也是通常的拟同构。
\end{exercise}

\begin{exercise}
\label{exercises-exercise-injective-resolution}
设 $\mathcal{A}$ 为 Abel 范畴，并假设 $\mathcal{A}$ 有足够多内射对象。
设 $A$ 为 $\text{Fil}^f(\mathcal{A})$ 的对象。证明存在
$\text{Fil}^f(\mathcal{A})$ 的复形 $I^\bullet$，以及态射
$A[0] \to I^\bullet$，使得
\begin{enumerate}
\item 每个 $I^p$ 都滤过内射；
\item 当 $p < 0$ 时 $I^p = 0$；
\item $A[0] \to I^\bullet$ 是滤过拟同构。
\end{enumerate}
\end{exercise}

\begin{exercise}
\label{exercises-exercise-injective-resolution-complex}
设 $\mathcal{A}$ 为 Abel 范畴，并假设 $\mathcal{A}$ 有足够多内射对象。
设 $K^\bullet$ 为 $\text{Fil}^f(\mathcal{A})$ 的对象所成的
下有界复形。证明存在滤过拟同构
$\alpha : K^\bullet \to I^\bullet$，其中 $I^\bullet$ 是
$\text{Fil}^f(\mathcal{A})$ 的复形，其各项 $I^n$ 滤过内射，
并且它下有界。事实上，可以选取 $\alpha$，使每个 $\alpha^n$
都是严格单态射。
\end{exercise}

\begin{exercise}
\label{exercises-exercise-morphisms-lift}
设 $\mathcal{A}$ 为 Abel 范畴。考察实线图
$$
\xymatrix{
K^\bullet \ar[r]_\alpha \ar[d]_\gamma & L^\bullet \ar@{-->}[dl]^\beta \\
I^\bullet
}
$$
其中各对象是 $\text{Fil}^f(\mathcal{A})$ 中的复形。假设
$K^\bullet$、$L^\bullet$、$I^\bullet$ 下有界，并且每个
$I^n$ 都是滤过内射对象。再假设 $\alpha$ 是滤过拟同构。
\begin{enumerate}
\item 存在复形的映射 $\beta$，使图在同伦意义下交换。
\item 若 $\alpha$ 在每个次数上都是严格单态射，则可以找到
使图严格交换的 $\beta$。
\end{enumerate}
\end{exercise}

\begin{exercise}
\label{exercises-exercise-acyclic-is-zero}
设 $\mathcal{A}$ 为 Abel 范畴。
% P12-ERR-0444
设 $K^\bullet$、$K^\bullet$ 为 $\text{Fil}^f(\mathcal{A})$ 中的复形。
假设
\begin{enumerate}
\item $K^\bullet$ 下有界且滤过无环；
\item $I^\bullet$ 下有界，并且由滤过内射对象组成。
\end{enumerate}
则任意态射 $K^\bullet \to I^\bullet$ 都同伦于零。
\end{exercise}

\begin{exercise}
\label{exercises-exercise-morphisms-equal-up-to-homotopy}
设 $\mathcal{A}$ 为 Abel 范畴。考察实线图
$$
\xymatrix{
K^\bullet \ar[r]_\alpha \ar[d]_\gamma & L^\bullet \ar@{-->}[dl]^{\beta_i} \\
I^\bullet
}
$$
其中各对象是 $\text{Fil}^f(\mathcal{A})$ 中的复形。
假设 $K^\bullet$、$L^\bullet$、$I^\bullet$ 下有界，每个 $I^n$
都是滤过内射对象，并且 $\alpha$ 是滤过拟同构。任何两个使图在
同伦意义下交换的态射 $\beta_1, \beta_2$ 彼此同伦。
\end{exercise}







\section{正则函数}
\label{exercises-section-regular-functions}

\begin{exercise}
\label{exercises-exercise-extra-function}
考察由 $\mathbf{C}^2$ 中方程 $t^2 = s^5 + 8$ 给出的仿射曲线
$X$，坐标为 $s, t$。设 $x \in X$ 为坐标 $(1, 3)$ 的点。
令 $U = X \setminus \{x\}$。证明 $U$ 上存在一个正则函数，
它不是 $\mathbf{C}^2$ 上正则函数的限制；也就是说，它不是
$s$ 与 $t$ 的多项式在 $U$ 上的限制。
\end{exercise}

\begin{exercise}
\label{exercises-exercise-no-extra-function}
设 $n \geq 2$，并设 $E \subset \mathbf{C}^n$ 为有限子集。
证明 $\mathbf{C}^n \setminus E$ 上的任意正则函数都是多项式。
\end{exercise}

\begin{exercise}
\label{exercises-exercise-cone}
设 $X \subset \mathbf{C}^n$ 为仿射簇。若
$x = (a_1, \ldots, a_n) \in X$ 且 $\lambda \in \mathbf{C}$
蕴含 $(\lambda a_1, \ldots, \lambda a_n) \in X$，则称 $X$
为{\it 锥}。当然，若 $\mathfrak p \subset \mathbf{C}[x_1, \ldots, x_n]$
是由 $x_1, \ldots, x_n$ 的齐次多项式生成的素理想，则仿射簇
$X = V(\mathfrak p) \subset \mathbf{C}^n$ 是锥。证明反过来，
与锥对应的素理想 $\mathfrak p \subset \mathbf{C}[x_1, \ldots, x_n]$
也可由 $x_1, \ldots, x_n$ 的齐次多项式生成。
\end{exercise}

\begin{exercise}
\label{exercises-exercise-extra-function-cone}
给出一个为锥的仿射簇 $X \subset \mathbf{C}^n$（见练习
\ref{exercises-exercise-cone}），以及 $U = X \setminus \{(0, \ldots, 0)\}$
上的正则函数 $f$，使它不是 $\mathbf{C}^n$ 上多项式函数的限制。
\end{exercise}

\begin{exercise}
\label{exercises-exercise-regular-functions}
本练习考察非代数闭域上的正则函数。设 $k$ 为域。
设 $Z \subset k^n$ 为 Zariski 局部闭子集；即存在理想
$I \subset J \subset k[x_1, \ldots, x_n]$，使得
$$
Z = \{a \in k^n \mid
f(a) = 0\ \forall\ f \in I,\ \exists\ g \in J,\ g(a) \not = 0\}.
$$
若对每个 $z \in Z$，都存在 Zariski 开邻域
$z \in U \subset Z$ 和多项式 $f, g \in k[x_1, \ldots, x_n]$，
使得对所有 $u \in U$ 都有 $g(u) \not = 0$，并且对所有
$u \in U$ 都有 $\varphi(u) = f(u)/g(u)$，则称函数
$\varphi : Z \to k$ {\it 正则}。
\begin{enumerate}
\item 若 $k = \bar k$ 且 $Z = k^n$，证明正则函数由多项式给出。
（只在以前未见过这一论证时完成此题。）
\item 若 $k$ 有限，证明（a）每个函数 $\varphi$ 都正则；
（b）正则函数环在 $k$ 上有限维。（如果愿意，可以取
$Z = k^n$，甚至取 $n = 1$。）
\item 若 $k = \mathbf{R}$，给出 $Z = \mathbf{R}$ 上不由多项式
给出的正则函数。
\item 若 $k = \mathbf{Q}_p$，给出 $Z = \mathbf{Q}_p$ 上不由
多项式给出的正则函数。
\end{enumerate}
\end{exercise}






\section{层}
\label{exercises-section-sheaves}

\noindent
范畴 $\mathcal{C}$ 的态射 $f : X \to Y$ 称为{\it 单态射}，
是指对每一对态射 $a, b : W \to X$，都有
$f \circ a = f \circ b \Rightarrow a = b$。集合范畴中的单态射
就是集合的单射。

\begin{exercise}
\label{exercises-exercise-mono-sheaves-sets}
% P12-ERR-0445
仔细证明：集合层的映射是单态射（在集合层范畴中），当且仅当它在
所有茎上诱导的映射都是单射。
\end{exercise}

\noindent
范畴 $\mathcal{C}$ 的态射 $f : X \to Y$ 称为{\it 同构}，
是指存在态射 $g : Y \to X$，使得
$f \circ g = \text{id}_Y$ 且 $g \circ f = \text{id}_X$。
集合范畴中的同构就是集合的双射。

\begin{exercise}
\label{exercises-exercise-isomorphism-sheaves-sets}
仔细证明：集合层的映射是同构（在集合层范畴中），当且仅当它在
所有茎上诱导的映射都是双射。
\end{exercise}

\noindent
范畴 $\mathcal{C}$ 的态射 $f : X \to Y$ 称为{\it 满态射}，
是指对每一对态射 $a, b : Y \to Z$，都有
$a \circ f = b \circ f \Rightarrow a = b$。集合范畴中的满态射
就是集合的满射。

\begin{exercise}
\label{exercises-exercise-epi-sheaves-sets}
仔细证明：集合层的映射是满态射（在集合层范畴中），当且仅当它在
所有茎上诱导的映射都是满射。
\end{exercise}

\begin{exercise}
\label{exercises-exercise-adjoint-push-pull}
设 $f : X \to Y$ 为拓扑空间的映射。证明{\bf 集合}层的推前 $f_\ast$
与拉回 $f^{-1}$ 构成一对伴随函子。
\end{exercise}

\begin{exercise}
\label{exercises-exercise-j-shriek}
设 $j : U \to X$ 为开浸入。证明：
\begin{enumerate}
\item 拉回 $j^{-1} : \Sh(X) \to \Sh(U)$ 有左伴随
$j_{!} : \Sh(U) \to \Sh(X)$，称为{\it 以空集延拓}。
\item 刻画 $j_{!}({\mathcal G})$ 的茎，其中
$\mathcal{G} \in \Sh(U)$。
\item 拉回 $j^{-1} : \textit{Ab}(X) \to \textit{Ab}(U)$ 有左伴随
$j_{!} : \textit{Ab}(U) \to \textit{Ab}(X)$，称为{\it 零延拓}。
\item 刻画 $j_{!}({\mathcal G})$ 的茎，其中
$\mathcal{G} \in \textit{Ab}(U)$。
\end{enumerate}
注意，零延拓与以空集延拓不同！
\end{exercise}

\begin{exercise}
\label{exercises-exercise-not-locally-generated-by-sections}
令 $X = \mathbf{R}$，取通常拓扑。令
$\mathcal{O}_X = \underline{\mathbf{Z}/2\mathbf{Z}}_X$。
令 $i : Z = \{0\} \to X$ 为包含映射，并令
$\mathcal{O}_Z = \underline{\mathbf{Z}/2\mathbf{Z}}_Z$。
证明下列各点（前三点直接来自定义；若定义尚不清楚，应借此阐明）：
\begin{enumerate}
\item $i_*\mathcal{O}_Z$ 是摩天层。
\item 存在典范满射
$\underline{\mathbf{Z}/2\mathbf{Z}}_X \to
i_*\underline{\mathbf{Z}/2\mathbf{Z}}_Z$。记其核为
$\mathcal{I} \subset \mathcal{O}_X$。
\item $\mathcal{I}$ 是 $\mathcal{O}_X$ 的理想层。
\item $X$ 上的层 $\mathcal{I}$ 不可能由截面局部生成
（含义见《模》定义 \ref{modules-definition-locally-generated}）。
\end{enumerate}
\end{exercise}

\begin{exercise}
\label{exercises-exercise-quotient-j-shriek-Z}
设 $X$ 为拓扑空间，${\mathcal F}$ 为 $X$ 上的 Abel 群层。
证明 ${\mathcal F}$ 是某个（可能非常大的）层直和的商，其中每一项
都形如
$$
j_{!}(\underline{{\mathbf Z}}_U)
$$
这里 $U \subset X$ 开，$\underline{{\mathbf Z}}_U$ 表示 $U$
上取值为 ${\mathbf Z}$ 的常值层。
\end{exercise}

\begin{remark}
\label{exercises-remark-direct-sum-stalk-abelian}
设 $X$ 为拓扑空间。在 Abel 群层范畴中，一族层
$\{{\mathcal F}_i\}_{i\in I}$ 的直和，是预层
% P12-ERR-0446
$U \mapsto \oplus {\mathcal F}_i(U)$ 的层化。因此直和在点
$x$ 处的茎，是各 ${\mathcal F}_i$ 在 $x$ 处之茎的直和。
\end{remark}

\begin{exercise}
\label{exercises-exercise-product-over-points}
设 $X$ 为拓扑空间，并给定由 $x \in X$ 标号的一族 Abel 群
$A_x$。证明规则 $U \mapsto \prod_{x \in U} A_x$ 配以显然的限制映射，
定义 Abel 群层 $\mathcal{G}$。用例子说明，通常对
$x \in X$ 并没有 $\mathcal{G}_x = A_x$。
\end{exercise}

\begin{exercise}
\label{exercises-exercise-modified-product-over-points}
设 $X$、$A_x$、$\mathcal{G}$ 同练习
\ref{exercises-exercise-product-over-points}。设 $\mathcal{B}$ 为 $X$
拓扑的一组基；见《拓扑》定义 \ref{topology-definition-base}。
对 $U \in \mathcal{B}$，设 $A_U$ 为子群
$A_U \subset \mathcal{G}(U) = \prod_{x \in U} A_x$。假设当
$U \subset V$ 且 $U, V \in \mathcal{B}$ 时，限制映射把
$A_V$ 映入 $A_U$。对开集 $U \subset X$，令
$$
\mathcal{F}(U) =
\left\{
(s_x)_{x \in U}
\middle|
\begin{matrix}
\text{ 对每个 }x\text{ 属于 }U\text{，存在 } V \in \mathcal{B} \\
x \in V \subset U\text{，使得 } (s_y)_{y \in V} \in A_V
\end{matrix}
\right\}
$$
证明 $\mathcal{F}$ 定义 $X$ 上的 Abel 群层。用例子说明，
通常对 $U \in \mathcal{B}$ 并没有 $\mathcal{F}(U) = A_U$。
\end{exercise}

\begin{exercise}
\label{exercises-exercise-exact-but-not-a-stalk-functor}
给出拓扑空间 $X$ 和函子
$$
F : \Sh(X) \longrightarrow \textit{Sets}
$$
的例子，使该函子正合（与有限积、等化子、有限余积及余等化子交换；
见《范畴》 \ref{categories-section-exact-functor} 节），但不存在
$x \in X$，使 $F$ 同构于茎函子
$\mathcal{F} \mapsto \mathcal{F}_x$。
\end{exercise}




\section{概形}
\label{exercises-section-schemes}

\noindent
令 $LRS$ 为局部环空间范畴。仿射概形是 $LRS$ 中与某个形如
$(\Spec(A), {\mathcal O}_{\Spec(A)})$ 的序对在 $LRS$ 中同构的对象。
概形是 $LRS$ 的对象 $(X, {\mathcal O}_X)$，并满足：每个点
$x\in X$ 都有开邻域 $U \subset X$，使得序对
$(U, {\mathcal O}_X|_U)$ 是仿射概形。

\begin{exercise}
\label{exercises-exercise-one-point}
找出一个不是概形的 $1$-点局部环空间。
\end{exercise}

\begin{exercise}
\label{exercises-exercise-two-points}
假设概形 $X$ 的底拓扑空间有 2 个点。证明 $X$ 是仿射概形。
\end{exercise}

\begin{exercise}
\label{exercises-exercise-discrete-finite-set-points}
假设概形 $X$ 的底拓扑空间是有限离散集。证明 $X$ 是仿射概形。
\end{exercise}

\begin{exercise}
\label{exercises-exercise-three-points}
证明存在具有三个点的非仿射概形。
\end{exercise}

\begin{exercise}
\label{exercises-exercise-quasi-compact-closed-point}
假设 $X$ 是非空拟紧概形。证明 $X$ 有闭点。
\end{exercise}

\begin{remark}
\label{exercises-remark-open-immersion}
若 $(X, {\mathcal O}_X)$ 是环空间，且 $U \subset X$ 是开子集，
则 $(U, {\mathcal O}_X|_U)$ 是环空间。记
${\mathcal O}_U = {\mathcal O}_X|_U$。
% P12-ERR-0447
存在典范的环空间态射
$$
j : (U, {\mathcal O}_U) \longrightarrow (X, {\mathcal O}_X).
$$
若 $(X, {\mathcal O}_X)$ 是局部环空间，则
$(U, {\mathcal O}_U)$ 也是局部环空间，且 $j$ 是局部环空间态射。
若 $(X, {\mathcal O}_X)$ 是概形，则 $(U, {\mathcal O}_U)$
也是概形，且 $j$ 是概形态射。称 $(U, {\mathcal O}_U)$ 为
$(X, {\mathcal O}_X)$ 的{\it 开子概形}，并称 $j$ 为{\it 开浸入}。
更一般地，任何态射
$j' : (V, {\mathcal O}_V) \to (X, {\mathcal O}_X)$，若与上述某个
态射 $j : (U, {\mathcal O}_U) \to (X, {\mathcal O}_X)$
{\it 同构}，也称为开浸入。
\end{remark}

\begin{exercise}
\label{exercises-exercise-open-affine-not-affine}
给出仿射概形 $(X, {\mathcal O}_X)$ 以及开集 $U \subset X$ 的例子，
使得
% P12-ERR-0448
$(U, {\mathcal O}_X|U)$ 不是仿射概形。
\end{exercise}

\begin{exercise}
\label{exercises-exercise-morphism-does-not-extend}
% P12-ERR-0449
给出一对仿射概形 $(X, {\mathcal O}_X)$、$(Y, {\mathcal O}_Y)$，
$X$ 的开子概形 $(U, {\mathcal O}_X|_U)$，以及概形态射
$(U, {\mathcal O}_X|_U) \to (Y, {\mathcal O}_Y)$，使它不能延拓为
概形态射 $(X, {\mathcal O}_X) \to (Y, {\mathcal O}_Y)$。
\end{exercise}

\begin{exercise}
\label{exercises-exercise-closed-subscheme-does-not-extend}
（这相当困难。）
% P12-ERR-0450
给出概形 $X$、开子概形 $U \subset X$ 和闭子概形 $Z \subset U$
的例子，使 $Z$ 不能延拓为 $X$ 的闭子概形。
\end{exercise}

\begin{exercise}
\label{exercises-exercise-not-morphism-schemes}
给出概形 $X$、域 $K$ 以及环空间态射
$\Spec(K) \to X$ 的例子，使该态射不是概形态射。
\end{exercise}

\begin{exercise}
\label{exercises-exercise-just-kidding}
完成 \cite[Chapter II]{H} 第 1 节与第 2 节中的所有练习……
\ \ 开玩笑！
\end{exercise}

\begin{definition}
\label{exercises-definition-integral}
若概形 $X$ 非空，并且对每个非空仿射开集 $U \subset X$，环
$\Gamma(U, \mathcal{O}_X) = \mathcal{O}_X(U)$ 都是整环，
则称 $X$ {\it 整}。
\end{definition}

\begin{exercise}
\label{exercises-exercise-morphism-integral-schemes-surjective-stalks-not-closed}
给出{\it 整}概形的态射 $f : X \to Y$ 的例子，使得对所有
$x\in X$，诱导映射
${\mathcal O}_{Y, f(x)} \to {\mathcal O}_{X, x}$ 都是满射，
但 $f$ 不是闭浸入。
\end{exercise}

\begin{exercise}
\label{exercises-exercise-fibre-product-affines-not-affine}
给出纤维积 $X \times_S Y$ 的例子，使 $X$ 与 $Y$ 都仿射，
但 $X \times_S Y$ 不仿射。
\end{exercise}

\begin{remark}
\label{exercises-remark-separated-base-fibre-product-affines-affine}
事实证明，当 $S$ 分离时，这不可能发生。你知道为什么吗？
\end{remark}

\begin{exercise}
\label{exercises-exercise-not-geometrically-integral}
% P12-ERR-0451
给出概形 $V$ 的例子，使它是 ${\mathbf Q}$ 上有限型的一维整概形，
但 $\Spec({\mathbf C}) \times_{\Spec({\mathbf Q})} V$ 不整。
\end{exercise}

\begin{exercise}
\label{exercises-exercise-not-geometrically-reduced}
% P12-ERR-0452
给出概形 $V$ 的例子，使它是域 $k$ 上有限型的一维整概形，
但对某个有限域扩张 $k'/k$，
$\Spec(k') \times_{\Spec(k)} V$ 不是既约的。
\end{exercise}

\begin{remark}
\label{exercises-remark-affine-dimension}
% P12-ERR-0453
若你的概形是仿射的，则其维数等于底层环的 Krull 维数。
因此可以利用上学期的结果来计算维数。
\end{remark}






\section{态射}
\label{exercises-section-morphisms}

\noindent
一个重要问题是：给定态射 $\pi : X \to S$，
该态射是否有截面或有理截面。下面是一些例题。

\begin{exercise}
\label{exercises-exercise-no-section}
考虑概形态射
$$
\pi :
X = \Spec(\mathbf{C}[x, t, 1/xt])
\longrightarrow
S = \Spec(\mathbf{C}[t]).
$$
\begin{enumerate}
\item 证明不存在态射 $\sigma : S \to X$
使得 $\pi \circ \sigma = \text{id}_S$。
\item 证明存在非空开集 $U \subset S$ 以及态射
$\sigma : U \to X$，使得 $\pi \circ \sigma = \text{id}_U$。
\end{enumerate}
\end{exercise}

\begin{exercise}
\label{exercises-exercise-no-rational-section}
考虑概形态射
$$
\pi :
X = \Spec(\mathbf{C}[x, t]/(x^2 + t))
\longrightarrow
S = \Spec(\mathbf{C}[t]).
$$
证明不存在非空开集 $U \subset S$ 以及态射
$\sigma : U \to X$，使得 $\pi \circ \sigma = \text{id}_U$。
\end{exercise}

\begin{exercise}
\label{exercises-exercise-has-rational-section}
设 $A, B, C \in \mathbf{C}[t]$ 为非零多项式。考虑概形态射
$$
\pi :
X = \Spec(\mathbf{C}[x, y, t]/(A + Bx^2 + Cy^2))
\longrightarrow
S = \Spec(\mathbf{C}[t]).
$$
证明存在非空开集 $U \subset S$ 以及态射
$\sigma : U \to X$，使得 $\pi \circ \sigma = \text{id}_U$。
（提示：在形式上，对某个 $d$，把
$x = X/Z$、$y = Y/Z$ 写成次数 $\leq d$ 的某些
$X, Y, Z \in \mathbf{C}[t]$，并求出它们解此方程的条件。
% P12-ERR-0454
然后用维数理论证明：若 $d >> 0$，便可找到解此方程的非零
$X, Y, Z$。）
\end{exercise}

\begin{remark}
\label{exercises-remark-tsen}
练习 \ref{exercises-exercise-has-rational-section}
是“Tsen 定理”的一个特例。
练习 \ref{exercises-exercise-no-section-curve} 表明，此方法仅限于低次方程
（当底与纤维的维数均为 1 时，即二次曲线）。
\end{remark}

\begin{exercise}
\label{exercises-exercise-no-section-curve}
考虑概形态射
$$
\pi :
X = \Spec(\mathbf{C}[x, y, t]
/(1 + t x^3 + t^2 y^3))
\longrightarrow
S = \Spec(\mathbf{C}[t])
$$
证明不存在非空开集 $U \subset S$ 以及态射
$\sigma : U \to X$，使得 $\pi \circ \sigma = \text{id}_U$。
\end{exercise}

\begin{exercise}
\label{exercises-exercise-no-section-surface}
考虑概形
$$
X = \Spec(\mathbf{C}[\{x_i\}_{i = 1}^{8}, s, t]
/(1 + s x_1^3 + s^2 x_2^3 +
t x_3^3 + st x_4^3 + s^2t x_5^3 +
t^2 x_6^3 + st^2 x_7^3 + s^2t^2 x_8^3))
$$
以及
$$
S = \Spec(\mathbf{C}[s, t])
$$
和概形态射
$$
\pi : X \longrightarrow S
$$
证明不存在非空开集 $U \subset S$ 以及态射
$\sigma : U \to X$，使得 $\pi \circ \sigma = \text{id}_U$。
\end{exercise}

\begin{exercise}
\label{exercises-exercise-for-number-theorists}
（给数论学家。）给出闭子概形
$$
Z \subset \Spec\left({\mathbf Z}[x, \frac{1 }{ x(x-1)(2x-1)}]\right)
$$
使得态射 $Z \to \Spec({\mathbf Z})$ 有限且满。
\end{exercise}

\begin{exercise}
\label{exercises-exercise-quasi-section}
若你不喜欢数论，可以尝试如下变体。考察
$$
\Spec\left({\mathbf F}_p[t, x, \frac{1 }{ x(x-t)(tx-1)}]\right)
\longrightarrow
\Spec({\mathbf F}_p[t])
$$
并试着找出上方概形的一个闭子概形，使其有限满射到下方概形。
% P12-ERR-0455
（这里取有限基域有理论上的原因；不过在这个特例中可能并非必要。）
\end{exercise}

\begin{remark}
\label{exercises-remark-interpretation-skolem-noether}
在练习 \ref{exercises-exercise-for-number-theorists} 与
\ref{exercises-exercise-quasi-section} 中，给定态射 $f : X \to S$，
其中 $S$ 是 Dedekind 环 $A$ 的谱，而 $f$ 平坦且满，并具有
几何不可约的泛纤维。在两种情形中，结论都是：存在有限满态射
$S' \to S$，使得基变换 $X_{S'} \to S'$ 确有截面。
事实证明，若 $A$ 优良，其极大理想处的剩余域是有限域的代数扩张，
且当 $A'$ 是 $A$ 在其分式域的有限扩张中的整闭包时
$\Pic(A')$ 为挠群，则上述结论成立；见 \cite{MB1}。
例如，$A = \mathbf{Z}$ 或 $A = \mathbf{F}_p[x]$，或它们的有限扩张，
都满足这一条件。然而，存在一个 Dedekind 环 $A$，其极大素理想处的
剩余域有限，而其 Picard 群含有非挠元；见 \cite{Goldman}。
对这种环的谱，上述结论为假。
练习 \ref{exercises-exercise-no-quasi-section} 给出了该结论不成立的几何例子。
\end{remark}

\begin{exercise}
\label{exercises-exercise-no-quasi-section}
% P12-ERR-0456
证明存在 $f \in \mathbf{C}[x, t]$，它不被任意
$\alpha \in \mathbf{C}$ 所对应的 $t - \alpha$ 整除，并且不存在
任何 $Z \subset \Spec(\mathbf{C}[x, t, 1/f])$
有限满射到 $\Spec(\mathbf{C}[t])$。
（我认为 $f(x, t) = (xt - 2)(x - t + 3)$ 可行。
不过我觉得，要证明任一候选式具有所需性质并不容易。）
\end{exercise}

\begin{exercise}
\label{exercises-exercise-finite}
设 $A \to B$ 为有限型环映射。假设
$\Spec(B) \to \Spec(A)$ 可分解为某个 $n$ 下的闭浸入
$\Spec(B) \to \mathbf{P}^n_A$。
证明 $A \to B$ 是有限环映射，即 $B$ 作为 $A$-模是有限的。
提示：若 $A$ Noether（请直接假设这一点），则可利用如下事实论证：
对每个闭子概形 $Z \subset \mathbf{P}^n_A$ 和每个
$i \in \mathbf{Z}$，$H^i(Z, \mathcal{O}_Z)$ 都是有限 $A$-模。
\end{exercise}

\begin{exercise}
\label{exercises-exercise-projective-finite}
设 $k$ 为代数闭域。设 $f : X \to Y$ 为射影簇的态射，
并且对每个闭点 $y \in Y$，$f^{-1}(\{y\})$ 都是有限集。
证明 $f$ 作为概形态射是有限的。
提示：（a）有限性是局部性质；（b）试着归约到练习
\ref{exercises-exercise-finite}；（c）利用闭浸入
$X \to \mathbf{P}^n_k$ 得到 $Y$ 上的闭浸入
$X \to \mathbf{P}^n_Y$。
\end{exercise}






\section{切空间}
\label{exercises-section-tangent-space}

\begin{definition}
\label{exercises-definition-dual-numbers}
对任意环 $R$，用 $R[\epsilon]$ 表示{\it 对偶数}环。
作为 $R$-模，它以 $1$、$\epsilon$ 为基自由。
环结构由规定 $\epsilon^2 = 0$ 得到。
\end{definition}

\begin{exercise}
\label{exercises-exercise-tangent-space-Zariski}
设 $f : X \to S$ 为概形态射。设 $x \in X$ 为一点，并令
$s = f(x)$。考虑实线所成的交换图
$$
\xymatrix{
\Spec(\kappa(x)) \ar[r] \ar[dr] \ar@/^1pc/[rr] &
\Spec(\kappa(x)[\epsilon]) \ar@{.>}[r] \ar[d]&
X \ar[d] \\
&
\Spec(\kappa(s)) \ar[r] &
S
}
$$
其中弯曲箭头是 $\Spec(\kappa(x))$ 到 $X$ 的典范态射。
% P12-ERR-0457
若 $\kappa(x) = \kappa(s)$，证明使该图交换的虚线箭头集合
与线性映射集合
$$
\Hom_{\kappa(x)}(
\frac{\mathfrak m_x}{\mathfrak m_x^2 + \mathfrak m_s\mathcal{O}_{X, x}},
\kappa(x))
$$
一一对应。换言之，请描述这样的双射。
（当 $\kappa(x) \supset \kappa(s)$ 为可分代数扩张时，
这一结论更一般地也成立。）
\end{exercise}

\begin{definition}
\label{exercises-definition-tangent-space}
设 $f : X \to S$ 为概形态射，$x \in X$。
我们把练习 \ref{exercises-exercise-tangent-space-Zariski} 中的虚线箭头集合
称为{\it $X$ 在 $S$ 上的切空间}，并记作 $T_{X/S, x}$。
此空间的一个元素称为 $X/S$ 在 $x$ 处的{\it 切向量}。
\end{definition}

\begin{exercise}
\label{exercises-exercise-simple-push-out}
对任意域 $K$，证明图
$$
\xymatrix{
\Spec(K) \ar[r] \ar[d] & \Spec(K[\epsilon_1]) \ar[d] \\
\Spec(K[\epsilon_2]) \ar[r] &
\Spec(K[\epsilon_1, \epsilon_2]/(\epsilon_1\epsilon_2))
}
$$
是概形范畴中的推出图。
（这里仍有 $\epsilon_i^2 = 0$。）
\end{exercise}

\begin{exercise}
\label{exercises-exercise-tangent-space-vectors-space}
设 $f : X \to S$ 为概形态射，$x \in X$。
利用练习 \ref{exercises-exercise-simple-push-out} 以及适当的态射
$$
\Spec(K[\epsilon])
\longrightarrow
\Spec(K[\epsilon_1, \epsilon_2]/(\epsilon_1\epsilon_2)).
$$
定义切向量的加法。类似地，定义切向量的标量乘法（这更容易）。
证明你的构造使 $T_{X/S, x}$ 成为 $\kappa(x)$-向量空间。
\end{exercise}

\begin{exercise}
\label{exercises-exercise-compute-TS}
设 $k$ 为域。考虑结构态射
$f : X = \mathbf{A}^1_k \to \Spec(k) = S$。
\begin{enumerate}
\item 设 $x \in X$ 为闭点。$T_{X/S, x}$ 的维数是多少？
\item 设 $\eta \in X$ 为泛点。$T_{X/S, \eta}$ 的维数是多少？
\item 现在把 $X$ 看作 $\Spec(\mathbf{Z})$ 上的概形。
$T_{X/\mathbf{Z}, x}$ 与 $T_{X/\mathbf{Z}, \eta}$ 的维数分别是多少？
\end{enumerate}
\end{exercise}

\begin{remark}
\label{exercises-remark-tangent-space-relative}
练习 \ref{exercises-exercise-compute-TS} 说明了为何必须考虑 $X$ 在 $S$
上的切空间，才能得到良好的概念。
\end{remark}

\begin{exercise}
\label{exercises-exercise-compute-TS-field}
考虑概形态射
$$
f : X = \Spec(\mathbf{F}_p(t))
\longrightarrow
\Spec(\mathbf{F}_p(t^p)) = S
$$
计算 $X/S$ 在 $X$ 的唯一一点处的切空间。
这不是很奇怪吗？若取与
$\mathbf{F}_p[t^p] \to \mathbf{F}_p[t]$ 对应的概形态射，
你认为会发生什么？
\end{exercise}

\begin{exercise}
\label{exercises-exercise-compute-TS-cusp}
设 $k$ 为域。对
$X = \Spec(k[x, y]/(x^2 - y^3))$，计算 $X/k$ 在点
$x = (0, 0)$ 处的切空间。
\end{exercise}

\begin{exercise}
\label{exercises-exercise-map-tangent-spaces}
设 $f : X \to Y$ 为 $S$ 上的概形态射。设 $x \in X$，
并令 $y = f(x)$。假设自然映射
$\kappa(y) \to \kappa(x)$ 为双射。利用定义证明 $f$
诱导自然线性映射
$$
\text{d}f : T_{X/S, x} \longrightarrow T_{Y/S, y}.
$$
在 $\kappa(x) = \kappa(s)$ 的情形下，通过练习
\ref{exercises-exercise-tangent-space-Zariski} 将它与局部环上的情形对应起来。
\end{exercise}

\begin{exercise}
\label{exercises-exercise-Jacobian}
设 $k$ 为代数闭域。设
% P12-ERR-0458
\begin{eqnarray*}
f : \mathbf{A}_k^n & \longrightarrow & \mathbf{A}^m_k \\
(x_1, \ldots, x_n) & \longmapsto & (f_1(x_i), \ldots, f_m(x_i))
\end{eqnarray*}
为 $k$ 上的概形态射。它由 $n$ 个变量中的 $m$ 个多项式
$f_1, \ldots, f_m$ 给出。考虑矩阵
% P12-ERR-0459
$$
A = \left( \frac{\partial f_j}{\partial x_i} \right)
$$
设 $x \in \mathbf{A}^n_k$ 为闭点，并令 $y =  f(x)$。
证明切空间上的映射
$T_{\mathbf{A}^n_k/k, x} \to T_{\mathbf{A}^m_k/k, y}$
由矩阵 $A$ 在点 $x$ 处的取值给出。
\end{exercise}













\section{拟凝聚层}
\label{exercises-section-quasi-coherent}

\begin{definition}
\label{exercises-definition-quasi-coherent}
设 $X$ 为概形。若对每个仿射开集
$\Spec(R) = U \subset X$，$\mathcal{O}_X$-模层
$\mathcal{F}$ 的限制 $\mathcal{F}|_U$ 都可写成某个
$R$-模 $M$ 所对应的 $\widetilde M$，则称它
是{\it 拟凝聚的}。
\end{definition}

\noindent
% P12-ERR-0460
只需在 $X$ 的一个仿射开覆盖的成员上检验此条件即可。
更多结果见概形篇第 \ref{schemes-section-quasi-coherent} 节。

\begin{definition}
\label{exercises-definition-specialization}
设 $X$ 为拓扑空间，$x, x' \in X$。
当且仅当 $x \in \overline{\{x'\}}$ 时，称 $x$ 是 $x'$ 的
{\it 特化}。
\end{definition}

\begin{exercise}
\label{exercises-exercise-quasi-coherent-specialization-points}
设 $X$ 为概形，$x, x' \in X$，并设 $\mathcal{F}$ 为
$\mathcal{O}_X$-模的拟凝聚层。假设（a）$x$ 是 $x'$ 的特化，
并且（b）$\mathcal{F}_{x'} \not = 0$。
证明 $\mathcal{F}_x \not = 0$。
\end{exercise}

\begin{exercise}
\label{exercises-exercise-O-module-specialization-points}
找出概形 $X$、点 $x, x' \in X$ 以及
$\mathcal{O}_X$-模层 $\mathcal{F}$ 的例子，使得
（a）$x$ 是 $x'$ 的特化，并且（b）
$\mathcal{F}_{x'} \not = 0$ 而 $\mathcal{F}_x = 0$。
\end{exercise}

\begin{definition}
\label{exercises-definition-Noetherian-scheme}
当且仅当对每个点 $x \in X$，都存在仿射开集
$\Spec(R) = U \subset X$，其中 $R$ 为 Noether 环时，
称概形 $X$ {\it 局部 Noether}。若概形局部 Noether 且拟紧，
则称其为 {\it Noether 概形}。
\end{definition}

\noindent
若 $X$ 局部 Noether，则 $X$ 的任意仿射开集都是 Noether 环的谱；
见性质篇引理 \ref{properties-lemma-locally-Noetherian}。

\begin{definition}
\label{exercises-definition-coherent}
设 $X$ 为局部 Noether 概形，并设 $\mathcal{F}$ 为
$\mathcal{O}_X$-模的拟凝聚层。若对每个点 $x \in X$，
都存在仿射开集 $\Spec(R) = U \subset X$，使得
$\mathcal{F}|_U$ 同构于某个有限 $R$-模 $M$ 所对应的
$\widetilde M$，则称 $\mathcal{F}$ {\it 凝聚}。
\end{definition}

\begin{exercise}
\label{exercises-exercise-extend-quasi-coherent}
设 $X = \Spec(R)$ 为仿射概形。
\begin{enumerate}
\item 设 $f \in R$，并设 $\mathcal{G}$ 为开子概形
$D(f)$ 上的拟凝聚 $\mathcal{O}_{D(f)}$-模层。
证明存在拟凝聚 $\mathcal{O}_X$-模层 $\mathcal{F}$，
使 $\mathcal{G} = \mathcal{F}|_{D(f)}$。
\item 设 $I \subset R$ 为理想。令 $i : Z \to X$ 为与 $I$
对应的 $X$ 的闭子概形。设 $\mathcal{G}$ 为闭子概形 $Z$ 上的
拟凝聚 $\mathcal{O}_Z$-模层。证明存在拟凝聚
$\mathcal{O}_X$-模层 $\mathcal{F}$，使得
$\mathcal{G} = i^*\mathcal{F}$。（为什么这很显然？）
\item 假设 $R$ Noether。设 $f \in R$，并设 $\mathcal{G}$
为开子概形 $D(f)$ 上的凝聚 $\mathcal{O}_{D(f)}$-模层。
证明存在凝聚 $\mathcal{O}_X$-模层 $\mathcal{F}$，
使 $\mathcal{G} = \mathcal{F}|_{D(f)}$。
\end{enumerate}
\end{exercise}

\begin{remark}
\label{exercises-remark-extend-off-open}
若 $U \to X$ 是拟紧浸入，则 $U$ 上任意拟凝聚层都是 $X$ 上
某个拟凝聚层的限制。若 $X$ 为 Noether 概形且 $U \subset X$
为开集，则 $U$ 上任意凝聚层都是 $X$ 上某个凝聚层的限制。
当然，上面的练习更容易，不应使用这些一般事实。
\end{remark}

















\section{Proj 与射影概形}
\label{exercises-section-proj}

\begin{exercise}
\label{exercises-exercise-graded-ring-specified-result}
给出分次环 $S$ 的例子，使得
\begin{enumerate}
\item $\text{Proj}(S)$ 仿射且非空；
\item $\text{Proj}(S)$ 整且非空，但对任意 $n\geq 0$、
任意环 $A$，均不同构于 ${\mathbf P}^n_A$。
\end{enumerate}
\end{exercise}

\begin{exercise}
\label{exercises-exercise-nonconstant-morphism-proj}
给出 $\Spec({\mathbf C})$ 上概形的非常值态射
${\mathbf P}^1_{\mathbf C} \to {\mathbf P}^5_{\mathbf C}$ 的例子。
\end{exercise}

\begin{exercise}
\label{exercises-exercise-isomorphism-P1-conic}
给出概形同构的例子
$$
{\mathbf P}^1_{\mathbf C} \to
\text{Proj}({\mathbf C}[X_0, X_1, X_2]/(X_0^2 + X_1^2 + X_2^2))
$$
\end{exercise}

\begin{exercise}
\label{exercises-exercise-family-special-fibre-different}
给出概形态射
$f : X \to {\mathbf A}^1_{\mathbf C} = \Spec({\mathbf C}[T])$
的例子，使 $f$ 在 $t \in {\mathbf A}^1_{\mathbf C}$ 上的
（概形论）纤维 $X_t$ 满足：
（a）当 $t$ 是不等于 $0$ 的闭点时，它同构于
${\mathbf P}^1_{\mathbf C}$；
（b）当 $t = 0$ 时，它不同构于 ${\mathbf P}^1_{\mathbf C}$。
我们把 $X_0$ 称为该态射的{\it 特殊纤维}。
这可以用许多许多种方式做到。试着给出还满足下列每一项附加限制的例子
（除非不可能）：
\begin{enumerate}
\item 能否使特殊纤维射影？
\item 能否使特殊纤维不可约且射影？
\item 能否使特殊纤维整且射影？
\item 能否使特殊纤维光滑且射影？
\item 能否使 $f$ 为平坦态射？
这仅表示：对每个仿射开集 $\Spec(A) \subset X$，诱导环映射
$\mathbf{C}[t] \to A$ 都平坦；在本例中，这表示 $t$ 的任意
非零多项式在 $A$ 上都是非零因子。
\item 能否使 $f$ 为平坦射影态射？
\item 能否使 $f$ 平坦、射影且特殊纤维既约？
\item 能否使 $f$ 平坦、射影且特殊纤维不可约？
\item 能否使 $f$ 平坦、射影且特殊纤维整？
\end{enumerate}
若把 ${\mathbf P}^1_{\mathbf C}$ 换成 ${\mathbf C}$ 上的另一个簇，
你认为会发生什么？
（取决于你要求上述哪个变体，这可能会变得非常困难。）
\end{exercise}

\begin{exercise}
\label{exercises-exercise-affine-onto-projective-space}
设 $n \geq 1$ 为任意正整数。给出以仿射概形 $X$ 为源的满态射
$X \to {\mathbf P}^n_{\mathbf C}$ 的例子。
\end{exercise}

\begin{exercise}
\label{exercises-exercise-morphism-proj}
$\text{Proj}$ 之间的映射。设 $R$ 与 $S$ 为分次环。
假设有环映射
$$
\psi : R \to S
$$
以及整数 $e \geq 1$，使得对所有 $d \geq 0$，
$\psi(R_d) \subset S_{de}$。
（按我们的约定，除非 $e = 1$，这不是分次环同态。）
\begin{enumerate}
\item 对哪些元素 $\mathfrak p \in \text{Proj}(S)$，
$\text{Proj}(R)$ 中有定义良好的对应点？换言之，找出适当的开集
$U \subset \text{Proj}(S)$，使 $\psi$ 定义连续映射
$r_\psi : U \to \text{Proj}(R)$。
\item 给出 $U \not = \text{Proj}(S)$ 的例子。
\item 给出 $U = \text{Proj}(S)$ 的例子。
\item （不要把它写下来。）说服自己：连续映射
$U \to \text{Proj}(R)$ 典范地带有层上的映射，从而 $r_\psi$
是概形态射：
$$
\text{Proj}(S) \supset U \longrightarrow \text{Proj}(R).
$$
\item 若
$R = \bigoplus_{d \geq 0} S_{de}$（作为分次环，把
$S_e$、$S_{2e}$ 等分别放在次数 $1$、$2$ 等处），且 $\psi$
是包含映射，你能对该映射说些什么？
\end{enumerate}
\end{exercise}

\noindent
{\bf 记号。}设 $R$ 为如上的分次环，$n \geq 0$ 为整数，
并令 $X = \text{Proj}(R)$。则 $X$ 上存在唯一的拟凝聚
${\mathcal O}_X$-模 ${\mathcal O}_X(n)$，使得对每个正次数齐次元
$f \in R$，${\mathcal O}_X |_{D_{+}(f)}$ 是与
$R_{(f)} = (R_f)_0$-模 $(R_f)_n$ 相伴的拟凝聚层
% P12-ERR-0461
（$ = $在 $R_f = R[1/f]$ 中次数为 $n$ 的齐次元）。
见 \cite[page 116+]{H}。注意，存在自然映射
$$
{\mathcal O}_X(n_1) \otimes_{{\mathcal O}_X} {\mathcal O}_X(n_2)
\longrightarrow
{\mathcal O}_X(n_1 + n_2)
$$

\begin{exercise}
\label{exercises-exercise-pathologies-proj}
$\text{Proj}$ 中的病态现象。
给出如上的 $R$ 的例子，使得
\begin{enumerate}
\item ${\mathcal O}_X(1)$ 不是可逆 ${\mathcal O}_X$-模。
\item ${\mathcal O}_X(1)$ 可逆，但自然映射
${\mathcal O}_X(1) \otimes_{{\mathcal O}_X} {\mathcal O}_X(1) \to
{\mathcal O}_X(2)$ 不是同构。
\end{enumerate}
\end{exercise}

\begin{exercise}
\label{exercises-exercise-finitely-many-points-in-affine}
设 $S$ 为分次环，并令 $X = \text{Proj}(S)$。
证明 $X$ 的任意有限点集都包含在某个标准仿射开集中。
\end{exercise}

\begin{exercise}
\label{exercises-exercise-prepare-glueing}
设 $S$ 为分次环，$X = \text{Proj}(S)$。
设 $Z, Z' \subset X$ 为两个闭子概形，并设
$\varphi : Z \to Z'$ 为同构。假设 $Z \cap Z' = \emptyset$。
证明：对任意 $z \in Z$，都存在仿射开集 $U \subset X$，使得
$z \in U$、$\varphi(z) \in U$ 且
$\varphi(Z \cap U) = Z' \cap U$。
（提示：使用练习 \ref{exercises-exercise-finitely-many-points-in-affine}
以及类似于概形篇引理
\ref{schemes-lemma-standard-open-two-affines} 的方法。）
\end{exercise}







\section{从射影直线出发的态射}
\label{exercises-section-from-P1}

\noindent
本节研究从 $\mathbf{P}^1$ 到射影概形的态射。

\begin{exercise}
\label{exercises-exercise-from-generic-point}
设 $k$ 为域，并设 $k[t] \subset k(t)$ 为多项式环到其分式域的包含。
设 $X$ 为 $k$ 上有限型概形。证明：对任意 $k$ 上的态射
$$
\varphi : \Spec(k(t)) \longrightarrow X
$$
都存在非零 $f \in k[t]$ 以及 $k$ 上的态射
$\psi : \Spec(k[t, 1/f]) \to X$，使 $\varphi$ 为复合
$$
\Spec(k(t)) \longrightarrow \Spec(k[t, 1/f]) \longrightarrow X
$$
\end{exercise}

\begin{exercise}
\label{exercises-exercise-from-generic-point-Pn}
设 $k$ 为域，并设 $k[t] \subset k(t)$ 为多项式环到其分式域的包含。
证明：对任意 $k$ 上的态射
$$
\varphi : \Spec(k(t)) \longrightarrow \mathbf{P}^n_k
$$
都存在 $k$ 上的态射
$\psi : \Spec(k[t]) \to \mathbf{P}^n_k$，使 $\varphi$ 为复合
$$
\Spec(k(t)) \longrightarrow \Spec(k[t]) \longrightarrow \mathbf{P}^n_k
$$
提示：$\varphi$ 的像包含在某个 $i$ 所对应的标准开集
$D_+(T_i)$ 中；然后证明可以“清除分母”。
\end{exercise}

\begin{exercise}
\label{exercises-exercise-from-generic-point-projective}
设 $k$ 为域，并设 $k[t] \subset k(t)$ 为多项式环到其分式域的包含。
设 $X$ 为 $k$ 上的射影概形。证明：对任意 $k$ 上的态射
$$
\varphi : \Spec(k(t)) \longrightarrow X
$$
都存在 $k$ 上的态射 $\psi : \Spec(k[t]) \to X$，
使 $\varphi$ 为复合
$$
\Spec(k(t)) \longrightarrow \Spec(k[t]) \longrightarrow X
$$
提示：使用练习 \ref{exercises-exercise-from-generic-point-Pn}。
\end{exercise}

\begin{exercise}
\label{exercises-exercise-from-generic-point-P1-to-projective}
设 $k$ 为域，$X$ 为 $k$ 上的射影概形。
令 $K$ 为 $\mathbf{P}^1_k$ 的函数域（见下方提示）。
证明：对任意 $k$ 上的态射
$$
\varphi : \Spec(K) \longrightarrow X
$$
都存在 $k$ 上的态射 $\psi : \mathbf{P}^1_k \to X$，
使 $\varphi$ 为复合
$$
\Spec(k(t)) \longrightarrow \mathbf{P}^1_k \longrightarrow X
$$
提示：对仿射开覆盖
$\mathbf{P}^1_k = D_+(T_0) \cup D_+(T_1)$ 的两个部分分别使用
练习 \ref{exercises-exercise-from-generic-point-projective}；利用
$D_+(T_0)$ 是多项式环的谱，并且 $K$ 是该多项式环的分式域。
\end{exercise}







\section{从曲面到曲线的态射}
\label{exercises-section-from-surfaces-to-curves}

\begin{exercise}
\label{exercises-exercise-points-projective-space}
设 $R$ 为环。设 $R \to k$ 为从 $R$ 到某个域的映射。
设 $n \geq 0$。证明
$$
\Mor_{\Spec(R)}(\Spec(k), \mathbf{P}^n_R)
=
(k^{n + 1} \setminus \{0\})/k^*
$$
其中 $k^*$ 通过标量乘法作用在 $k^{n + 1}$ 上。
此后，用 $(x_0 : \ldots : x_n)$ 表示与元素
$(x_0, \ldots, x_n) \in k^{n + 1} \setminus \{0\}$ 的等价类
对应的态射 $\Spec(k) \to \mathbf{P}^n_k$。
\end{exercise}

\begin{exercise}
\label{exercises-exercise-curve-projective-plane}
设 $k$ 为域，并设 $Z \subset \mathbf{P}^2_k$ 为不可约既约闭子概形。
证明下列三种情形之一成立：（a）$Z$ 是闭点；或
% P12-ERR-0462
（b）存在次数 $> 0$ 的齐次不可约多项式
$F \in k[X_0, X_1, X_2]$，使 $Z = V_{+}(F)$；或
（c）$Z = \mathbf{P}^2_k$。（提示：考察一个标准仿射开集。）
\end{exercise}

\begin{exercise}
\label{exercises-exercise-bezout}
% P12-ERR-0463
设 $k$ 为域。设 $Z_1, Z_2 \subset \mathbf{P}^2_k$
为形如 $V_{+}(F)$ 的不可约闭子概形，其中
$F_i \in k[X_0, X_1, X_2]$ 是次数 $> 0$ 的齐次不可约多项式。
证明 $Z_1 \cap Z_2$ 非空。
（提示：用维数理论估计
$k[X_0, X_1, X_2]/(F_1, F_2)$ 在 $0$ 处局部环的维数。）
\end{exercise}

\begin{exercise}
\label{exercises-exercise-no-nonconstant-morphism-proj}
证明不存在 $\Spec(\mathbf{C})$ 上概形的非常值态射
$\mathbf{P}^2_{\mathbf{C}} \to \mathbf{P}^1_{\mathbf{C}}$。
这里，{\it 常值态射}是指其像仅为一个点的态射。
（提示：若态射不是常值的，考察 $0$ 与 $\infty$ 上的纤维，
并论证它们必然相交，从而得到矛盾。）
\end{exercise}

\begin{exercise}
\label{exercises-exercise-nonconstant-morphism}
设 $k$ 为域。假设 $X \subset \mathbf{P}^3_k$ 是由单个齐次方程
$F \in k[X_0, X_1, X_2, X_3]$ 给出的闭子概形。换言之，
$$
X = \text{Proj}(k[X_0, X_1, X_2, X_3]/(F)) \subset \mathbf{P}^3_k
$$
如课程中所述。假设
$$
F = X_0 G + X_1 H
$$
其中 $G, H \in k[X_0, X_1, X_2, X_3]$ 是正次数齐次多项式。
证明：若 $X_0, X_1, G, H$ 没有公共零点，则存在
$\Spec(k)$ 上概形的非常值态射
$$
X \longrightarrow \mathbf{P}^1_k
$$
它在域值点上（见练习 \ref{exercises-exercise-points-projective-space}），
只要 $x_0$ 或 $x_1$ 非零，就表现为
$(x_0 : x_1 : x_2 : x_3) \mapsto (x_0 : x_1)$。
\end{exercise}





\section{可逆层}
\label{exercises-section-invertible-sheaves}

\begin{definition}
\label{exercises-definition-invertible-sheaf}
设 $X$ 为局部环空间。$X$ 上的一个
{\it 可逆 ${\mathcal O}_X$-模}，是 ${\mathcal O}_X$-模层
${\mathcal L}$，使得每一点都有开邻域 $U \subset X$，
且 ${\mathcal L}|_U$ 作为 ${\mathcal O}_U$-模同构于
${\mathcal O}_U$。若 ${\mathcal L}$ 作为
${\mathcal O}_X$-模同构于 ${\mathcal O}_X$，则称它平凡。
\end{definition}

\begin{exercise}
\label{exercises-exercise-general-facts-invertible}
一般事实。
\begin{enumerate}
\item 证明概形 $X$ 上的可逆 ${\mathcal O}_X$-模是拟凝聚的。
\item 假设 $X\to Y$ 是局部环空间的态射，并且
% P12-ERR-0464
${\mathcal L}$ 是可逆 ${\mathcal O}_Y$-模。
% P12-ERR-0465
证明 $f^\ast {\mathcal L}$ 是可逆 ${\mathcal O}_X$ 模。
\end{enumerate}
\end{exercise}

\begin{exercise}
\label{exercises-exercise-invertible-algebra}
代数。
\begin{enumerate}
\item 证明仿射概形 $\Spec(A)$ 上的可逆
${\mathcal O}_X$-模对应于 $A$-模 $M$，后者
（i）有限，（ii）射影，（iii）局部自由且秩为 1，
从而（iv）平坦，并且（v）有限表现。
% P12-ERR-0466
（可以随意引用上学期课程中的结论，或代数教材中的结论。）
\item 假设 $A$ 为整环，$M$ 是（a）中那样的模。
% P12-ERR-0467
证明 $M$ 作为 $A$-模同构于某个理想 $I \subset A$，
且对每个素理想 ${\mathfrak p}$，$IA_{\mathfrak p}$ 都是主理想。
\end{enumerate}
\end{exercise}

\begin{definition}
\label{exercises-definition-invertible-module}
设 $R$ 为环。所谓{\it 可逆模 $M$}，是指 $R$-模 $M$
满足：$\widetilde M$ 是 $R$ 的谱上的可逆层。
若 $M$ 作为 $R$-模满足 $M \cong R$，则称该模{\it 平凡}。
\end{definition}

\noindent
换言之，当且仅当 $M$ 满足以下全部条件时，它才可逆：
它平坦、有限表现、射影，并且局部自由且秩为 1。
（当然，仅要求它局部自由且秩为 1 就足够了。）

\begin{exercise}
\label{exercises-exercise-simple-examples-invertible}
简单例子。
\begin{enumerate}
\item
\label{exercises-item-affine-line}
设 $k$ 为域，$A = k[x]$。
证明 $X = \Spec(A)$ 仅有平凡的可逆
${\mathcal O}_X$-模。换言之，证明每个可逆 $A$-模都自由且秩为 1。
\item
\label{exercises-item-affine-line-with-0-and-1-identified}
设 $A$ 为环
$$
A = \{ f\in k[x] \mid f(0) = f(1) \}.
$$
证明除非 $k = {\mathbf F}_2$，否则存在非平凡可逆 $A$-模。
（提示：把 $\Spec(A)$ 视为在
${\mathbf A}^1_k = \Spec(k[x])$ 中将 $0$ 与 $1$ 粘合。）
\item
\label{exercises-item-affine-line-with-cusp}
对环 $A = k[x^2, x^3] \subset k[x]$ 回答
（\ref{exercises-item-affine-line-with-0-and-1-identified}）中的同一问题
（但现在 $k = {\mathbf F}_2$ 也可行）。
\end{enumerate}
\end{exercise}

\begin{exercise}
\label{exercises-exercise-higher-dimension-invertible}
更高维情形。
\begin{enumerate}
\item 证明二维仿射空间上的每个可逆层都平凡。更精确地，设
${\mathbf A}^2_k = \Spec(k[x, y])$，其中 $k$ 为域。
证明 ${\mathbf A}^2_k$ 上的每个可逆层都平凡。
（提示：一种做法是考察对应模 $M$，再考察
$M \otimes_{k[x, y]} k(x)[y]$，然后用练习
\ref{exercises-exercise-simple-examples-invertible}
（\ref{exercises-item-affine-line}）找出它的一个生成元；之后你仍需思考。
% P12-ERR-0468
另一种做法是使用练习 \ref{exercises-exercise-invertible-algebra}，
并利用我们关于多项式环中理想的知识：高度一的素理想由一个
不可约多项式生成；之后你仍需思考。）
\item 证明二维仿射空间的任意开子概形上的每个可逆层都平凡。
更精确地，设 $U \subset {\mathbf A}^2_k$ 为开子概形，其中 $k$
为域。证明 $U$ 上的每个可逆层都平凡。提示：证明 $U$ 上每个
可逆层都可延拓为 ${\mathbf A}^2_k$ 上的可逆层。这并不容易，
但可在 \cite{H} 中找到。
\item 找出域上去掉顶点的锥面上的一个非平凡可逆层。
更精确地，设 $k$ 为域，$C = \Spec(k[x, y, z]/(xy-z^2))$。
令 $U = C \setminus \{ (x, y, z) \}$。找出 $U$ 上的一个
非平凡可逆层。提示：计算 $U$ 上可逆层的同构类群，可能比
仅找出一个更容易。注意，$U$ 被开集
$\Spec(k[x, y, z, 1/x]/(xy-z^2))$ 与
$\Spec(k[x, y, z, 1/y]/(xy-z^2))$ 覆盖，
而它们是“容易”处理的。
\end{enumerate}
\end{exercise}

\begin{definition}
\label{exercises-definition-picard-group}
设 $X$ 为局部环空间。$X$ 的{\it Picard 群}是可逆
$\mathcal{O}_X$-模的同构类集合 $\Pic(X)$，其加法由张量积给出。
见模篇定义 \ref{modules-definition-pic}。
对环 $R$，令 $\Pic(R) = \Pic(\Spec(R))$。
\end{definition}

\begin{exercise}
\label{exercises-exercise-traverso}
设 $R$ 为环。
\begin{enumerate}
\item 证明：若 $R$ 是 Noether 正规整环，则
$\Pic(R) = \Pic(R[t])$。[提示：存在映射 $R[t] \to R$，
$t \mapsto 0$，它是映射 $R \to R[t]$ 的左逆。
因此只需证明：若可逆 $R[t]$-模 $M$ 满足
$M/tM \cong R$，则它自由且秩为 $1$。令 $K$ 为 $R$ 的分式域。
选取平凡化 $K[t] \to M \otimes_{R[t]} K[t]$；由练习
\ref{exercises-exercise-simple-examples-invertible}（\ref{exercises-item-affine-line}），
这是可能的。调整它，使其与上面 $M/tM$ 的平凡化相容。
证明它事实上是 $M$ 在 $R[t]$ 上的一个平凡化
（正规性在此处起作用）。]
\item 设 $k$ 为域。证明
$\Pic(k[x^2, x^3, t]) \not = \Pic(k[x^2, x^3])$。
\end{enumerate}
\end{exercise}














\section{{\v C}ech 上同调}
\label{exercises-section-cech-cohomology}

\begin{exercise}
\label{exercises-exercise-cech-cohomology}
{\v C}ech 上同调。这里 $k$ 为域。
\begin{enumerate}
\item 设 $X$ 为带有开覆盖
${\mathcal U} : X = U_1 \cup U_2$ 的概形，其中
$U_1 = \Spec(k[x])$、$U_2 =  \Spec(k[y])$，
$U_1 \cap U_2 = \Spec(k[z, 1/z])$，而开浸入
$U_1 \cap U_2 \to U_1$ 与 $U_1 \cap U_2 \to U_2$
分别由 $x \mapsto z$ 与 $y \mapsto z$ 决定
（我确实就是这个意思）。
（我们已在讲课中见过这样的 $X$ 存在；它是原点加倍的仿射直线。）
计算 ${\check H}^1({\mathcal U}, {\mathcal O})$；
例如，给出它作为 $k$-向量空间的一组基。
\item 对
${\check H}^1({\mathcal U}, {\mathcal O})$ 中的每个元素，
构造 ${\mathcal O}_X$-模层的正合列
$$
0 \to {\mathcal O}_X \to E \to {\mathcal O}_X \to 0
$$
使边界 $\delta(1) \in {\check H}^1({\mathcal U},
{\mathcal O})$ 等于给定元素。
（问题的一部分是赋予这句话意义；另见下文。
抽象地证明这样的对象必然存在也可以。）
\end{enumerate}
\end{exercise}

\begin{definition}
\label{exercises-definition-delta}
（delta 的定义。）假设
$$
0 \to {\mathcal F}_1 \to {\mathcal F}_2 \to {\mathcal F}_3 \to 0
$$
是任意拓扑空间 $X$ 上 Abel 层的短正合列。边界映射
$\delta : H^0(X, {\mathcal F}_3) \to {\check H}^1(X, {\mathcal F}_1)$
定义如下。取元素 $\tau \in H^0(X, {\mathcal F}_3)$。
选取开覆盖 ${\mathcal U} : X = \bigcup_{i\in I} U_i$，使对每个
$i$，都存在截面 $\tilde \tau_i \in {\mathcal F}_2$，
它提升 $\tau$ 在 $U_i$ 上的限制。然后考虑赋值
$$
(i_0, i_1) \longmapsto
\tilde \tau_{i_0}|_{U_{i_0i_1}} - \tilde \tau_{i_1}|_{U_{i_0i_1}}.
$$
这是 {\v C}ech 复形
${\check C}^\ast({\mathcal U}, {\mathcal F}_2)$ 中显然的 1-余边界。
但注意到（把 ${\mathcal F}_1$ 看作 ${\mathcal F}_2$ 的子层），
右端总是 $U_{i_0i_1}$ 上 ${\mathcal F}_1$ 的截面。
% P12-ERR-0469
因此可见，该赋值在复形
${\check C}^\ast({\mathcal U}, {\mathcal F}_2)$ 中定义了一个
1-上链。按照定义，这个 1-上链的上同调类就是
{\it $\delta(\tau)$}。
\end{definition}




\section{上同调}
\label{exercises-section-cohomology}

\begin{exercise}
\label{exercises-exercise-cohomology-not-zero}
令 $X = \mathbf{R}$，赋予通常的 Euclid 拓扑。
仅使用上同调的形式性质（函子性与长正合上同调列），证明存在
$X$ 上的层 $\mathcal{F}$，使 $H^1(X, \mathcal{F})$ 非零。
\end{exercise}

\begin{exercise}
\label{exercises-exercise-mayer-vietoris}
设拓扑空间 $X = U \cup V$ 写成两个开集的并。
则有 Mayer--Vietoris 长正合列
$$
0 \to
H^0(X, \mathcal{F}) \to
H^0(U, \mathcal{F}) \oplus H^0(V, \mathcal{F}) \to
H^0(U \cap V, \mathcal{F}) \to
H^1(X, \mathcal{F}) \to \ldots
$$
内射层的什么性质对于构造 Mayer--Vietoris 长正合列至关重要？
为什么该性质成立？
\end{exercise}

\begin{exercise}
\label{exercises-exercise-cohomology-two-point-space}
设 $X$ 为拓扑空间。
\begin{enumerate}
\item 若 $X$ 至多有 $2$ 个点，证明对 $i > 0$，
$H^i(X, \mathcal{F})$ 为零。
\item 若 $X$ 有 $3$ 个点，又如何？
\end{enumerate}
\end{exercise}

\begin{exercise}
\label{exercises-exercise-cohomology-spec-local-ring}
设 $X$ 为局部环的谱。证明：对任意 Abel 群层
$\mathcal{F}$ 和任意 $i > 0$，$H^i(X, \mathcal{F})$ 为零。
\end{exercise}

\begin{exercise}
\label{exercises-exercise-affine-morphism}
设 $f : X \to Y$ 为概形的仿射态射。证明：对任意拟凝聚
$\mathcal{O}_X$-模 $\mathcal{F}$，
$H^i(X, \mathcal{F}) = H^i(Y, f_*\mathcal{F})$。
你可以随意对 $X$ 与 $Y$ 加上一些条件，并使用如下相合性：
对分离概形的仿射开覆盖与拟凝聚层，{\v C}ech 上同调与上同调相合。
\end{exercise}

\begin{exercise}
\label{exercises-exercise-affine-morphism-to-Pn}
设 $A$ 为环，并设
$\mathbf{P}^n_A = \text{Proj}(A[T_0, \ldots, T_n])$
为 $A$ 上的射影空间。

令
$\mathbf{A}^{n + 1}_A = \Spec(A[T_0, \ldots, T_n])$，并令
$$
U = \bigcup\nolimits_{i = 0, \ldots, n} D(T_i) \subset \mathbf{A}^{n + 1}_A
$$
为闭浸入 $0 : \Spec(A) \to \mathbf{A}^{n + 1}_A$ 的像的补集。
构造仿射满态射
$$
f : U \longrightarrow \mathbf{P}^n_A
$$
并证明
$f_*\mathcal{O}_U =
\bigoplus_{d \in \mathbf{Z}} \mathcal{O}_{\mathbf{P}^n_A}(d)$。
更一般地，证明对分次 $A[T_0, \ldots, T_n]$-模 $M$ 有
$$
f_*(\widetilde{M}|_U) =
\bigoplus\nolimits_{d \in \mathbf{Z}} \widetilde{M(d)}
$$
其中左端的 $\widetilde{M}$ 是 $\mathbf{A}^{n + 1}_A$ 上
与 $M$ 相伴的拟凝聚层，而右端的 $\widetilde{M(d)}$
是与分次模 $M(d)$ 相伴的拟凝聚层。
\end{exercise}

\begin{exercise}
\label{exercises-exercise-compute-Pn}
设 $A$ 为环，并设
$\mathbf{P}^n_A = \text{Proj}(A[T_0, \ldots, T_n])$
为 $A$ 上的射影空间。仔细计算 $\mathbf{P}^n_A$ 上结构层的
Serre 扭曲 $\mathcal{O}_{\mathbf{P}^n_A}(d)$ 的上同调。
你可以使用 {\v C}ech 上同调，以及分离概形的仿射开覆盖与
拟凝聚层所具有的 {\v C}ech 上同调和上同调相合性。
\end{exercise}

\begin{exercise}
\label{exercises-exercise-cohomology-hypersurface}
设 $A$ 为环，并设
$\mathbf{P}^n_A = \text{Proj}(A[T_0, \ldots, T_n])$
为 $A$ 上的射影空间。设 $F \in A[T_0, \ldots, T_n]$
为次数 $d$ 的齐次多项式。令 $X \subset \mathbf{P}^n_A$
为与 $A[T_0, \ldots, T_n]$ 的分次理想 $(F)$ 对应的闭子概形。
关于 $H^i(X, \mathcal{O}_X)$，你能说些什么？
\end{exercise}

\begin{exercise}
\label{exercises-exercise-characterize-finite-product-fields}
设环 $R$ 满足：对任意左正合函子
$F : \text{Mod}_R \to \textit{Ab}$，对 $i > 0$ 都有
$R^iF = 0$。证明 $R$ 是有限多个域的乘积。
\end{exercise}







\section{更多上同调}
\label{exercises-section-more-cohomology}

\begin{exercise}
\label{exercises-exercise-cohomology-coordinate-axes}
设 $k$ 为域。设 $X \subset \mathbf{P}^n_k$ 为“坐标十字”。
具体地，令 $X$ 由齐次方程
$$
T_i T_j = 0\text{ 对 }i > j > 0
$$
定义；照例写
$\mathbf{P}^n_k = \text{Proj}(k[T_0, \ldots, T_n])$。
换言之，$X$ 是与多项式环的商
$k[T_0, \ldots, T_n]/(T_iT_j; i > j > 0)$ 对应的闭子概形。
对所有 $i$ 计算 $H^i(X, \mathcal{O}_X)$。
提示：使用 {\v C}ech 上同调。
\end{exercise}

\begin{exercise}
\label{exercises-exercise-compute-cohomology-punctured}
设 $A$ 为环，$I = (f_1, \ldots, f_t)$ 为 $A$ 的有限生成理想。
令 $U \subset \Spec(A)$ 为 $V(I)$ 的补集。
对任意 $A$-模 $M$，写出一个 $A$-模复形
（用 $A$、$f_1, \ldots, f_t$、$M$ 表示），
使其上同调群给出 $H^n(U, \widetilde{M})$。
\end{exercise}

\begin{exercise}
\label{exercises-exercise-compute-cohomology-affine-space-punctured}
设 $k$ 为域。令 $U \subset \mathbf{A}^d_k$ 为
$\mathbf{A}^d_k$ 中闭点 $0$ 的补集。
对所有 $n$ 计算 $H^n(U, \mathcal{O}_U)$。
\end{exercise}

\begin{exercise}
\label{exercises-exercise-find-curve-genus-one-hundred}
设 $k$ 为域。显式找出 $k$ 上射影、维数为 $1$ 的概形 $X$，
使 $H^0(X, \mathcal{O}_X) = k$ 且
$\dim_k H^1(X, \mathcal{O}_X) = 100$。
\end{exercise}

\begin{exercise}
\label{exercises-exercise-degree-2-cover}
设 $f : X \to Y$ 为次数 $2$ 的有限局部自由态射。
假设 $X$ 与 $Y$ 都是整概形，并且 $2$ 在 $Y$ 的结构层中可逆，
即 $2 \in \Gamma(Y, \mathcal{O}_Y)$ 可逆。
证明 $\mathcal{O}_Y$-模映射
$$
f^\sharp : \mathcal{O}_Y \longrightarrow f_*\mathcal{O}_X
$$
有左逆，即存在 $\mathcal{O}_Y$-模映射
$\tau : f_*\mathcal{O}_X \to \mathcal{O}_Y$，使得
$\tau \circ f^\sharp = \text{id}$。
由此推出 $H^n(Y, \mathcal{O}_Y) \to H^n(X, \mathcal{O}_X)$
为单射\footnote{确实存在整概形之间次数为 $2$ 的有限局部自由态射
$X \to Y$，使得映射
$H^1(Y, \mathcal{O}_Y) \to H^1(X, \mathcal{O}_X)$ 不是单射。}。
\end{exercise}

\begin{exercise}
\label{exercises-exercise-pic}
设 $X$ 为概形（或局部环空间）。规则
$U \mapsto \mathcal{O}_X(U)^*$ 定义一个群层，记为
$\mathcal{O}_X^*$。简要说明为何 $X$ 的 Picard 群
（定义 \ref{exercises-definition-picard-group}）等于
$H^1(X, \mathcal{O}_X^*)$。
\end{exercise}

\begin{exercise}
\label{exercises-exercise-pic-nontrivial}
给出具有非平凡 $\Pic(X)$ 的仿射概形 $X$ 的例子。
利用练习 \ref{exercises-exercise-pic} 推出：对任意这样的 $X$，
$H^1(X, -)$ 都不是零函子。
\end{exercise}

\begin{exercise}
\label{exercises-exercise-kill-cohomology-complement}
设 $A$ 为环，$I = (f_1, \ldots, f_t)$ 为 $A$ 的有限生成理想，
并令 $U \subset \Spec(A)$ 为 $V(I)$ 的补集。
给定拟凝聚 $\mathcal{O}_{\Spec(A)}$-模 $\mathcal{F}$
以及满足 $p > 0$ 的 $\xi \in H^p(U, \mathcal{F})$，
证明存在 $n > 0$，使得对 $i = 1, \ldots, t$ 有
$f_i^n \xi = 0$。提示：一种可能的做法是使用你在练习
\ref{exercises-exercise-compute-cohomology-punctured} 中找到的复形。
\end{exercise}

\begin{exercise}
\label{exercises-exercise-h0-complement}
设 $A$ 为环，$I = (f_1, \ldots, f_t)$ 为 $A$ 的有限生成理想，
并令 $U \subset \Spec(A)$ 为 $V(I)$ 的补集。
设 $M$ 为 $A$-模，其 $I$-挠子模为零，即
$0 = \Ker((f_1, \ldots, f_t) : M \to M^{\oplus t})$。
证明存在典范同构
$$
H^0(U, \widetilde{M}) = \colim \Hom_A(I^n, M).
$$
警告：这并非显然。
\end{exercise}

\begin{exercise}
\label{exercises-exercise-Noetherian}
设 $A$ 为 Noether 环，$I$ 为 $A$ 的理想，$M$ 为 $A$-模。
令 $M[I^\infty]$ 为由下式定义的 $I$-幂挠元集合：
$$
M[I^\infty] = \{x \in M \mid
\text{ there exists an }n \geq 1\text{ 使得 }I^nx = 0\}
$$
令 $M' = M/M[I^\infty]$。则 $M'$ 的 $I$-幂挠子模为零。
证明
$$
\colim \Hom_A(I^n, M) = \colim \Hom_A(I^n, M').
$$
警告：这并非显然。提示：（1）试着归约到 $M$ 有限；
（2）若 $N = M[I^\infty]$ 且 $M$ 有限，证明
$\Ext^1_A(I^n, N)$ 的任意元素对某个 $m > 0$ 映到
$\Ext^1_A(I^{n + m}, N)$ 中的零；
（3）对 $\Hom_A(I^n, N)$ 证明与（2）相同的结论；
% P12-ERR-0470
（3）考察长正合列
$$
0 \to \Hom_A(I^n, M[I^\infty]) \to \Hom_A(I^n, M) \to \Hom_A(I^n, M')
\to \Ext^1_A(I^n, M[I^\infty])
$$
其中 $M$ 有限，并与 $I^{n + m}$ 所对应的正合列比较以得出结论。
\end{exercise}








\section{再论上同调}
\label{exercises-section-more-more-cohomology}


\begin{exercise}
\label{exercises-exercise-nonkillable}
构造如下例子：一个域 $k$、$k$ 上的一条曲线 $X$、
一个可逆 $\mathcal{O}_X$-模 $\mathcal{L}$，以及上同调类
$\xi \in H^1(X, \mathcal{L})$，使其具有如下性质：
对概形的每个有限满态射 $\pi : Y \to X$，元素 $\xi$
都拉回为 $H^1(Y, \pi^*\mathcal{L})$ 中的非零元素。
提示：构造 $X$、$k$、$\mathcal{L}$，使得存在短正合列
$0 \to \mathcal{L} \to \mathcal{O}_X \to i_*\mathcal{O}_Z \to 0$，
其中 $Z \subset X$ 是由多于 $1$ 个闭点组成的闭子概形。
% P12-ERR-0471
然后考察拉回时会发生什么。
\end{exercise}

\begin{exercise}
\label{exercises-exercise-cohomology-cycle-curves}
设 $k$ 为代数闭域，$X$ 为 $1$-维射影概形。
假设 $X$ 含有一个曲线圈，即假设存在 $n \geq 2$、
两两不同的 $1$-维整闭子概形 $C_1, \ldots, C_n$ 以及两两不同的闭点
$x_1, \ldots, x_n \in X$，使得 $x_n \in C_n \cap C_1$，
且对 $i = 1, \ldots, n - 1$ 有 $x_i \in C_i \cap C_{i + 1}$。
证明 $H^1(X, \mathcal{O}_X)$ 非零。
提示：令 $\mathcal{F}$ 为映射
% P12-ERR-0472
$\mathcal{O}_X \to \bigoplus \mathcal{O}_{C_i}$ 的像，
并利用 $\kappa(x_i) = k$ 与
$H^0(C_i, \mathcal{O}_{C_i}) = k$，证明
$H^1(X, \mathcal{F})$ 非零。
还要利用 Grothendieck 定理所给出的 $H^2(X, -) = 0$。
\end{exercise}

\begin{exercise}
\label{exercises-exercise-surface}
设 $X$ 为代数闭域 $k$ 上的射影曲面。
证明存在真闭子概形 $Z \subset X$，使 $H^1(Z, \mathcal{O}_Z)$
非零。提示：使用练习 \ref{exercises-exercise-cohomology-cycle-curves}。
\end{exercise}

\begin{exercise}
\label{exercises-exercise-surface-better}
设 $X$ 为代数闭域 $k$ 上的射影曲面。
证明对每个 $n \geq 0$，都存在真闭子概形 $Z \subset X$，
使得 $\dim_k H^1(Z, \mathcal{O}_Z) > n$。
只需说明如何修改练习 \ref{exercises-exercise-surface} 与
\ref{exercises-exercise-cohomology-cycle-curves} 中的论证；不必给出全部细节。
\end{exercise}

\begin{exercise}
\label{exercises-exercise-surface-h2-nonzero}
设 $X$ 为代数闭域 $k$ 上的射影曲面。
证明存在凝聚 $\mathcal{O}_X$-模 $\mathcal{F}$，使得
$H^2(X, \mathcal{F})$ 非零。
提示：使用练习 \ref{exercises-exercise-surface-better} 的结论以及一个巧妙选取的
正合列。
\end{exercise}

\begin{exercise}
\label{exercises-exercise-kunneth}
设 $X$ 与 $Y$ 为域 $k$ 上的概形
（可以随意假设 $X$ 与 $Y$ 性质良好，例如在 $k$ 上 qcqs 或射影）。
令 $X \times Y = X \times_{\Spec(k)} Y$，其投影为
$p : X \times Y \to X$ 与 $q : X \times Y \to Y$。
对拟凝聚 $\mathcal{O}_X$-模 $\mathcal{F}$ 和拟凝聚
$\mathcal{O}_Y$-模 $\mathcal{G}$，证明
$$
H^n(X \times Y,
p^*\mathcal{F} \otimes_{\mathcal{O}_{X \times Y}} q^*\mathcal{G}) =
\bigoplus\nolimits_{a + b = n}
H^a(X, \mathcal{F}) \otimes_k H^b(Y, \mathcal{G})
$$
或者只证明取维数后该等式成立。
“干净”的解答可获额外加分。
\end{exercise}

\begin{exercise}
\label{exercises-exercise-Oab}
设 $k$ 为域。令
% P12-ERR-0473
$X = \mathbf{P}|^1 \times \mathbf{P}^1$
为 $k$ 上射影直线与自身的乘积，投影为
$p : X \to \mathbf{P}^1_k$ 与 $q : X \to \mathbf{P}^1_k$。
令
$$
\mathcal{O}(a, b) = p^*\mathcal{O}_{\mathbf{P}^1_k}(a)
\otimes_{\mathcal{O}_X} q^*\mathcal{O}_{\mathbf{P}^1_k}(b)
$$
对所有 $i, a, b$，计算 $H^i(X, \mathcal{O}(a, b))$ 的维数。
提示：使用练习 \ref{exercises-exercise-kunneth}。
\end{exercise}











\section{上同调与 Hilbert 多项式}
\label{exercises-section-cohomology-hilbert-polynomials}

\begin{situation}
\label{exercises-situation-hilbert-polynomial}
设 $k$ 为域。设 $X = \mathbf{P}^n_k$ 为 $n$-维射影空间，
并设 $\mathcal{F}$ 为凝聚 $\mathcal{O}_X$-模。回顾
$$
\chi(X, \mathcal{F}) =
\sum\nolimits_{i = 0}^n (-1)^i \dim_k H^i(X, \mathcal{F})
$$
回顾，$\mathcal{F}$ 的{\it Hilbert 多项式}是函数
$$
t \longmapsto \chi(X, \mathcal{F}(t))
$$
还回顾
$\mathcal{F}(t) = \mathcal{F} \otimes_{\mathcal{O}_X} \mathcal{O}_X(t)$，
其中 $\mathcal{O}_X(t)$ 是结构层的第 $t$ 次扭曲，如构造篇定义
\ref{constructions-definition-twist}。
我们已在簇篇第 \ref{varieties-subsection-hilbert} 小节证明，
Hilbert 多项式是 $t$ 的多项式。
\end{situation}

\begin{exercise}
\label{exercises-exercise-hilbert-pol-easy}
在情形 \ref{exercises-situation-hilbert-polynomial} 中：
\begin{enumerate}
\item 若 $P(t)$ 是 $\mathcal{F}$ 的 Hilbert 多项式，
$\mathcal{F}(-13)$ 的 Hilbert 多项式是什么？
\item 若 $P_i$ 是 $\mathcal{F}_i$ 的 Hilbert 多项式，
$\mathcal{F}_1 \oplus \mathcal{F}_2$ 的 Hilbert 多项式是什么？
\item 若 $P_i$ 是 $\mathcal{F}_i$ 的 Hilbert 多项式，
且 $\mathcal{F}$ 是满射 $\mathcal{F}_1 \to \mathcal{F}_2$ 的核，
$\mathcal{F}$ 的 Hilbert 多项式是什么？
\end{enumerate}
\end{exercise}

\begin{exercise}
\label{exercises-exercise-find-given-hilbert-pol-dim-1}
在情形 \ref{exercises-situation-hilbert-polynomial} 中，假设 $n \geq 1$。
找出 Hilbert 多项式为 $t - 101$ 的凝聚层。
\end{exercise}

\begin{exercise}
\label{exercises-exercise-find-given-hilbert-pol-dim-2}
在情形 \ref{exercises-situation-hilbert-polynomial} 中，假设 $n \geq 2$。
找出 Hilbert 多项式为 $t^2/2 + t/2 - 1$ 的凝聚层。
（这有点棘手；只证明存在这样的凝聚层即可。）
\end{exercise}

\begin{exercise}
\label{exercises-exercise-bound-genus-in-degree}
在情形 \ref{exercises-situation-hilbert-polynomial} 中，假设 $n \geq 2$
且 $k$ 代数闭。设 $C \subset X$ 为维数 $1$ 的整闭子概形。
换言之，$C$ 是射影曲线。
设 $d t + e$ 为 $\mathcal{O}_C$ 作为 $X$ 上凝聚层的
Hilbert 多项式。
\begin{enumerate}
\item 给出 $e$ 的一个上界。（提示：利用
$\mathcal{O}_C(t)$ 仅在次数 $0$ 与 $1$ 有上同调，
并研究 $H^0(C, \mathcal{O}_C)$。）
\item 选取 $\mathcal{O}_X(1)$ 的整体截面 $s$，使它与 $C$
横截相交，即存在两两不同的闭点 $c_1, \ldots, c_r \in C$
以及短正合列
$$
0 \to \mathcal{O}_C \xrightarrow{s} \mathcal{O}_C(1) \to
\bigoplus\nolimits_{i = 1, \ldots, r} k_{c_i} \to 0
$$
其中 $k_{c_i}$ 是在 $c_i$ 处取值为 $k$ 的摩天层。
（这样的 $s$ 存在；请直接使用这一点。）
证明 $r = d$。（提示：扭曲该正合列并观察所得结果。）
\item 扭曲该短正合列后，对任意 $t$ 得到 $k$-线性映射
$\varphi_t : \Gamma(C, \mathcal{O}_C(t)) \to \bigoplus_{i = 1, \ldots, d} k$。
% P12-ERR-0474
证明若 $t \geq d - 1$，该映射为满射。
\item 用 $d$ 给出 $e$ 的一个下界。（提示：使用（3）的结果和消失性，
证明 $H^1(C, \mathcal{O}_C(d - 2)) = 0$。）
\end{enumerate}
\end{exercise}

\begin{exercise}
\label{exercises-exercise-three-quadrics-in-plane}
在情形 \ref{exercises-situation-hilbert-polynomial} 中，假设 $n = 2$。
设 $s_1, s_2, s_3 \in \Gamma(X, \mathcal{O}_X(2))$
为三个二次方程。考虑凝聚层
$$
\mathcal{F} = \Coker\left(\mathcal{O}_X(-2)^{\oplus 3}
\xrightarrow{s_1, s_2, s_3} \mathcal{O}_X\right)
$$
列出这样的 $\mathcal{F}$ 所有可能的 Hilbert 多项式。
（试着把射影平面中二次曲线的交可视化。）
\end{exercise}





\section{曲线}
\label{exercises-section-curves}

\begin{exercise}
\label{exercises-exercise-closed-subset-vanshing}
设 $k$ 为代数闭域，$X$ 为 $k$ 上的射影曲线。
设 $\mathcal{L}$ 为可逆 $\mathcal{O}_X$-模，并设
$s_0, \ldots, s_n \in H^0(X, \mathcal{L})$ 为 $\mathcal{L}$
的整体截面。证明存在自然闭子概形
$$
Z \subset
\mathbf{P}^n \times X
$$
使闭点 $((\lambda_0 : \ldots : \lambda_n), x)$ 属于 $Z$
当且仅当截面 $\lambda_0 s_0 + \ldots + \lambda_n s_n$
在 $x$ 处消失。（提示：在仿射局部描述 $Z$。）
\end{exercise}

\begin{exercise}
\label{exercises-exercise-diagonal-cartier}
设 $k$ 为代数闭域，$X$ 为 $k$ 上的光滑曲线。
设 $r \geq 1$。证明闭子集
$$
D \subset X \times X^r = X^{r + 1}
$$
具有可逆理想层；其闭点是满足某个 $i$ 下 $x = x_i$ 的元组
$(x, x_1, \ldots, x_r)$。换言之，证明 $D$ 是有效 Cartier 除子。
% P12-ERR-0475
提示：归约到 $r = 1$，并利用 $X$ 是光滑曲线这一点来研究对角线
（可查阅书籍）。
\end{exercise}

\begin{exercise}
\label{exercises-exercise-one-divisor-in-another}
设 $k$ 为代数闭域，$X$ 为 $k$ 上的光滑射影曲线。
设 $T$ 为 $k$ 上有限型概形，并设
$$
D_1 \subset X \times T
\quad\text{且}\quad
D_2 \subset X \times T
$$
为两个有效 Cartier 除子，使对 $t \in T$，纤维
$D_{i, t} \subset X_t$ 不稠密（即不含曲线 $X_t$ 的泛点）。
证明存在典范闭子概形 $Z \subset T$，使闭点 $t \in T$ 属于
$Z$，当且仅当 $D_1$、$D_2$ 的概形论纤维
$D_{1, t}$、$D_{2, t}$ 满足
$$
D_{1, t} \subset D_{2, t}
$$
（作为 $X_t$ 的闭子概形）。
提示：证明可以在缩小 $T$ 后假设 $T = \Spec(A)$ 仿射，
并存在仿射开集 $U \subset X$ 使 $D_i \subset U \times T$。
然后证明 $M_1 = \Gamma(D_1, \mathcal{O}_{D_1})$ 是有限局部自由
$A$-模（这里需要一些非平凡的代数——问问你的朋友）。
缩小 $T$ 后，可假设 $M_1$ 是自由 $A$-模。然后考察
$$
\Gamma(U \times T, \mathcal{I}_{D_2}) \to M_1 = A^{\oplus N}
$$
并把 $Z$ 定义为由该映射像中向量的系数所生成的理想切出的闭子概形。
说明这为何成立（这里可能还需要一点交换代数）。
\end{exercise}

\begin{exercise}
\label{exercises-exercise-closed-subset-vanshing-r}
设 $k$ 为代数闭域，$X$ 为 $k$ 上的光滑射影曲线。
设 $\mathcal{L}$ 为可逆 $\mathcal{O}_X$-模，并设
$s_0, \ldots, s_n \in H^0(X, \mathcal{L})$ 为 $\mathcal{L}$
的整体截面。设 $r \geq 1$。证明存在自然闭子概形
$$
Z \subset
\mathbf{P}^n \times X \times \ldots \times X = \mathbf{P}^n \times X^r
$$
使闭点 $((\lambda_0 : \ldots : \lambda_n), x_1, \ldots, x_r)$
属于 $Z$，当且仅当截面
$s_\lambda = \lambda_0 s_0 + \ldots + \lambda_n s_n$
在除子 $D = x_1 + \ldots + x_r$ 上消失，即截面 $s_\lambda$
属于 $\mathcal{L}(-D)$。
% P12-ERR-0476
提示：说明如何由练习 \ref{exercises-exercise-diagonal-cartier} 与
\ref{exercises-exercise-one-divisor-in-another} 的结论结合得出这一点。
\end{exercise}

\begin{exercise}
\label{exercises-exercise-effective}
设 $k$ 为代数闭域，$X$ 为 $k$ 上的光滑射影曲线。
设 $\mathcal{L}$ 为可逆 $\mathcal{O}_X$-模。
证明存在自然闭子集
$$
Z \subset X^r
$$
使 $X^r$ 的闭点 $(x_1, \ldots, x_r)$ 属于 $Z$，当且仅当
$\mathcal{L}(-x_1 - \ldots -x_r)$ 有非零整体截面。
提示：使用练习 \ref{exercises-exercise-closed-subset-vanshing-r}。
\end{exercise}

\begin{exercise}
\label{exercises-exercise-difference-effective}
设 $k$ 为代数闭域，$X$ 为 $k$ 上的光滑射影曲线。
设 $r \geq s$ 为整数。证明存在自然闭子集
$$
Z \subset X^r \times X^s
$$
使 $X^r \times X^s$ 的闭点
$(x_1, \ldots, x_r, y_1, \ldots, y_s)$ 属于 $Z$，当且仅当
$x_1 + \ldots + x_r - y_1 - \ldots - y_s$
线性等价于某个有效除子。
% P12-ERR-0477
提示：选取次数非常高的辅助可逆模 $\mathcal{L}$，使对任意次数
为 $r$ 的有效除子 $D$，$\mathcal{L}(-D)$ 都有非消失截面。
然后两次使用练习 \ref{exercises-exercise-effective} 的结论。
\end{exercise}

\begin{exercise}
\label{exercises-exercise-genus-7}
选择你最喜欢的代数闭域 $k$。尽你所能，确定某条亏格为 $7$ 的
曲线上可能存在的所有 $\mathfrak g^r_d$。同时还要试着：
\begin{enumerate}
\item 确定在哪些情形下 $\mathfrak g^r_d$ 无基点；
\item 确定在哪些情形下 $\mathfrak g^r_d$ 给出到
$\mathbf{P}^r$ 的闭嵌入。
\end{enumerate}
若假设曲线是“一般的”，完成同样的工作
（自行给出“一般”的概念——这可能比上面的问题容易）。
若假设曲线是超椭圆曲线，完成同样的工作。
若假设曲线是三角曲线（且不是超椭圆曲线），完成同样的工作。等等。
\end{exercise}





\section{模空间}
\label{exercises-section-moduli}

\noindent
本节考察代数几何对象的模空间的一些朴素处理方式。

\medskip\noindent
设 $k$ 为代数闭域。假设 $M$ 是 $k$ 上的一个模问题。
这里不精确定义其含义，但在每个练习中，我们的意思应当是清楚的。
要理解下面的内容，只需知道 $M$ 在 $k$ 上有哪些对象、$M$
在 $k$ 上的对象之间有哪些同构，以及簇上的 $M$-对象族是什么。
% P12-ERR-0478
若下列命题成立，就称 $M$ 的{\it 模数}为 $d \geq 0$：
\begin{enumerate}
\item 存在有限多个族 $X_i \to V_i$（$i = 1, \ldots, n$），
使 $M$ 在 $k$ 上的每个对象都同构于其中某个族的一条纤维，
且 $\max \dim(V_i) = d$；
\item 不可能用更小的 $d$ 做到这一点。
\end{enumerate}
这实际上只是文献中更好概念的一种非常粗略的近似。

\begin{exercise}
\label{exercises-exercise-number-moduli-lines-conics}
设 $k$ 为代数闭域，并设 $d \geq 1$、$n \geq 1$。
规定 $P^n$ 中 $d$ 次超曲面的模空间如下：
\begin{enumerate}
\item 一个对象是次数为 $d$ 的超曲面
$X \subset \mathbf{P}^n_k$；
\item 两个对象 $X \subset \mathbf{P}^n_k$ 与
$Y \subset \mathbf{P}^n_k$ 之间的一个同构，是满足
% P12-ERR-0479
$g(X) = Y$ 的元素 $g \in \text{PGL}_n(k)$；
\item 簇 $V$ 上的一个超曲面族，是闭子概形
$X \subset \mathbf{P}^n_V$，使对所有 $v \in V$，
$X \to V$ 的概形论纤维 $X_v$ 是 $\mathbf{P}^n_v$ 中的一条
% P12-ERR-0480
超曲面。
\end{enumerate}
计算下列模数（如果可以的话——它们会逐渐变难）：
\begin{enumerate}
\item $n = 1$ 而 $d$ 任意时的模数；
\item $n = 2$ 且 $d = 1$ 时的模数；
\item $n = 2$ 且 $d = 2$ 时的模数；
\item $n \geq 1$ 且 $d = 2$ 时的模数；
\item $n = 2$ 且 $d = 3$ 时的模数；
\item $n = 3$ 且 $d = 3$ 时的模数；
\item $n = 2$ 且 $d = 4$ 时的模数。
\end{enumerate}
\end{exercise}

\begin{exercise}
\label{exercises-exercise-number-moduli-hyperelliptic}
设 $k$ 为代数闭域，并设 $g \geq 2$。
规定亏格 $g$ 的超椭圆曲线的模空间如下：
\begin{enumerate}
\item 一个对象是亏格 $g$ 的光滑射影超椭圆曲线 $X$；
\item 两个对象 $X$ 与 $Y$ 之间的一个同构，是 $k$ 上概形的同构
$X \to Y$；
\item 簇 $V$ 上亏格 $g$ 的一个超椭圆曲线族，是一个
真平坦\footnote{去掉此假设也不会改变问题的答案。}态射
% P12-ERR-0481
$X \to Y$，使 $X \to V$ 的所有概形论纤维都是亏格 $g$ 的
光滑射影超椭圆曲线。
\end{enumerate}
证明其模数为 $2g - 1$。
\end{exercise}





\section{整体 Ext}
\label{exercises-section-global-exts}

\begin{exercise}
\label{exercises-exercise-ext-line-plane-p3}
设 $k$ 为域，$X = \mathbf{P}^3_k$。设 $L \subset X$ 与
$P \subset X$ 分别为直线和平面，把它们看作分别由 $1$ 个和
$2$ 个线性方程切出的闭子概形。对所有 $i$，计算
$$
\Ext^i_X(\mathcal{O}_L, \mathcal{O}_P)
$$
的维数。务必同时处理 $L$ 包含在 $P$ 中以及 $L$ 不包含在
$P$ 中的情形。
\end{exercise}

\begin{exercise}
\label{exercises-exercise-ext-point}
设 $k$ 为域，$X = \mathbf{P}^n_k$。
设 $Z \subset X$ 为一个闭的 $k$-有理点，并把它看作闭子概形；
例如，齐次坐标为 $(1 : 0 : \ldots : 0)$ 的点。
对所有 $i$，计算
$$
\Ext^i_X(\mathcal{O}_Z, \mathcal{O}_Z)
$$
的维数。
\end{exercise}

\begin{exercise}
\label{exercises-exercise-ext-pairings}
设 $X$ 为环空间。对 $\mathcal{O}_X$-模
$\mathcal{F}, \mathcal{G}, \mathcal{H}$ 定义杯积映射
$$
\Ext^i_X(\mathcal{G}, \mathcal{H})
\times
\Ext^j_X(\mathcal{F}, \mathcal{G})
\longrightarrow
\Ext^{i + j}_X(\mathcal{F}, \mathcal{H})
$$
（提示：这是一个极其一般的构造。）
\end{exercise}

\begin{exercise}
\label{exercises-exercise-move-out-locally-free}
设 $X$ 为环空间。设 $\mathcal{E}$ 为有限局部自由
$\mathcal{O}_X$-模，其对偶为
$\mathcal{E}^\vee = \SheafHom_{\mathcal{O}_X}(\mathcal{E}, \mathcal{O}_X)$。
证明下列命题：
\begin{enumerate}
\item $\SheafExt^i_{\mathcal{O}_X}(
\mathcal{F} \otimes_{\mathcal{O}_X} \mathcal{E}, \mathcal{G}) =
\SheafExt^i_{\mathcal{O}_X}(
\mathcal{F},
\mathcal{E}^\vee \otimes_{\mathcal{O}_X} \mathcal{G}) =
\SheafExt^i_{\mathcal{O}_X}(
\mathcal{F}, \mathcal{G}) \otimes_{\mathcal{O}_X} \mathcal{E}^\vee$；
\item $\Ext^i_X(
\mathcal{F} \otimes_{\mathcal{O}_X} \mathcal{E}, \mathcal{G}) =
\Ext^i_X(\mathcal{F},
\mathcal{E}^\vee \otimes_{\mathcal{O}_X} \mathcal{G})$。
\end{enumerate}
这里 $\mathcal{F}$ 与 $\mathcal{G}$ 是 $\mathcal{O}_X$-模。
推出
$$
\Ext^i_X(\mathcal{E}, \mathcal{G}) =
H^i(X, \mathcal{E}^\vee \otimes_{\mathcal{O}_X} \mathcal{G})
$$
\end{exercise}

\begin{exercise}
\label{exercises-exercise-trace-on-ext}
设 $X$ 为环空间，$\mathcal{E}$ 为有限局部自由
$\mathcal{O}_X$-模。对所有 $i$ 构造迹映射
$$
\Ext^i_X(\mathcal{E}, \mathcal{E}) \to H^i(X, \mathcal{O}_X)
$$
把它推广为对任意 $\mathcal{O}_X$-模 $\mathcal{F}$ 的迹映射
$$
\Ext^i_X(\mathcal{E}, \mathcal{E} \otimes_{\mathcal{O}_X} \mathcal{F})
\to H^i(X, \mathcal{F})
$$
\end{exercise}

\begin{exercise}
\label{exercises-exercise-duality}
设 $k$ 为域，$X = \mathbf{P}^d_k$。令
$\omega_{X/k} = \mathcal{O}_X(-d - 1)$。
证明：对有限局部自由模 $\mathcal{E}$、$\mathcal{F}$，
Ext 上的杯积与 Ext 上的迹映射结合得到
$$
\Ext^i_X(\mathcal{E}, \mathcal{F} \otimes_{\mathcal{O}_X} \omega_{X/k})
\times
\Ext^{d - i}_X(\mathcal{F}, \mathcal{E})
\to
\Ext_X^d(\mathcal{F}, \mathcal{F} \otimes_{\mathcal{O}_X} \omega_{X/k})
\to
H^d(X, \omega_{X/k}) = k
$$
是非退化配对。提示：你可以在此设定中重新证明对偶性，
也可以归约到层上同调，再应用讲课中证明的 Serre 对偶定理。
\end{exercise}








\section{除子}
\label{exercises-section-divisors}

\noindent
为方便起见，我们把所有相关定义汇集于此。

\begin{definition}
\label{exercises-definition-divisor}
以下始终令 $S$ 为任意概形，并令 $X$ 为 Noether 整概形。
\begin{enumerate}
\item $X$ 上的一个{\it Weil 除子}是素除子 $Z_i$ 的整数系数形式线性组合
$\Sigma n_i[Z_i]$。
\item 所谓{\it 素除子}，是整闭子概形 $Z \subset X$，其泛点
$\xi \in Z$ 满足 ${\mathcal O}_{X, \xi}$ 的维数为 $1$。
当 $\xi \in Z \subset X$ 如上时，我们使用记号
${\mathcal O}_{X, Z} = {\mathcal O}_{X, \xi}$。
注意，${\mathcal O}_{X, Z} \subset K(X)$ 是 $X$ 的函数域的子环。
\item 与有理函数 $f \in K(X)^\ast$ {\it 相伴的 Weil 除子}
是和 $\Sigma v_Z(f)[Z]$。这里 $v_Z(f)$ 定义如下：
\begin{enumerate}
\item 若 $f \in {\mathcal O}_{X, Z}^\ast$，则 $v_Z(f) = 0$。
\item 若 $f \in {\mathcal O}_{X, Z}$，则
$$
v_Z(f) = \text{length}_{{\mathcal O}_{X, Z}}({\mathcal O}_{X, Z}/(f)).
$$
\item 若 $f = \frac{a}{b}$，其中
$a, b \in {\mathcal O}_{X, Z}$，则
$$
v_Z(f) = \text{length}_{{\mathcal O}_{X, Z}}({\mathcal O}_{X, Z}/(a)) -
\text{length}_{{\mathcal O}_{X, Z}}({\mathcal O}_{X, Z}/(b)).
$$
\end{enumerate}
\item 概形 $S$ 上的一个{\it 有效 Cartier 除子}是闭子概形
$D \subset S$，使得每个点 $d\in D$ 都有 $S$ 中的仿射开邻域
% P12-ERR-0482
$\Spec(A) = U \subset S$，并且 $D \cap U = \Spec(A/(f))$，
其中 $f \in A$ 为非零因子。
\item 与 Noether 整概形 $X$ 的有效 Cartier 除子 $D \subset X$
{\it 相伴的 Weil 除子 $[D]$}定义为和
$\Sigma v_Z(D)[Z]$，其中 $v_Z(D)$ 定义如下：
\begin{enumerate}
\item 若 $Z$ 的泛点 $\xi$ 不属于 $D$，则 $v_Z(D) = 0$。
\item 若 $Z$ 的泛点 $\xi$ 属于 $D$，则
$$
v_Z(D) = \text{length}_{{\mathcal O}_{X, Z}}({\mathcal O}_{X, Z}/(f))
$$
其中 $f \in {\mathcal O}_{X, Z} = {\mathcal O}_{X, \xi}$ 是在
$\xi$ 的某个仿射邻域中定义 $D$ 的非零因子
（如上面的（4））。
\end{enumerate}
\item 设 $S$ 为概形。{\it 全商环层 ${\mathcal K}_S$}是
${\mathcal O}_S$-代数层，它是如下预层 ${\mathcal K}'$ 的层化。
对开集 $U \subset S$，令
${\mathcal K}'(U) = S_U^{-1}{\mathcal O}_S(U)$，其中
$S_U \subset {\mathcal O}_S(U)$ 是由如下截面组成的乘法子集：
$f \in {\mathcal O}_S(U)$，且对每个 $u\in U$，$f$ 在
${\mathcal O}_{S, u}$ 中的芽都是非零因子。
特别地，$S_U$ 的元素全都是非零因子。因此
${\mathcal O}_S$ 是 ${\mathcal K}_S$ 的子层，并得到短正合列
$$
0 \to {\mathcal O}_S^\ast \to {\mathcal K}_S^\ast \to
{\mathcal K}_S^\ast/{\mathcal O}_S^\ast \to 0.
$$
\item 概形 $S$ 上的一个{\it Cartier 除子}是商层
${\mathcal K}_S^\ast/{\mathcal O}_S^\ast$ 的整体截面。
\item 在 Noether 整概形 $X$ 上，与 Cartier 除子
$\tau \in \Gamma(X, {\mathcal K}_X^\ast/{\mathcal O}_X^\ast)$
{\it 相伴的 Weil 除子}是和 $\Sigma v_Z(\tau)[Z]$，其中
% P12-ERR-0483
$v_Z(\tau)$ 按下列步骤定义：
\begin{enumerate}
\item 若 $\tau$ 在 $Z$ 的泛点 $\xi$ 处的芽为零——换言之，
$\tau$ 在茎 $({\mathcal K}^\ast/{\mathcal O}^\ast)_\xi$ 中的像
为“零”——则 $v_Z(\tau) = 0$。
\item 找出仿射开邻域 $\Spec(A) = U \subset X$，使
$\tau|_U$ 是某个截面 $f \in {\mathcal K}(U)$ 的像，
而且 $f = a/b$，其中 $a, b \in A$。然后令
$$
v_Z(f) = \text{length}_{{\mathcal O}_{X, Z}}({\mathcal O}_{X, Z}/(a)) -
\text{length}_{{\mathcal O}_{X, Z}}({\mathcal O}_{X, Z}/(b)).
$$
\end{enumerate}
\end{enumerate}
\end{definition}

\begin{remarks}
\label{exercises-remarks-divisors}
下面是一些简单的说明。
\begin{enumerate}
\item 在 Noether 整概形 $X$ 上，层 ${\mathcal K}_X$ 是取值为
函数域 $K(X)$ 的常层。
\item 为使上述定义有意义，需要证明：对一维 Noether 局部整环
${\mathcal O}$ 的任意一对非零元素 $(a, b)$，有
$$
\text{length}_{\mathcal O}({\mathcal O}/(ab)) =
\text{length}_{\mathcal O}({\mathcal O}/(a)) +
\text{length}_{\mathcal O}({\mathcal O}/(b))
$$
这将在讲课中完成。
\end{enumerate}
\end{remarks}

\begin{exercise}
\label{exercises-exercise-effective-cartier-cartier}
（在任意概形上。）描述如何给一个有效 Cartier 除子指派
一个 Cartier 除子。
\end{exercise}

\begin{exercise}
\label{exercises-exercise-rational-function-cartier}
（在整概形上。）描述如何给有理函数 $f$ 指派 Cartier 除子
$D$，使与 $D$ 相伴的 Weil 除子和与 $f$ 相伴的 Weil 除子相同。
（这很显然。）
\end{exercise}

\begin{exercise}
\label{exercises-exercise-weil-not-cartier}
给出簇上的 Weil 除子的例子，使它不是与任何 Cartier 除子相伴的
Weil 除子。
\end{exercise}

\begin{exercise}
\label{exercises-exercise-weil-Q-cartier}
给出簇上的 Weil 除子 $D$ 的例子，使它不是与任何 Cartier 除子
相伴的 Weil 除子，但对某个 $n > 1$，$nD$ 是与一个 Cartier
除子相伴的 Weil 除子。
\end{exercise}

\begin{exercise}
\label{exercises-exercise-weil-not-Q-cartier}
给出簇上的 Weil 除子 $D$ 的例子，使它不是与任何 Cartier 除子
相伴的 Weil 除子，并且对任意 $n > 1$，$nD$ 都不是与任何
Cartier 除子相伴的 Weil 除子。
（提示：考虑一个锥，例如 $\mathbf{A}^4_k$ 中的
$X : xy - zw = 0$。试着证明 $D = [x = 0, z = 0]$ 可行。）
\end{exercise}

\begin{exercise}
\label{exercises-exercise-cartier-not-difference-effective-cartier}
在域上有限型分离概形 $X$ 上：
给出一个 Cartier 除子，使它不能写成两个有效 Cartier 除子的差。
提示：找一个没有任何非空有效 Cartier 除子的 $X$，例如
\cite[III Exercise 5.9]{H} 中构造的概形。
甚至有 $X$ 为簇的例子——即练习 \ref{exercises-exercise-nonprojective}
中的簇。
\end{exercise}

\begin{exercise}
\label{exercises-exercise-nonprojective}
非射影真簇的例子。设 $k$ 为域。设 $L \subset \mathbf{P}^3_k$
为直线，$C \subset \mathbf{P}^3_k$ 为非奇异二次曲线。
假设 $C \cap L = \emptyset$。选取同构 $\varphi : L \to C$。
令 $X$ 为通过 $\varphi$ 将 $C$ 与 $L$ 粘合所得的 $k$-簇。
换言之，存在满的真双有理态射
$$
\pi : \mathbf{P}^3_k \longrightarrow X
$$
以及开集 $U \subset X$，使
$\pi : \pi^{-1}(U) \to U$ 是同构，
$\pi^{-1}(U) = \mathbf{P}^3_k \setminus (L \cup C)$，
且 $\pi|_L = \pi|_C \circ \varphi$。
（这些条件尚不能唯一确定 $X$。为此还需指定 $X$ 在
$Z = X \setminus U$ 的点处的结构层。）
证明 $X$ 存在、是真簇，但不是射影的。
（提示：存在性使用练习 \ref{exercises-exercise-prepare-glueing} 的结论；
非射影性使用 $\Pic(\mathbf{P}^3_k) = \mathbf{Z}$，证明 $X$
不可能有充足可逆层。）
\end{exercise}




\section{微分}
\label{exercises-section-differentials}

\noindent
{\bf 定义与结果。}K\"ahler 微分。
\begin{enumerate}
\item 设 $R \to A$ 为环映射。$A$ 在 $R$ 上的
{\it K\"ahler 微分模}记作 $\Omega_{A/R}$。
它由元素 $\text{d}a$（$a \in A$）生成，并满足关系：
$$
\text{d}(a_1 + a_2) = \text{d}a_1 + \text{d}a_2,\quad
\text{d}(a_1a_2) = a_1\text{d}a_2 + a_2\text{d}a_1,\quad
\text{d}r = 0
$$
典范泛 $R$-导子 $\text{d} : A \to \Omega_{A/R}$
把 $a\mapsto \text{d}a$。
\item 考虑定义理想 $I$ 的短正合列
$$
0 \to I \to A \otimes_R A \to A \to 0
$$
存在典范导子 $\text{d} : A \to I/I^2$，把 $a$ 映到
$a \otimes 1 - 1 \otimes a$ 的类。这给出了 $A$ 在 $R$ 上的
微分模的另一种表示，换言之，
$$
(I/I^2, \text{d}) \cong (\Omega_{A/R}, \text{d}).
$$
\item 对满足 $S_R \subset R$ 的像包含在 $S_A \subset A$ 中的
乘法子集 $S_R$、$S_A$，有
$$
\Omega_{S_A^{-1}A / S_R^{-1}R} =
S_A^{-1}\Omega_{A/R}.
$$
\item 若 $A$ 是有限表现 $R$-代数，则 $\Omega_{A/R}$ 是有限表现
$A$-模。因此，在这种情形下，$\Omega_{A/R}$ 的
{\it Fitting} 理想有定义。
\item 设 $f : X \to S$ 为概形态射。存在拟凝聚
${\mathcal O}_X$-模层 $\Omega_{X/S}$ 以及
% P12-ERR-0484
${\mathcal O}_S$-线性导子
$$
\text{d} : {\mathcal O}_X \longrightarrow \Omega_{X/S}
$$
使得对满足 $f(U) \subset V$ 的任意仿射开集
$\Spec(A) = U \subset X$、$\Spec(R) = V \subset S$，
都有
$$
\Gamma(\Spec(A), \Omega_{X/S}) = \Omega_{A/R}
$$
并且与 $\text{d}$ 相容。
\end{enumerate}

\begin{exercise}
\label{exercises-exercise-dual-numbers}
设 $k[\epsilon]$ 为域 $k$ 上的对偶数环，即
$\epsilon^2 = 0$。
\begin{enumerate}
\item 考虑环映射
$$
R = k[\epsilon] \to A = k[x, \epsilon]/(\epsilon x)
$$
证明 $\Omega_{A/R}$ 的 Fitting 理想为
（从第零 Fitting 理想开始）
$$
(\epsilon), A, A, \ldots
$$
\item 考虑映射
$R = k[t] \to A = k[x, y, t]/(x(y-t)(y-1), x(x-t))$。
证明 $\Omega_{A/R}$ 在 $A$ 中的 Fitting 理想为
（为简单起见，假设 $k$ 的特征为零）
$$
x(2x-t)(2y-t-1)A, \ (x, y, t)\cap (x, y-1, t), \ A, \ A, \ldots
$$
% P12-ERR-0485
因此，第 $0$-个 Fitting 理想由 $A$ 的单个元素切出，
第 $1$ 个 Fitting 理想定义 $\Spec(A)$ 的两个闭点，
其余理想全都平凡。
\item 考虑映射 $R = k[t] \to A = k[x, y, t]/(xy-t^n)$。
计算 $\Omega_{A/R}$ 的 Fitting 理想。
\end{enumerate}
\end{exercise}

\begin{remark}
\label{exercises-remark-fitting-omega-not-sings}
$\Omega_{X/S}$ 的第 $k$ 个 Fitting 理想常用于定义态射
$X \to S$ 的奇异概形，其中 $X$ 在 $S$ 上的相对维数为 $k$。
% P12-ERR-0486
但正如（a）表明的，当族的纤维维数不是“常数”（例如它不平坦）时，
这样做必须小心。正如（b）表明的，平坦性也不保证这种做法有效
（而且，这确实是一个平坦族）。在“良好情形”中——例如（c）——
对于曲线族，我们期望第 $0$ 个 Fitting 理想为零，而第 $1$ 个
Fitting 理想在概形论意义下定义奇异轨迹。
\end{remark}

\begin{exercise}
\label{exercises-exercise-formally-smooth}
假设 $R$ 是环且
$$
A = R[x_1, \ldots, x_n]/(f_1, \ldots, f_n).
$$
注意，我们假设 $A$ 的表示中方程数与变量数相同。因此偏导矩阵
$$
( \partial f_i / \partial x_j )
$$
是 $n \times n$ 矩阵，即方阵。假设其行列式作为 $A$ 的元素可逆。
注意，在有 $n$ 个生成元与 $n$ 个关系的本情形中，这恰好是条件
$\Omega_{A/R} = (0)$。设 $\pi : B' \to B$ 为 $R$-代数的满射，
其核 $J$ 的平方为零（作为 $B'$ 的理想）。
设 $\varphi : A \to B$ 为 $R$-代数同态。
证明存在唯一的 $R$-代数同态 $\varphi' : A \to B'$，
使得 $\varphi = \pi \circ \varphi'$。
\end{exercise}

\begin{exercise}
\label{exercises-exercise-formally-smooth-one-equation}
把练习 \ref{exercises-exercise-formally-smooth} 的结论推广到
$A = R[x, y]/(f)$ 的情形。
\end{exercise}

\begin{exercise}
\label{exercises-exercise-Jacobian-criterion}
设 $k$ 为域，
$f_1, \ldots, f_c \in k[x_1, \ldots, x_n]$，并令
$A = k[x_1, \ldots, x_n]/(f_1, \ldots, f_c)$。
假设 $f_j(0, \ldots, 0) = 0$。这意味着
$\mathfrak m = (x_1, \ldots, x_n)A$ 是极大理想。
证明：若矩阵
$$
(\partial f_j/ \partial x_i)|_{(x_1, \ldots, x_n) = (0, \ldots, 0)}
$$
的秩为 $c$，则局部环 $A_\mathfrak m$ 正则。
在这种情形下，$A_\mathfrak m$ 的维数是多少？
给出 $A_\mathfrak m$ 正则但上述秩小于 $c$ 的例子，
以证明逆命题为假；在你的例子中，$A_\mathfrak m$ 的维数是多少？
\end{exercise}





\section{概形，2007 年秋季期末考试}
\label{exercises-section-final-exam-fall-2007}

\noindent
% P12-ERR-0487
这些是哥伦比亚大学 2007 年春季一门概形课程的期末考试题。

\begin{exercise}[定义]
\label{exercises-exercise-definitions}
给出下列概念的定义。
\begin{enumerate}
\item $X$ 是{\it 概形}
\item 概形态射 $f : X \to Y$ 是{\it 有限的}
% P12-ERR-0488
\item 概形态射 $f : X \to Y$ 是{\it 有限型的}
\item 概形 $X$ 是{\it Noether 的}
\item 概形 $X$ 上的 ${\mathcal O}_X$-模 ${\mathcal L}$ 是
{\it 可逆的}
\item 代数闭域上非奇异射影曲线的{\it 亏格}
\end{enumerate}
\end{exercise}

\begin{exercise}
\label{exercises-exercise-kill-global-sections}
设 $X = \Spec({\mathbf Z}[x, y])$，并设 ${\mathcal F}$ 为拟凝聚
${\mathcal O}_X$-模。假设 ${\mathcal F}$ 限制到标准仿射开集
$D(x)$ 后为零。
\begin{enumerate}
\item 证明 ${\mathcal F}$ 的每个整体截面 $s$ 都被 $x$ 的某个幂
消去，即对某个 $n\in {\mathbf N}$ 有 $x^ns = 0$。
\item 若不假设 ${\mathcal F}$ 拟凝聚，你认为同一结论仍成立吗？
\end{enumerate}
\end{exercise}

\begin{exercise}
\label{exercises-exercise-empty-fibre-empty}
假设 $X \to \Spec(R)$ 是真态射，$R$ 是剩余域为 $k$ 的离散赋值环。
假设 $X \times_{\Spec(R)} \Spec(k)$ 是空概形。
证明 $X$ 是空概形。
\end{exercise}

\begin{exercise}
\label{exercises-exercise-curve-p1-p1}
考虑复数域 ${\mathbf C}$ 上的射影\footnote{射影嵌入为
$((X_0, X_1), (Y_0, Y_1))\mapsto (X_0Y_0, X_0Y_1, X_1Y_0, X_1Y_1)$，
换言之，$(x, y)\mapsto (1, y, x, xy)$。}簇
$$
{\mathbf P}^1 \times {\mathbf P}^1 = {\mathbf P}^1_{{\mathbf C}}
\times_{\Spec({\mathbf C})} {\mathbf P}^1_{\mathbf C}
$$
它被四个仿射块覆盖，分别对应于 ${\mathbf P}^1_{\mathbf C}$
两个因子的标准仿射块对。例如，在第一个因子上使用齐次坐标
$X_0, X_1$，在第二个因子上使用 $Y_0, Y_1$。
令 $x = X_1/X_0$、$y = Y_1/Y_0$。则四个仿射开块是下列环的谱：
$$
{\mathbf C}[x, y], \quad
{\mathbf C}[x^{-1}, y], \quad
{\mathbf C}[x, y^{-1}], \quad
{\mathbf C}[x^{-1}, y^{-1}].
$$
设 $X \subset {\mathbf P}^1 \times {\mathbf P}^1$ 为闭子概形，
它是第一个仿射块中由方程
$$
y^3(x^4 + 1) = x^4 -1.
$$
给出的闭子集的闭包。
\begin{enumerate}
\item 证明 $X$ 包含在四个仿射开块中第一个与最后一个的并中。
\item 证明 $X$ 是非奇异射影曲线。
% P12-ERR-0489
\item 考虑态射 $pr_2 : X \to {\mathbf P}^1$（到第一个因子的投影）。
在第一个仿射块上，它是映射 $(x, y) \mapsto x$。
简要说明为何它的次数为 $3$。
\item 计算映射 $pr_2 : X \to {\mathbf P}^1$ 的分歧点与分歧指数。
\item 计算 $X$ 的亏格。
\end{enumerate}
\end{exercise}

\begin{exercise}
\label{exercises-exercise-finite-type-over-Z}
设 $X \to \Spec({\mathbf Z})$ 为有限型态射。
假设有无穷多个素数 $p$，使得
$X \times_{\Spec({\mathbf Z})} \Spec({\mathbf F}_p)$ 非空。
\begin{enumerate}
\item 证明 $X \times_{\Spec({\mathbf Z})}\Spec(\mathbf{Q})$ 非空。
\item 若把“有限型”条件换成“局部有限型”条件，你认为同一结论仍成立吗？
\end{enumerate}
\end{exercise}





\section{概形，2009 年春季期末考试}
\label{exercises-section-final-exam-spring-2009}

\noindent
这些是哥伦比亚大学 2009 年春季一门概形课程的期末考试题。

\begin{exercise}
\label{exercises-exercise-Noetherian-coherent}
设 $X$ 为 Noether 概形，$\mathcal{F}$ 为 $X$ 上的凝聚层，
$x \in X$ 为一点。假设
$\text{Supp}(\mathcal{F}) = \{ x \}$。
\begin{enumerate}
\item 证明 $x$ 是 $X$ 的闭点。
\item 证明 $H^0(X, \mathcal{F})$ 非零。
\item 证明 $\mathcal{F}$ 由整体截面生成。
\item 证明对 $p > 0$，$H^p(X, \mathcal{F}) = 0$。
\end{enumerate}
\end{exercise}

\begin{remark}
\label{exercises-remark-invertible-projective-space}
设 $k$ 为域，并令
$\mathbf{P}^2_k = \text{Proj}(k[X_0, X_1, X_2])$。
$\mathbf{P}^2_k$ 上的任意可逆层都同构于某个
$n \in \mathbf{Z}$ 所对应的
$\mathcal{O}_{\mathbf{P}^2_k}(n)$。回顾
$$
\Gamma(\mathbf{P}^2_k, \mathcal{O}_{\mathbf{P}^2_k}(n)) =
k[X_0, X_1, X_2]_n
$$
是多项式环的次数 $n$ 部分。对 $\mathbf{P}^2_k$ 上的拟凝聚层
$\mathcal{F}$，照例令
$\mathcal{F}(n) =
\mathcal{F}
\otimes_{\mathcal{O}_{\mathbf{P}^2_k}}
\mathcal{O}_{\mathbf{P}^2_k}(n)$。
\end{remark}

\begin{exercise}
\label{exercises-exercise-nonsplit-vectorbundle}
设 $k$ 为域。设 $\mathcal{E}$ 为 $\mathbf{P}^2_k$ 上的向量丛，
即有限局部自由 $\mathcal{O}_{\mathbf{P}^2_k}$-模。
% P12-ERR-0490
若 $\mathcal{E}$ 同构于可逆
$\mathcal{O}_{\mathbf{P}^2_k}$-模的直和，则称 $\mathcal{E}$
{\it 分裂}。
\begin{enumerate}
\item 证明 $\mathcal{E}$ 分裂当且仅当 $\mathcal{E}(n)$ 分裂。
\item 证明若 $\mathcal{E}$ 分裂，则对所有 $n \in \mathbf{Z}$，
$H^1({\mathbf{P}^2_k}, \mathcal{E}(n)) = 0$。
\item 设
$$
\varphi :
\mathcal{O}_{\mathbf{P}^2_k}
\longrightarrow
\mathcal{O}_{\mathbf{P}^2_k}(1) \oplus
\mathcal{O}_{\mathbf{P}^2_k}(1) \oplus
\mathcal{O}_{\mathbf{P}^2_k}(1)
$$
由线性形式
$L_0, L_1, L_2 \in \Gamma(\mathbf{P}^2_k, \mathcal{O}_{\mathbf{P}^2_k}(1))$
给出。假设对某个 $i$ 有 $L_i \not = 0$。
$L_0, L_1, L_2$ 应满足什么条件，才能使 $\varphi$ 的余核为向量丛？
为什么？
% P12-ERR-0491
\item 给出这样的 $\varphi$ 的例子。
\item 证明 $\Coker(\varphi)$ 不分裂（若它是向量丛）。
\end{enumerate}
\end{exercise}

\begin{remark}
\label{exercises-remark-recall-dimension-theory}
可以自由使用下列维数理论事实（若需要，还可自行添加）。
\begin{enumerate}
\item 概形的维数是不可约闭子集链长度的上确界。
\item 域上有限型概形的维数是其仿射开集维数的最大值。
\item Noether 概形的维数是其不可约分支维数的最大值。
\item 仿射概形的维数等于对应环的维数。
\item 设 $k$ 为域，$A$ 为有限型 $k$-代数。
% P12-ERR-0492
若 $A$ 是整环且 $x \not = 0$，则
$\dim(A) = \dim(A/xA) + 1$。
\end{enumerate}
\end{remark}

\begin{exercise}
\label{exercises-exercise-irreducible-fibres-same-dimension-irreducible}
设 $k$ 为域，$X$ 为 $k$ 上射影既约概形。
设 $f : X \to \mathbf{P}^1_k$ 为 $k$ 上的概形态射。
假设存在整数 $d \geq 0$，使对每个点
$t \in \mathbf{P}^1_k$，纤维 $X_t = f^{-1}(t)$
都不可约且维数为 $d$。（回顾：不可约空间非空。）
\begin{enumerate}
\item 证明 $\dim(X) = d + 1$。
\item 设 $X_0 \subset X$ 为 $X$ 的一个维数为 $d + 1$
的不可约分支。证明对每个 $t \in \mathbf{P}^1_k$，
纤维 $X_{0, t}$ 的维数为 $d$。
\item 由上面各项，关于 $X_t$ 与 $X_{0, t}$ 能得出什么结论？
\item 证明 $X$ 不可约。
\end{enumerate}
\end{exercise}

\begin{remark}
\label{exercises-remark-chi}
给定域 $k$ 上的射影概形 $X$ 以及 $X$ 上的凝聚层
$\mathcal{F}$，令
$$
\chi(X, \mathcal{F}) =
\sum\nolimits_{i \geq 0} (-1)^i\dim_k H^i(X, \mathcal{F}).
$$
\end{remark}

\begin{exercise}
\label{exercises-exercise-complete-intersection}
设 $k$ 为域。写
$\mathbf{P}^3_k = \text{Proj}(k[X_0, X_1, X_2, X_3])$。
设 $C \subset \mathbf{P}^3_k$ 为
{\it $(5, 6)$ 型完全交曲线}。这意味着存在
$F \in k[X_0, X_1, X_2, X_3]_5$ 与
$G \in k[X_0, X_1, X_2, X_3]_6$，使
$$
C = \text{Proj}(k[X_0, X_1, X_2, X_3]/(F, G))
$$
是维数为 $1$ 的簇。（簇意味着既约且不可约，但若愿意，
可以直接假设 $C$ 非奇异。）
令 $i : C \to \mathbf{P}^3_k$ 为相应的闭浸入。
完全交还意味着
$$
\xymatrix{
0 \ar[r] &
\mathcal{O}_{\mathbf{P}^3_k}(-11)
\ar[r]^-{
\left(
\begin{matrix}
-G \\
F
\end{matrix}
\right)
} &
\mathcal{O}_{\mathbf{P}^3_k}(-5) \oplus \mathcal{O}_{\mathbf{P}^3_k}(-6)
\ar[r]^-{(F, G)} &
\mathcal{O}_{\mathbf{P}^3_k} \ar[r] &
i_*\mathcal{O}_C \ar[r] &
0
}
$$
是层的正合列。请使用这些事实：
\begin{enumerate}
\item 对任意 $n \in \mathbf{Z}$ 计算
$\chi(C, i^*\mathcal{O}_{\mathbf{P}^3_k}(n))$；
\item 计算 $H^1(C, \mathcal{O}_C)$ 的维数。
\end{enumerate}
\end{exercise}

\begin{exercise}
\label{exercises-exercise-glueing}
设 $k$ 为域。考虑环
\begin{align*}
A & = k[x, y]/(xy) \\
B & = k[u, v]/(uv) \\
C & = k[t, t^{-1}] \times k[s, s^{-1}]
\end{align*}
以及 $k$-代数映射
$$
\begin{matrix}
A \longrightarrow C, &
x \mapsto (t, 0), &
y \mapsto (0, s) \\
B \longrightarrow C, &
u \mapsto (t^{-1}, 0), &
v \mapsto (0, s^{-1})
\end{matrix}
$$
事实是，这些映射诱导同构 $A_{x + y} \to C$ 与
$B_{u + v} \to C$。因此映射 $A \to C$ 与 $B \to C$
把 $\Spec(C)$ 分别等同于 $\Spec(A)$ 与 $\Spec(B)$ 的开子集。
令 $X$ 为沿 $\Spec(C)$ 粘合 $\Spec(A)$ 与 $\Spec(B)$
所得的概形：
$$
X = \Spec(A) \amalg_{\Spec(C)} \Spec(B).
$$
正如课程中所见，这样的概形存在，并有仿射开集
$\Spec(A) \subset X$ 与 $\Spec(B) \subset X$；
它们的交恰好是 $\Spec(C)$，通过上述映射与二者各自的一个开集等同。
\begin{enumerate}
\item 为什么 $X$ 分离？
\item 为什么 $X$ 在 $k$ 上有限型？
\item 计算 $H^1(X, \mathcal{O}_X)$，或者说，它的维数是多少？
\item 如何以更几何的方式描述 $X$？
\end{enumerate}
\end{exercise}





\section{概形，2010 年秋季期末考试}
\label{exercises-section-final-exam-fall-2010}

\noindent
这些是哥伦比亚大学 2010 年秋季一门概形课程的期末考试题。

\begin{exercise}[定义]
\label{exercises-exercise-definitions-fall-2010}
给出下列概念的定义。
\begin{enumerate}
\item 分离概形；
\item 概形的拟紧态射；
\item 概形的仿射态射；
\item 环的乘法子集；
\item Noether 概形；
\item 簇。
\end{enumerate}
\end{exercise}

\begin{exercise}
\label{exercises-exercise-prime-avoidance}
素理想回避。
\begin{enumerate}
\item 设 $A$ 为环，$I \subset A$ 为理想，并设
$\mathfrak q_1$、$\mathfrak q_2$ 为满足
$I \not \subset \mathfrak q_i$ 的素理想。
证明 $I \not \subset \mathfrak q_1 \cup \mathfrak q_2$。
\item （1）的几何解释是什么？
\item 设 $X = \text{Proj}(S)$，其中 $S$ 为某个分次环。
设 $x_1, x_2 \in X$。证明存在同时包含 $x_1$ 与 $x_2$
的标准开集 $D_{+}(F)$。
\end{enumerate}
\end{exercise}

\begin{exercise}
\label{exercises-exercise-compose-affine}
为什么仿射态射的复合仍是仿射态射？
\end{exercise}

\begin{exercise}[例子]
\label{exercises-exercise-examples}
给出下列对象的例子：
\begin{enumerate}
\item 域 $k$ 上的可约射影概形；
\item 具有 100 个点的概形；
\item 概形的非仿射态射。
\end{enumerate}
\end{exercise}

\begin{exercise}
\label{exercises-exercise-chevalley-hilbert-nullstellensatz}
Chevalley 定理与 Hilbert 零点定理。
\begin{enumerate}
\item 设 $\mathfrak p \subset \mathbf{Z}[x_1, \ldots, x_n]$
为极大理想。Chevalley 定理对
$\mathfrak p \cap \mathbf{Z}$ 蕴含什么？
\item 进而，Hilbert 零点定理对 $\kappa(\mathfrak p)$ 蕴含什么？
\end{enumerate}
\end{exercise}

\begin{exercise}
\label{exercises-exercise-P0}
设 $A$ 为环。令 $S = A[X]$，把它看成分次 $A$-代数，
其中 $X$ 的次数为 $1$。证明
$\text{Proj}(S) \cong \Spec(A)$（作为 $A$ 上的概形）。
\end{exercise}

\begin{exercise}
\label{exercises-exercise-finite-is-projective}
设 $A \to B$ 为有限环映射。
证明 $\Spec(B)$ 是 $\Spec(A)$ 上的 H-射影概形。
\end{exercise}

\begin{exercise}
\label{exercises-exercise-not-geometrically-irreducible}
给出域 $k$ 上概形 $X$ 的例子，使 $X$ 不可约，
但对某个有限扩张 $k'/k$，基变换
$X_{k'} = X \times_{\Spec(k)} \Spec(k')$ 连通但可约。
\end{exercise}





\section{概形，2011 年春季期末考试}
\label{exercises-section-final-exam-spring-2011}

\noindent
这些是哥伦比亚大学 2011 年春季一门概形课程的期末考试题。

\begin{exercise}[定义]
\label{exercises-exercise-definitions-spring-2011}
给出斜体概念的定义。
\begin{enumerate}
\item {\it 分离}概形；
\item 概形的{\it 普遍闭}态射；
\item 对包含在同一个域中的局部环 $A, B$，所谓
{\it $A$ 控制 $B$}；
\item 概形 $X$ 的{\it 维数}；
\item 概形 $X$ 的不可约闭子概形 $Y$ 的{\it 余维数}。
\end{enumerate}
\end{exercise}

\begin{exercise}[结果]
\label{exercises-exercise-results-spring-2011}
陈述某个与课程中讨论的事实形式等价的命题。
\begin{enumerate}
\item 例如，簇的态射 $X \to Y$ 的真性赋值判别法。
\item 对域 $k$ 上的簇 $X$，$\dim(X)$ 与 $X$ 的函数域
$k(X)$ 之间的关系。
\item 填空：$k$ 上非奇异射影曲线与非常值态射的范畴，
反等价于 $\ldots\ldots\ldots$。
\item Noether 正规化。
\item Jacobian 判别法。
\end{enumerate}
\end{exercise}

\begin{exercise}
\label{exercises-exercise-genus-plane-curve}
设 $k$ 为域。设 $F \in k[X_0, X_1, X_2]$ 为次数 $d$ 的齐次形式。
假设 $C = V_{+}(F) \subset \mathbf{P}^2_k$ 是 $k$ 上的光滑曲线。
% P12-ERR-0493
记 $i : C \to \mathbf{P}^2_k$ 为相应的闭浸入。
\begin{enumerate}
\item 证明 $\mathbf{P}^2_k$ 上存在凝聚层的短正合列
$$
0 \to \mathcal{O}_{\mathbf{P}^2_k}(-d)
\to \mathcal{O}_{\mathbf{P}^2_k} \to i_*\mathcal{O}_C \to 0
$$
请说明映射是什么，并简要说明为何该列正合。
\item 推出 $H^0(C, \mathcal{O}_C) = k$。
\item 计算 $C$ 的亏格。
\item 现在假设 $P = (0 : 0 : 1)$ 不在 $C$ 上。
证明由 $(a_0 : a_1 : a_2) \mapsto (a_0 : a_1)$ 给出的
$\pi : C \to \mathbf{P}^1_k$ 的次数为 $d$。
\item 假设 $k$ 代数闭，所有分歧指数（“$e_i$”）均为
$1$ 或 $2$，且 $k$ 的特征不等于 $2$。
$\pi : C \to \mathbf{P}^1_k$ 有多少个分歧点？
\item 用 $F$ 表示，你认为 $\pi$ 的分歧点集合的一组方程是什么？
\item 你能猜出 $K_C$ 吗？
\end{enumerate}
\end{exercise}

\begin{exercise}
\label{exercises-exercise-Pic-triangle}
设 $k$ 为域。设 $X$ 为 $k$ 上的一个“三角形”，即取三份
$\mathbf{A}^1_k$，把第一份的 $0$ 与第二份的 $1$ 粘合，
第二份的 $0$ 与第三份的 $1$ 粘合，再把第三份的 $0$ 与第一份的
$1$ 粘合，从而得到 $X$。事实证明，
$X$ 同构于 $\Spec(k[x, y]/(xy(x + y + 1)))$；可以直接使用这一点。
计算 $X$ 的 Picard 群。
\end{exercise}

\begin{exercise}
\label{exercises-exercise-birational-morphism-curves-ample}
设 $k$ 为域。设 $\pi : X \to Y$ 为曲线的有限双有理态射，
其中 $X$ 是 $k$ 上的射影非奇异曲线。
由课程内容可知，$Y$ 是真曲线，$\pi$ 是 $Y$ 的正规化态射。
课程中还已看到，存在稠密开集 $V \subset Y$，使
$U = \pi^{-1}(V)$ 是 $X$ 中的稠密开集，且
$\pi : U \to V$ 是同构。
\begin{enumerate}
\item 证明存在有效 Cartier 除子 $D \subset X$，使
$D \subset U$，并且 $\mathcal{O}_X(D)$ 在 $X$ 上充足。
\item 设 $D$ 如（1）。证明 $E = \pi(D)$ 是 $Y$ 上的有效
Cartier 除子。
\item 简要说明为何：
\begin{enumerate}
\item 映射 $\mathcal{O}_Y \to \pi_*\mathcal{O}_X$ 有凝聚余核
$Q$，且它支撑在 $Y \setminus V$ 上；
\item 对每个 $n$，都有相应映射
$\mathcal{O}_Y(nE) \to \pi_*\mathcal{O}_X(nD)$，
其余核同构于 $Q$。
\end{enumerate}
\item 证明
$\dim_k H^0(X, \mathcal{O}_X(nD)) - \dim_k H^0(Y, \mathcal{O}_Y(nE))$
有界（被什么量所界？），并推出当 $n$ 很大时，可逆层
$\mathcal{O}_Y(nE)$ 有许多截面（为什么？）。
\end{enumerate}
\end{exercise}






\section{概形，2011 年秋季期末考试}
\label{exercises-section-final-exam-fall-2011}

\noindent
这些是哥伦比亚大学 2011 年秋季一门交换代数课程的期末考试题。


\begin{exercise}[定义]
\label{exercises-exercise-definitions-fall-2011}
给出斜体概念的定义。
\begin{enumerate}
\item {\it Noether} 环；
\item {\it Noether} 概形；
\item {\it 有限}环同态；
\item 概形的{\it 有限}态射；
\item 环的{\it 维数}。
\end{enumerate}
\end{exercise}

\begin{exercise}[结果]
\label{exercises-exercise-results-fall-2011}
陈述某个与课程中讨论的事实形式等价的命题。
\begin{enumerate}
\item Zariski 主定理。
\item Noether 正规化。
\item 中国剩余定理。
\item 有限环映射的上升定理。
\end{enumerate}
\end{exercise}

\begin{exercise}
\label{exercises-exercise-dimension-of-ring}
设 $(A, \mathfrak m, \kappa)$ 为 Noether 局部环，
其剩余域的特征不为 $2$。
假设 $\mathfrak m$ 由三个元素 $x, y, z$ 生成，
并且在 $A$ 中 $x^2 + y^2 + z^2 = 0$。
\begin{enumerate}
\item $\dim(A)$ 的可能取值是什么？
\item 给出例子，说明每个取值都可能出现。
\item 证明若 $\dim(A) = 2$，则 $A$ 是整环。
（提示：考察 $\bigoplus_{n \geq 0} \mathfrak m^n/\mathfrak m^{n + 1}$。）
\end{enumerate}
\end{exercise}

\begin{exercise}
\label{exercises-exercise-localization}
设 $A$ 为环，$S \subset T \subset A$ 为乘法子集。假设
$$
\{\mathfrak q \mid \mathfrak q \cap S = \emptyset\} =
\{\mathfrak q \mid \mathfrak q \cap T = \emptyset\}.
$$
证明 $S^{-1}A \to T^{-1}A$ 是同构。
\end{exercise}

\begin{exercise}
\label{exercises-exercise-locus-of-rank-1}
设 $k$ 为代数闭域。令
$$
V_0 = \{ A \in \text{Mat}(3 \times 3, k) \mid \text{rank}(A) = 1\}
\subset \text{Mat}(3 \times 3, k) = k^9.
$$
\begin{enumerate}
\item 证明 $V_0$ 是某个（Zariski）局部闭子集
$V \subset \mathbf{A}^9_k$ 的闭点集合。
\item $V$ 不可约吗？
\item $\dim(V)$ 是多少？
\end{enumerate}
\end{exercise}

\begin{exercise}
\label{exercises-exercise-not-complete-intersection}
证明 $\mathbf{C}[x, y]$ 中的理想 $(x^2, xy, y^2)$
不能由 $2$ 个元素生成。
\end{exercise}

\begin{exercise}
\label{exercises-exercise-finite-projection}
设 $f \in \mathbf{C}[x, y]$ 为非常值多项式。
证明对某些 $\alpha, \beta \in \mathbf{C}$，$\mathbf{C}$-代数映射
$$
\mathbf{C}[t] \longrightarrow \mathbf{C}[x, y]/(f),\quad
t \longmapsto \alpha x + \beta y
$$
是有限的。
\end{exercise}

\begin{exercise}
\label{exercises-exercise-union-of-two-affines}
证明：给定有限多个点
$p_1, \ldots, p_n \in \mathbf{C}^2$，概形
$\mathbf{A}^2_\mathbf{C} \setminus \{p_1, \ldots, p_n\}$
是两个仿射开集的并。
\end{exercise}

\begin{exercise}
\label{exercises-exercise-surjection-A1-P1}
证明存在概形的满态射
$\mathbf{A}^1_\mathbf{C} \to \mathbf{P}^1_\mathbf{C}$。
（“满”仅表示在底层点集上为满射。）
\end{exercise}

\begin{exercise}
\label{exercises-exercise-almost-surjective}
设 $k$ 为代数闭域，$A \subset B$ 为整环扩张，并且二者都是
有限型 $k$-代数。用下列步骤证明
$\Spec(B) \to \Spec(A)$ 的像包含 $\Spec(A)$ 的一个非空开子集：
\begin{enumerate}
\item 当 $A \to B$ 还是有限映射时证明该结论。
\item 当 $B$ 的分式域是 $A$ 的分式域的有限扩张时证明该结论。
\item 把命题归约到前一种情形。
\end{enumerate}
\end{exercise}






\section{概形，2013 年秋季期末考试}
\label{exercises-section-final-exam-fall-2013}

\noindent
这些是哥伦比亚大学 2013 年秋季一门交换代数课程的期末考试题。

\begin{exercise}[定义]
\label{exercises-exercise-definitions-fall-2013}
给出斜体概念的定义。
\begin{enumerate}
\item 环的{\it 根理想}；
\item {\it 有限型}环同态；
\item {\it Weil 式微分}；
\item {\it 概形}。
\end{enumerate}
\end{exercise}

\begin{exercise}[结果]
\label{exercises-exercise-results-fall-2013}
陈述某个与课程中讨论的事实形式等价的命题。
\begin{enumerate}
\item 关于分次模的 Hilbert 多项式的结果。
\item Noether 局部环 $(R, \mathfrak m)$ 的维数与
$\bigoplus_{n \geq 0} \mathfrak m^n/\mathfrak m^{n + 1}$。
\item Riemann--Roch 定理。
\item Clifford 定理。
\item Chevalley 定理。
\end{enumerate}
\end{exercise}

\begin{exercise}
\label{exercises-exercise-surjective-after-localization}
设 $A \to B$ 为环映射，$S \subset A$ 为乘法子集。
假设 $A \to B$ 有限型，并且
$S^{-1}A \to S^{-1}B$ 为满射。证明存在 $f \in S$，
使 $A_f \to B_f$ 为满射。
\end{exercise}

\begin{exercise}
\label{exercises-exercise-injective-local-ring-map-not-surjective}
给出局部环的单射局部同态 $A \to B$ 的例子，使得
$\Spec(B) \to \Spec(A)$ 不满。
\end{exercise}

\begin{situation}[平面曲线的记号]
\label{exercises-situation-curve-in-the-plane}
设 $k$ 为代数闭域。设
$F(X_0, X_1, X_2) \in k[X_0, X_1, X_2]$ 为次数 $d$ 的
齐次不可约多项式。令
$$
D = V(F) \subset \mathbf{P}^2
$$
为 $F$ 的零点所给出的射影平面曲线。

令 $x = X_1/X_0$、$y = X_2/X_0$，并令
$f(x, y) = X_0^{-d}F(X_0, X_1, X_2) = F(1, x, y)$。

% P12-ERR-0494
用 $K$ 表示整环 $k[x, y]/(f)$ 的分式域。
令 $C$ 为与 $K$ 对应的抽象曲线。
回顾（来自讲课）：存在满映射 $C \to D$，它在 $D$ 的非奇异轨迹上
为双射，并且当 $D$ 非奇异时为同构。
令 $f_x = \partial f/\partial x$、$f_y = \partial f/\partial y$。
最后，用 $\omega = \text{d}x/f_y = - \text{d}y/f_x$ 表示讲课中讨论的
$\Omega_{K/k}$ 的元素。用 $K_C$ 表示 $\omega$ 的零点与极点除子。
\end{situation}

\begin{exercise}
\label{exercises-exercise-node-in-the-plane}
在情形 \ref{exercises-situation-curve-in-the-plane} 中，假设 $d \geq 3$，
且曲线 $D$ 恰有一个奇点，即 $P = (1 : 0 : 0)$。
进一步假设在 $P = (0, 0)$ 附近有展开式
$$
f(x, y) = xy + h.o.t
$$
则 $C$ 在 $P$ 上方有两个点 $v$ 与 $w$，由下式刻画：
$$
v(x) = 1, v(y) > 1
\quad\text{且}\quad
w(x) > 1, w(y) = 1
$$
\begin{enumerate}
\item 证明 $\Omega_{K/k}$ 的元素
$\omega = \text{d}x/f_y = - \text{d}y/f_x$
在 $v$ 与 $w$ 处都有一阶极点。
（$\omega$ 在非奇异点处的行为如讲课中所述。）
\item 讲课中已证明，$\omega$ 在通过映射 $C \to D$ 拉回到
$C$ 的除子 $X_0 = 0$ 处消失至 $d - 3$ 阶。
结合（1）的信息，$C$ 上 $\omega$ 的零点与极点除子的次数是多少？
\item 曲线 $C$ 的亏格是多少？
\end{enumerate}
\end{exercise}

\begin{exercise}
\label{exercises-exercise-smooth-plane-curve-linear-system}
在情形 \ref{exercises-situation-curve-in-the-plane} 中，假设 $d = 5$，
且曲线 $C = D$ 非奇异。讲课中已证明 $C$ 的亏格为 $6$，
而线性系 $K_C$ 由
$$
L(K_C) = \{h\omega \mid h \in k[x, y],\ \deg(h) \leq 2\}
$$
给出，其中 $\deg$ 表示总次数\footnote{这里得到 $\leq 2$，
因为 $d - 3 = 5 - 3 = 2$。}。
设 $P_1, P_2, P_3, P_4, P_5 \in D$ 为位于仿射开集
$X_0 \not = 0$ 中的两两不同的点。
用 $\sum P_i = P_1 + P_2 + P_3 + P_4 + P_5$ 表示 $C$ 的相应除子。
\begin{enumerate}
\item 用多项式描述 $L(K_C - \sum P_i)$。
\item $l(\sum P_i)$ 有哪些可能值？
\end{enumerate}
\end{exercise}

\begin{exercise}
\label{exercises-exercise-rational-curve-high-degree}
写出情形 \ref{exercises-situation-curve-in-the-plane} 中的一个 $F$，
使 $d = 100$ 且 $C$ 的亏格为 $0$。
\end{exercise}

\begin{exercise}
\label{exercises-exercise-high-degree-curve-quadratic-equation}
设 $k$ 为代数闭域。设 $K/k$ 为超越次数 $1$ 的有限生成域扩张。
% P12-ERR-0495
令 $C$ 为与 $K$ 对应的抽象曲线。设 $V \subset K$ 为一个
$g^r_d$，$\Phi : C \to \mathbf{P}^r$ 为相应态射。
证明：若 $V$ 是完备线性系，且 $d$ 相对于 $C$ 的亏格足够大，
则 $C$ 的像包含在某个二次超曲面中\footnote{二次超曲面是次数为
$2$ 的超曲面，即 $\mathbf{P}^r$ 中某个次数为 $2$ 的齐次多项式的
零点集。}。（加分题：给出所需次数的良好界。）
\end{exercise}

\begin{exercise}
\label{exercises-exercise-surjective-map-a2-p2}
记号如情形 \ref{exercises-situation-curve-in-the-plane}。
设 $U \subset \mathbf{P}^2_k$ 为补集是 $D$ 的开子概形。
描述 $k$-代数 $A = \mathcal{O}_{\mathbf{P}^2_k}(U)$。
给出 $A$ 作为 $k$-代数的生成元个数的上界。
\end{exercise}




\section{概形，2014 年春季期末考试}
\label{exercises-section-final-exam-spring-2014}

\noindent
这些是哥伦比亚大学 2014 年春季一门概形课程的期末考试题。

\begin{exercise}[定义]
\label{exercises-exercise-definitions-spring-2014}
设 $(X, \mathcal{O}_X)$ 为概形。给出斜体概念的定义。
\begin{enumerate}
\item $X$ 在点 $x$ 处的{\it 局部环}；
\item $\mathcal{O}_X$-模的{\it 拟凝聚}层；
% P12-ERR-0496
\item $\mathcal{O}_X$-模的{\it 凝聚}层（请假设 $X$ 局部 Noether；
\item $X$ 的{\it 仿射开集}；
\item {\it 概形的有限态射} $X \to Y$。
\end{enumerate}
\end{exercise}

\begin{exercise}[定理]
\label{exercises-exercise-results-spring-2014}
对每一项，精确陈述讲课中讨论的一个相关非平凡事实。
\begin{enumerate}
\item 关于簇的多重亏格的双有理不变性；
\item 仿射态射是局部性质；
\item 态射的概形论纤维的拓扑；
\item 真性赋值判别法。
\end{enumerate}
\end{exercise}

\begin{exercise}
\label{exercises-exercise-miss-curve}
设 $X = \mathbf{A}^2_\mathbf{C}$，其中 $\mathbf{C}$ 为复数域。
所谓{\it 直线}，是指由一个线性方程 $ax + by + c = 0$
定义的 $X$ 的闭子概形，其中
$a, b, c \in \mathbf{C}$ 且 $(a, b) \not = (0, 0)$。
所谓{\it 曲线}，是指维数为 $1$ 的不可约（因而非空）闭子概形
$C \subset X$。所谓{\it 二次曲线}，是指由一个二次方程
$ax^2 + bxy + cy^2 + dx + ey + f = 0$ 定义的曲线，其中
$a, b, c, d, e, f \in \mathbf{C}$ 且
$(a, b, c) \not = (0, 0, 0)$。
\begin{enumerate}
\item 找一条曲线 $C$，使每条直线都与 $C$ 有非空交。
\item 找一条曲线 $C$，使每条直线与每条二次曲线都与 $C$ 有非空交。
\item 证明对每条曲线 $C$，都存在另一条曲线，使
$C \cap C' = \emptyset$。
\end{enumerate}
\end{exercise}

\begin{exercise}
\label{exercises-exercise-normal-bundle-exceptional-curve}
设 $k$ 为域。设 $b : X \to \mathbf{A}^2_k$ 为仿射平面在原点处的爆破。
换言之，若 $\mathbf{A}^2_k = \Spec(k[x, y])$，则
$X = \text{Proj}(\bigoplus_{n \geq 0} \mathfrak m^n)$，
其中 $\mathfrak m = (x, y) \subset k[x, y]$。
证明下列命题：
\begin{enumerate}
\item $b$ 在原点上的概形论纤维 $E$ 同构于 $\mathbf{P}^1_k$；
\item $E$ 是 $X$ 上的有效 Cartier 除子；
\item $\mathcal{O}_X(-E)$ 限制到 $E$ 后是次数为 $1$ 的线丛。
\end{enumerate}
（回顾：$\mathcal{O}_X(-E)$ 是 $E$ 在 $X$ 中的理想层。）
\end{exercise}

\begin{exercise}
\label{exercises-exercise-surjective-map-affine-variety-projective-variety}
设 $k$ 为域，$X$ 为 $k$ 上的射影簇。
证明存在 $k$ 上的仿射簇 $U$ 以及簇的满态射 $U \to X$。
\end{exercise}

\begin{exercise}
\label{exercises-exercise-vandermonde}
设 $k$ 为特征 $p > 0$ 且不同于 $2,3$ 的域。
考虑 $\mathbf{P}^n_k$ 中由下列方程定义的闭子概形 $X$：
$$
\sum\nolimits_{i = 0, \ldots, n} X_i = 0,\quad
\sum\nolimits_{i = 0, \ldots, n} X_i^2 = 0,\quad
\sum\nolimits_{i = 0, \ldots, n} X_i^3 = 0
$$
% P12-ERR-0497
对哪些数对 $(n, p)$，这个簇是奇异的？
\end{exercise}







\section{交换代数，2016 年秋季期末考试}
\label{exercises-section-final-exam-fall-2016}

\noindent
这些是哥伦比亚大学 2016 年秋季一门交换代数课程的期末考试题。

\begin{exercise}[定义]
\label{exercises-exercise-definitions-fall-2016}
设 $R$ 为环。给出斜体概念的定义。
\begin{enumerate}
\item $R$ 在素理想 $\mathfrak p$ 处的{\it 局部环}；
\item {\it 有限} $R$-模；
\item {\it 有限表现} $R$-模；
\item $R$ 是{\it 正则的}；
\item $R$ 是{\it 链状的}；
\item $R$ 是{\it Cohen--Macaulay 的}。
\end{enumerate}
\end{exercise}

\begin{exercise}[定理]
\label{exercises-exercise-results-fall-2016}
对每一项，精确陈述讲课中讨论的一个相关非平凡事实。
\begin{enumerate}
\item 正则环；
\item Cohen--Macaulay 模的相伴素理想；
\item 域上有限型整环的维数；
\item Chevalley 定理。
\end{enumerate}
\end{exercise}

\begin{exercise}
\label{exercises-exercise-finite-injective}
设 $A \to B$ 为满足下列条件的环映射：
\begin{enumerate}
\item $A$ 是极大理想为 $\mathfrak m$ 的局部环；
\item $A \to B$ 是有限环映射\footnote{回顾：这表示 $B$
作为 $A$-模是有限的。}；
\item $A \to B$ 是单射（把 $A$ 看作 $B$ 的子环）。
\end{enumerate}
证明存在素理想 $\mathfrak q \subset B$，使
$\mathfrak m = A \cap \mathfrak q$。
\end{exercise}

\begin{exercise}
\label{exercises-exercise-find-components}
设 $k$ 为域，$R = k[x, y, z, w]$。考虑理想
$I = (xy, xz, xw)$。$V(I) \subset \Spec(R)$ 的不可约分支是什么，
其维数分别是多少？
\end{exercise}

\begin{exercise}
\label{exercises-exercise-no-nonconstant-morphism}
设 $k$ 为域，$A = k[x, x^{-1}]$、$B = k[y]$。
证明任意 $k$-代数映射 $\varphi : A \to B$ 都把 $x$ 映到常数。
\end{exercise}

\begin{exercise}
\label{exercises-exercise-regular-over-Z}
考虑环 $R = \mathbf{Z}[x, y]/(xy - 7)$。
证明 $R$ 正则。
\end{exercise}

\noindent
给定 Noether 局部环 $(R, \mathfrak m, \kappa)$，对
$n \geq 0$ 令
$\varphi_R(n) = \dim_\kappa(\mathfrak m^n/\mathfrak m^{n + 1})$。

\begin{exercise}
\label{exercises-exercise-hilbert-function-allowed}
是否存在 Noether 局部环 $R$，使对所有 $n \geq 0$ 都有
$\varphi_R(n) = n + 1$？
\end{exercise}

\begin{exercise}
\label{exercises-exercise-hilbert-function-prohibited}
设 $R$ 为 Noether 局部环。假设
$\varphi_R(0) = 1$、$\varphi_R(1) = 3$、
$\varphi_R(2) = 5$。证明 $\varphi_R(3) \leq 7$。
\end{exercise}





\section{概形，2017 年春季期末考试}
\label{exercises-section-final-exam-spring-2017}

\noindent
这些是哥伦比亚大学 2017 年春季一门概形课程的期末考试题。

\begin{exercise}[定义]
\label{exercises-exercise-definitions-spring-2017}
设 $f : X \to Y$ 为概形态射。简要给出斜体概念的定义。
\begin{enumerate}
\item $f$ 在 $y \in Y$ 处的{\it 概形论纤维}；
\item $f$ 是{\it 有限态射}；
\item {\it 拟凝聚} $\mathcal{O}_X$-模；
% P12-ERR-0498
\item $X$ 是{\it 簇}；
\item $f$ 是{\it 光滑态射}；
\item $f$ 是{\it 真态射}。
\end{enumerate}
\end{exercise}

\begin{exercise}[定理]
\label{exercises-exercise-results-spring-2017}
对每一项，精确而简要地陈述讲课中讨论的一个相关非平凡事实。
\begin{enumerate}
\item 拟凝聚层的前推；
\item 射影簇上凝聚层的上同调；
\item 域上射影概形的 Serre 对偶；
\item Riemann--Hurwitz 公式。
\end{enumerate}
\end{exercise}

\begin{exercise}
\label{exercises-exercise-compute-degree}
设 $k$ 为代数闭域。设 $\ell > 100$ 为不同于 $k$ 的特征的素数。
设 $X$ 为仿射曲线
$$
y^\ell = x(x - 1)^3
$$
在 $\mathbf{A}^2_k$ 中所给出的非奇异射影模型。回答下列问题：
\begin{enumerate}
\item $X$ 的亏格是多少？
\item 给出 $X$ 的脉性\footnote{脉性是从 $X$ 到
$\mathbf{P}^1_k$ 的非常值态射的最小次数。}的一个上界。
\end{enumerate}
\end{exercise}

\begin{exercise}
\label{exercises-exercise-count-components}
设 $k$ 为代数闭域。设 $X$ 为 $k$ 上既约射影概形，
其所有不可约分支的维数都等于 $1$。
设 $\omega_{X/k}$ 为相对对偶模。
证明若 $\dim_k H^1(X, \omega_{X/k}) > 1$，则 $X$ 不连通。
\end{exercise}

\begin{exercise}
\label{exercises-exercise-sections-L-L-inverse}
给出概形 $X$ 与非平凡可逆 $\mathcal{O}_X$-模
$\mathcal{L}$ 的例子，使 $H^0(X, \mathcal{L})$ 与
$H^0(X, \mathcal{L}^{\otimes -1})$ 均非零。
\end{exercise}

\begin{exercise}
\label{exercises-exercise-in-product}
设 $k$ 为代数闭域，$g \geq 3$。
设 $X$ 与 $X'$ 为 $k$ 上亏格分别为 $g$ 与 $g + 1$ 的
光滑射影曲线。设 $Y \subset X \times X'$ 为曲线，并且投影
$Y \to X$ 与 $Y \to X'$ 都不是常值的。
证明 $Y$ 的非奇异射影模型的亏格 $\geq 2g + 1$。
\end{exercise}

\begin{exercise}
\label{exercises-exercise-finitely-many}
设 $k$ 为有限域，$g > 1$。勾勒下列命题的证明：
$k$ 上亏格为 $g$ 的光滑射影曲线仅有有限多个同构类。
（即使只是尝试作答也可得分。）
\end{exercise}








\section{交换代数，2017 年秋季期末考试}
\label{exercises-section-final-exam-fall-2017}

\noindent
这些是哥伦比亚大学 2017 年秋季一门交换代数课程的期末考试题。

\begin{exercise}[定义]
\label{exercises-exercise-definitions-fall-2017}
简要给出斜体概念的定义。
\begin{enumerate}
\item 函子 $F : \mathcal{A} \to \mathcal{B}$ 的{\it 左伴随}；
\item 域扩张 $L/K$ 的{\it 超越次数}；
\item 经典仿射簇 $X \subset k^n$ 上的{\it 正则函数}；
\item 拓扑空间上的{\it 层}；
\item {\it 局部环}；
\item 概形态射 $f : X \to Y$ 是{\it 仿射的}。
\end{enumerate}
\end{exercise}

\begin{exercise}[定理]
\label{exercises-exercise-results-fall-2017}
对每一项，精确而简要地陈述讲课中讨论的一个相关非平凡事实
（若不止一个，只需任选其一）。
\begin{enumerate}
\item Yoneda 引理；
\item Mayer--Vietoris；
\item 维数与上同调；
\item Hilbert 多项式；
\item 射影空间的对偶性。
\end{enumerate}
\end{exercise}

\begin{exercise}
\label{exercises-exercise-compute-dimension}
设 $k$ 为代数闭域。考虑带 Zariski 拓扑与坐标
$x_1, x_2, x_3, x_4, x_5$ 的 $k^5$ 中由方程
$$
x_1^2 - x_4 = 0,\quad
x_2^5 - x_5 = 0,\quad
x_3^2 + x_3 + x_4 + x_5 = 0
$$
给出的闭子集 $X$。$X$ 的维数是多少？为什么？
\end{exercise}

\begin{exercise}
\label{exercises-exercise-can-there-be}
设 $k$ 为域，$X = \mathbf{P}^1_k$ 为 $k$ 上的 $1$-维射影空间。
设 $\mathcal{E}$ 为有限局部自由 $\mathcal{O}_X$-模。
对 $d \in \mathbf{Z}$，用
$\mathcal{E}(d) = \mathcal{E} \otimes_{\mathcal{O}_X} \mathcal{O}_X(d)$
表示 $\mathcal{E}$ 的第 $d$ 次 Serre 扭曲，并令
$h^i(X, \mathcal{E}(d)) = \dim_k H^i(X, \mathcal{E}(d))$。
\begin{enumerate}
\item 为什么不存在满足
$h^0(X, \mathcal{E}) = 5$ 且 $h^0(X, \mathcal{E}(1)) = 4$
的 $\mathcal{E}$？
\item 为什么不存在满足
$h^1(X, \mathcal{E}(1)) = 5$ 且 $h^1(X, \mathcal{E}) = 4$
的 $\mathcal{E}$？
\item 对哪些 $a \in \mathbf{Z}$，可能存在 $X$ 上的向量丛
$\mathcal{E}$，使
$$
\begin{matrix}
h^0(X, \mathcal{E})\phantom{(1)} = 1 &
h^1(X, \mathcal{E})\phantom{(1)} = 1 \\
h^0(X, \mathcal{E}(1)) = 2 &
h^1(X, \mathcal{E}(1)) = 0 \\
h^0(X, \mathcal{E}(2)) = 4 &
h^1(X, \mathcal{E}(2)) = a
\end{matrix}
$$
\end{enumerate}
欢迎并鼓励不完整的答案。
\end{exercise}

\begin{exercise}
\label{exercises-exercise-banana}
设拓扑空间 $X$ 是两个闭子集 $Y$ 与 $Z$ 的并
$X = Y \cup Z$，其交记为 $W = Y \cap Z$。
用 $i : Y \to X$、$j : Z \to X$ 与 $k : W \to X$
表示包含映射。
\begin{enumerate}
\item 证明存在层的短正合列
$$
0 \to \underline{\mathbf{Z}}_X \to
i_*(\underline{\mathbf{Z}}_Y) \oplus
j_*(\underline{\mathbf{Z}}_Z) \to
k_*(\underline{\mathbf{Z}}_W) \to 0
$$
其中 $\underline{\mathbf{Z}}_X$ 表示 $X$ 上取值为
$\mathbf{Z}$ 的常层，等等。
\item 关于 $\mathbf{Z}$-系数的 $X$、$Y$、$Z$、$W$
上同调群之间的关系，能得出什么结论？
\end{enumerate}
\end{exercise}

\begin{exercise}
\label{exercises-exercise-cohomology-infinite-punctured}
设 $k$ 为域。令 $A = k[x_1, x_2, x_3, \ldots]$ 为无穷多个变量的
多项式环。用 $\mathfrak m$ 表示由所有变量生成的 $A$ 的极大理想。
令 $X = \Spec(A)$，$U = X \setminus \{\mathfrak m\}$。
\begin{enumerate}
\item 证明 $H^1(U, \mathcal{O}_U) = 0$。
提示：作 {\v C}ech 上同调计算。
\item 对 $i \geq 1$，你猜测 $H^i(U, \mathcal{O}_U)$ 是什么？
\end{enumerate}
\end{exercise}

\begin{exercise}
\label{exercises-exercise-principal}
设 $A$ 为局部环，$a \in A$ 为非零因子。
设 $I, J \subset A$ 为满足 $IJ = (a)$ 的理想。
证明理想 $I$ 是主理想，即由一个元素生成
（该元素也将是非零因子）。
\end{exercise}




\section{概形，2018 年春季期末考试}
\label{exercises-section-final-exam-spring-2018}

\noindent
这些是哥伦比亚大学 2018 年春季一门概形课程的期末考试题。

\begin{exercise}[定义]
\label{exercises-exercise-definitions-spring-2018}
简要给出斜体概念的定义。
设 $k$ 为代数闭域，$X$ 为 $k$ 上的射影曲线。
\begin{enumerate}
\item $k$ 上的{\it 光滑}代数；
\item $X$ 上可逆 $\mathcal{O}_X$-模的{\it 次数}；
\item $X$ 的{\it 亏格}；
\item $X$ 的{\it Weil 除子类群}；
\item $X$ 是{\it 超椭圆的}；
\item $k$ 上光滑射影曲面中两条曲线的{\it 交数}。
\end{enumerate}
\end{exercise}

\begin{exercise}[定理]
\label{exercises-exercise-results-spring-2018}
对每一项，精确而简要地陈述讲课中讨论的一个相关非平凡事实
（若不止一个，只需任选其一）。
\begin{enumerate}
\item Riemann--Hurwitz 定理；
\item Clifford 定理；
\item 光滑射影曲面之间映射的分解；
\item Hodge 指标定理；
\item 有限域上曲线的 Riemann 猜想。
\end{enumerate}
\end{exercise}

\begin{exercise}
\label{exercises-exercise-hyperelliptic}
设 $k$ 为代数闭域。设 $X \subset \mathbf{P}^3_k$ 为次数 $d$、
亏格 $\geq 2$ 的光滑曲线。假设 $X$ 不包含在任何平面中，
并且 $\mathbf{P}^3_k$ 中有一条直线 $\ell$ 与 $X$ 交于
$d - 2$ 个点。证明 $X$ 是超椭圆曲线。
\end{exercise}

\begin{exercise}
\label{exercises-exercise-singular}
设 $k$ 为代数闭域。设 $X$ 为射影曲线，其两两不同的奇点为
$p_1, \ldots, p_n$。说明为何 $X$ 的正规化的亏格至多为
$-n + \dim_k H^1(X, \mathcal{O}_X)$。
\end{exercise}

\begin{exercise}
\label{exercises-exercise-how-many-blowups}
设 $k$ 为域，$X = \Spec(k[x, y])$ 为仿射 $2$-空间。令
$$
I = (x^3, x^2y, xy^2, y^3) \subset k[x, y].
$$
令 $Y \subset X$ 为与 $I$ 对应的闭子概形。
设 $b : X' \to X$ 为理想 $(x, y)$ 的爆破，即仿射空间在原点处的
爆破。
\begin{enumerate}
\item 证明概形论逆像 $b^{-1}Y \subset X'$ 是有效 Cartier 除子。
% P12-ERR-0499
\item 给出理想 $J \subset k[x, y]$ 的例子，使
$I \subset J \subset (x, y)$，并且若 $Z \subset X$ 是与 $J$
对应的闭子概形，则概形论逆像 $b^{-1}Z$ 不是有效 Cartier 除子。
\end{enumerate}
\end{exercise}

\begin{exercise}
\label{exercises-exercise-rational-surface}
设 $k$ 为代数闭域。考虑下列类型的曲面：
\begin{enumerate}
\item $S = C_1 \times C_2$，其中 $C_1$ 与 $C_2$ 是光滑射影曲线；
\item $S = C_1 \times C_2$，其中 $C_1$ 与 $C_2$ 是光滑射影曲线，
且 $C_1$ 的亏格 $> 0$；
\item $S \subset \mathbf{P}^3_k$ 是次数 $4$ 的超曲面；
\item $S \subset \mathbf{P}^3_k$ 是次数 $4$ 的光滑超曲面。
\end{enumerate}
对每种类型，简要说明为何该类型的曲面类包含或不包含有理曲面。
\end{exercise}

\begin{exercise}
\label{exercises-exercise-self-square-line}
设 $k$ 为代数闭域。设 $S \subset \mathbf{P}^3_k$ 为次数 $d$
的光滑超曲面。假设 $S$ 包含直线 $\ell$。
% P12-ERR-0500
把 $\ell$ 看作 $S$ 上的除子时，它的自交数是多少？
\end{exercise}

\section{交换代数，2019 年秋季期末考试}
\label{exercises-section-final-exam-fall-2019}

\noindent
以下是哥伦比亚大学 2019 年秋季交换代数课程期末考试的试题。

\begin{exercise}[定义]
\label{exercises-exercise-definitions-fall-2019}
简要定义下列斜体概念。
\begin{enumerate}
\item Noether 拓扑空间的一个 {\it 可构造子集}，
\item $R$-模 $M$ 在素理想 $\mathfrak p$ 处的 {\it 局部化}，
\item Noether 局部环 $(A, \mathfrak m, \kappa)$ 上模的 {\it 长度}，
\item 环 $R$ 上的 {\it 射影模}，以及
\item Noether 局部环 $(A, \mathfrak m, \kappa)$ 上的
{\it Cohen--Macaulay 模}。
\end{enumerate}
\end{exercise}

\begin{exercise}[定理]
\label{exercises-exercise-results-fall-2019}
对下列每一项，准确而简要地陈述一个课堂上讨论过的相关非平凡事实
（若有多个，只需任选一个）。
\begin{enumerate}
\item 可构造集的像，
\item Hilbert 零点定理，
\item 域上有限型代数的维数，
\item Noether 正规化，以及
\item 正则局部环。
\end{enumerate}
\end{exercise}

\noindent
回顾：对于环 $R$ 与理想 $I \subset R$，$V(I)$ 表示点
$\mathfrak p \in \Spec(R)$ 所成的集合，其中 $I \subset \mathfrak p$。

\begin{exercise}[构造素理想]
\label{exercises-exercise-infinitely-many-primes}
构造无穷多个互不相同的素理想
$\mathfrak p \subset \mathbf{C}[x, y]$，
使 $V(\mathfrak p)$ 包含 $(x, y)$ 与 $(x - 1, y - 1)$。
\end{exercise}

\begin{exercise}[不存在这样的素理想]
\label{exercises-exercise-no-prime}
令 $R = \mathbf{C}[x, y, z]/(xy)$。简要论证：不存在素理想
$\mathfrak p \subset R$，使得
$V(\mathfrak p)$ 同时包含 $(x, y - 1, z - 5)$
与 $(x - 1, y, z - 7)$。
\end{exercise}

\begin{exercise}[Frobenius]
\label{exercises-exercise-frobenius}
令 $p$ 为素数（为简化公式，你可以假设 $p = 2$）。
令 $R$ 为满足 $R$ 中 $p = 0$ 的环。
\begin{enumerate}
\item 证明映射 $F : R \to R$，$x \mapsto x^p$，是环同态。
\item 证明 $\Spec(F) : \Spec(R) \to \Spec(R)$ 是恒等映射。
\end{enumerate}
\end{exercise}

\noindent
回顾：拓扑空间中点的一个 {\it 特化} $x \leadsto y$，仅仅意指
$y$ 位于 $x$ 的闭包中。若不存在异于 $x$ 与 $y$ 的点 $z$，
使 $x \leadsto z$ 且 $z \leadsto y$，则称 $x \leadsto y$
为 {\it 直接特化}。

\begin{exercise}[维数]
\label{exercises-exercise-dimension}
设索伯拓扑空间 $X$ 含有 $5$ 个互不相同的点
$x, y, z, u, v$，它们之间有如下特化关系：
$$
\xymatrix{
x \ar[d] \ar[r] & u & v \ar[l] \ar[dl] \\
y \ar[r] & z
}
$$
这样的 $X$ 至少可以是多少维？若 $X$ 是域上有限型代数的谱，
且 $x \leadsto u$ 是直接特化，那么关于特化
$v \leadsto z$ 你能得出什么结论？
\end{exercise}

\begin{exercise}[Tor 计算]
\label{exercises-exercise-tor-computation}
令 $R = \mathbf{C}[x, y, z]$。令 $M = R/(x, z)$ 且
$N = R/(y, z)$。对于哪些 $i \in \mathbf{Z}$，
$\text{Tor}_i^R(M, N)$ 非零？
\end{exercise}

\begin{exercise}
\label{exercises-exercise-depth-goes-up}
令 $A \to B$ 为 Noether 局部环之间的平坦局部同态。
证明：若 $A$ 的深度为 $k$，则 $B$ 的深度至少为 $k$。
\end{exercise}

\section{代数几何，2020 年春季期末考试}
\label{exercises-section-final-exam-spring-2020}

\noindent
以下是哥伦比亚大学 2020 年春季代数几何课程期末考试的试题。

\begin{exercise}[定义]
\label{exercises-exercise-definitions-spring-2020}
简要定义下列斜体概念。
\begin{enumerate}
\item 一个 {\it 概形}，
\item 概形的一个 {\it 态射}，
\item 概形上的一个 {\it 拟凝聚模}，
\item 域 $k$ 上的一个 {\it 代数簇}，
\item 域 $k$ 上的一条 {\it 曲线}，
\item 概形的一个 {\it 有限态射}，
\item 拓扑空间上阿贝尔群层 $\mathcal{F}$ 的 {\it 上同调}，其所在空间为 $X$，
\item 概形 $X$（维数为 $d$，在域 $k$ 上适当）上的 {\it 对偶层}，以及
\item 从代数簇 $X$ 到代数簇 $Y$ 的一个 {\it 有理映射}。
\end{enumerate}
\end{exercise}

\begin{exercise}[定理]
\label{exercises-exercise-results-spring-2020}
对下列每一项，准确而简要地陈述一个课堂上讨论过的相关非平凡事实
（若有多个，只需任选一个）。
\begin{enumerate}
\item Noether 拓扑空间 $X$（维数为 $d$）上阿贝尔层的上同调，
\item 域 $k$ 上光滑代数簇的微分层 $\Omega^1_{X/k}$，
\item 光滑射影代数簇的对偶层 $\omega_X$（代数簇为 $X$，定义在域 $k$ 上），
\item 代数闭域 $k$ 上一条光滑适当的亏格 $0$ 曲线，以及
\item 次数为 $d$ 的平面曲线的亏格。
\end{enumerate}
\end{exercise}

\begin{exercise}
\label{exercises-exercise-coh-P1-infty-doubled}
设 $k$ 为域，$X$ 为 $k$ 上的概形。假设 $X = X_1 \cup X_2$
是一个开覆盖，其中 $X_1$、$X_2$ 都同构于 $\mathbf{P}^1_k$，
而 $X_1 \cap X_2$ 同构于 $\mathbf{A}^1_k$。（例如，将
$\mathbf{P}^1_k$ 的 $\infty$ 点复制一份就可得到这样的概形。）
证明 $\dim_k H^1(X, \mathcal{O}_X)$ 是无穷的。
\end{exercise}

\begin{exercise}
\label{exercises-exercise-no-morphism}
设 $k$ 为代数闭域。令 $Y$ 为亏格 $10$ 的光滑射影曲线。
为一条光滑射影曲线 $X$ 的亏格给出一个好的下界，使得存在一个
非恒定态射 $f : X \to Y$，且它不是同构。
\end{exercise}

\begin{exercise}
\label{exercises-exercise-element-order-n-in-pic}
设 $k$ 为特征 $0$ 的代数闭域。令
$$
X : T_0^d + T_1^d - T_2^d = 0 \subset \mathbf{P}^2_k
$$
为次数 $d \geq 3$ 的 Fermat 曲线。考虑 $X$ 上的闭点
$p = [1 : 0 : 1]$ 与 $q = [0 : 1 : 1]$。置 $D = [p] - [q]$。
\begin{enumerate}
\item 证明 $D$ 在 Weil 除子类群中非平凡。
\item 证明 $d D$ 在 Weil 除子类群中平凡。
（提示：尝试证明 $d[p]$ 与 $d[q]$ 都是平面中一条直线
与 $X$ 的交。）
\end{enumerate}
\end{exercise}

\begin{exercise}
\label{exercises-exercise-number-equations}
设 $k$ 为代数闭域。考虑 $2$-重嵌入
$$
\varphi : \mathbf{P}^2 \longrightarrow \mathbf{P}^5
$$
按照课堂材料/记号，这是态射
$$
\varphi = \varphi_{\mathcal{O}_{\mathbf{P}^2}(2)} :
\mathbf{P}^2
\longrightarrow
\mathbf{P}(\Gamma(\mathbf{P}^2, \mathcal{O}_{\mathbf{P}^2}(2)))
$$
用齐次坐标表示，它为
$$
[a_0 : a_1 : a_2] \longmapsto
[a_0^2 : a_0a_1 : a_0a_2 : a_1^2 : a_1a_2 : a_2^2]
$$
这是一个闭浸入（请直接使用这一点）。令
$I \subset k[T_0, \ldots, T_5]$ 为 $\varphi(\mathbf{P}^2)$ 的齐次理想；
也就是说，齐次部分 $I_d$ 的元素是次数为 $d$ 的齐次多项式
$F(T_0, \ldots, T_5)$，它们在闭子概形 $\varphi(\mathbf{P}^2)$
上的限制为零。计算 $\dim_k I_d$ 关于 $d$ 的函数。
\end{exercise}

\begin{exercise}
\label{exercises-exercise-dualizing-sheaf-rank-1}
设 $k$ 为代数闭域。令 $X$ 为维数为 $d$ 的 $k$-上适当概形，
其对偶模为 $\omega_X$。已知：
\begin{enumerate}
\item
$\text{Ext}^i_X(\mathcal{F}, \omega_X) \times
H^{d - i}(X, \mathcal{F}) \to H^d(X, \omega_X) \xrightarrow{t} k$
对所有 $i$ 以及所有相干 $\mathcal{O}_X$-模 $\mathcal{F}$，
均为非退化的，以及
\item $\omega_X$ 是某个秩 $r$ 的有限局部自由模。
\end{enumerate}
证明 $r = 1$。（提示：考察取一个支集在闭点上的适当模时会发生什么，
例如取 $\mathcal{F}$。）
\end{exercise}

\section{交换代数，2021 年秋季期末考试}
\label{exercises-section-final-exam-fall-2021}

\noindent
以下是哥伦比亚大学 2021 年秋季交换代数课程期末考试的试题。

\begin{exercise}[定义]
\label{exercises-exercise-definitions-fall-2021}
简要定义下列斜体概念。
\begin{enumerate}
\item 环 $A$ 的一个 {\it 乘法子集}，
\item 一个 {\it Artinian 环} $A$，
\item 环 $A$ 的 {\it 谱}作为拓扑空间，
\item 一个 {\it 平坦环映射} $A \to B$，
\item $A$ 中素理想 $\mathfrak p$ 的 {\it 高度}，以及
\item 环 $A$ 上的函子 {\it $\text{Tor}^A_i(-, -)$}。
\end{enumerate}
\end{exercise}

\begin{exercise}[定理]
\label{exercises-exercise-results-fall-2021}
对下列每一项，准确而简要地陈述一个课堂上讨论过的相关非平凡事实
（若有多个，只需任选一个）。
\begin{enumerate}
\item Artinian 环，
\item 平坦性与素理想，
\item 对于 Noether 局部环 $(A, \mathfrak m)$，$A/\mathfrak m^n$ 的长度，
\item 万有链式 Noether 环的维数公式，
\item Noether 局部环的完备化，以及
\item Artinian 局部环的 Matlis 对偶。
\end{enumerate}
\end{exercise}

\begin{exercise}[单位]
\label{exercises-exercise-simple-units}
$\mathbf{Z}[x, 1/x]$ 的单位群作为阿贝尔群具有什么结构？
无需解释。
\end{exercise}

\begin{exercise}[理想]
\label{exercises-exercise-ideals}
令 $A = \mathbf{F}_2[x, y]/(x^2, xy, y^2)$，并将
$\overline{x}$ 与 $\overline{y}$ 记为 $x$ 与 $y$ 在 $A$ 中的像。
列出 $A$ 的理想。无需解释。
\end{exercise}

\begin{exercise}[Tor 与 Ext]
\label{exercises-exercise-compute-tor-ext}
设 $(A, \mathfrak m, \kappa)$ 为 Noether 局部环。
置 $\varphi(n) = \dim_\kappa \mathfrak m^n/\mathfrak m^{n + 1}$。
\begin{enumerate}
\item 证明 $\text{Tor}_1^A(A/\mathfrak m^n, \kappa)$
作为 $\kappa$-向量空间的维数为 $\varphi(n)$。
\item 证明 $\text{Ext}^1_A(A/\mathfrak m^n, \kappa)$
作为 $\kappa$-向量空间的维数为 $\varphi(n)$。
\end{enumerate}
\end{exercise}

\begin{exercise}[两个向量]
\label{exercises-exercise-two-vectors}
% P12-ERR-0501
令 $A = \mathbf{Z}[a_1, a_2, a_3, b_1, b_2, b3]$。
令 $a = (a_1, a_2, a_3)$ 且 $b = (b_1, b_2, b_3)$ 为 $A^{\oplus 3}$ 中的向量。
考虑集合
$$
Z = \{\mathfrak p \in \Spec(A) \mid a, b
\text{ map to linearly dependent vectors of }
\kappa(\mathfrak p)^{\oplus 3}\}
$$
\begin{enumerate}
% P12-ERR-0502
\item 证明 $Z$ 是 $\Spec(A)$ 的闭子集。
\item $Z$ 的维数 $\dim(Z)$ 是多少？
\item 如果将 $\mathbf{Z}$ 替换为一个域，$\dim(Z)$ 会怎样变化？
\end{enumerate}
\end{exercise}

\begin{exercise}[内射模]
\label{exercises-exercise-injective-artinian-local}
设 $(A, \mathfrak m, \kappa)$ 为 Artinian 局部环。
% P12-ERR-0503
假设 $A$ 作为 $A$-模是内射的。证明
$\Hom_A(\kappa, A)$ 的维数为 $1$ 是一个
$\kappa$-向量空间。
\end{exercise}

\section{代数几何，2022 年春季期末考试}
\label{exercises-section-final-exam-spring-2022}

\noindent
以下是哥伦比亚大学 2022 年春季代数几何课程期末考试的试题。

\begin{exercise}[定义]
\label{exercises-exercise-definitions-spring-2022}
简要定义下列斜体概念。
\begin{enumerate}
\item 一个 {\it 概形}，
\item 概形 $X$ 上的一个 {\it 拟凝聚模}，
\item 概形的一个 {\it 平坦}态射 $X \to Y$，
\item 概形的一个 {\it 有限}态射 $X \to Y$，
\item 基概形 $S$ 上的一个 {\it 群概形} $G$，
\item 基概形 $S$ 上的一个 {\it 代数簇族}，
\item 域 $k$ 上代数簇 $X$ 的闭点 $x$ 的 {\it 次数}，
\item 点 $p = (a_0 : \ldots : a_n)$ 在 $\mathbf{P}^n(\mathbf{Q})$
中的通常 {\it 对数高度}，以及
\item 一个 {\it $C_i$ 域}。
\end{enumerate}
\end{exercise}

\begin{exercise}[定理]
\label{exercises-exercise-results-spring-2022}
对下列每一项，准确而简要地陈述一个课堂上讨论过的相关非平凡事实
（若有多个，只需任选一个）。
\begin{enumerate}
\item 从概形 $X$ 到仿射概形 $\Spec(A)$ 的态射，
\item 仿射概形 $X$ 上拟凝聚模 $\mathcal{F}$ 的上同调，
\item 域 $k$ 上 $\mathbf{P}^1_k$ 的 Picard 群，
\item 当 $S$ 为 Noether 概形时，平坦适当态射
$X \to S$ 的纤维维数，
\item 概形 $S$ 上的 $\mathbf{G}_m$-等变模，以及
\item 交的 Bézout 定理（若愿意，可限制在某个特殊情形）。
\end{enumerate}
\end{exercise}

\begin{exercise}[三次超曲面]
\label{exercises-exercise-cubics}
令 $F \in \mathbf{C}[T_0, \ldots, T_n]$ 为次数 $3$ 的齐次多项式。
给定 $3$ 个向量 $x, y, z \in \mathbf{C}^{n + 1}$，考虑条件
$$
(*)\quad
F(\lambda x + \mu y + \nu z) = 0 \text{ 于 } \mathbf{C}[\lambda, \mu, \nu]
$$
\begin{enumerate}
\item 所有 $x, y, z$ 的选择所成空间的维数是多少？
\item 条件 (*) 对 $x$、$y$、$z$ 的坐标给出了多少个方程？
\item 满足 (*) 的所有三元组 $x, y, z$ 所成空间的预期维数是多少？
\item $x, y, z$ 线性相关的所有三元组所成空间的维数是多少？
\item 推出：在 $\mathbf{P}^n$ 中次数为 $3$ 的超曲面上，
当 $n \geq a$（其中 $a$ 由你求出）时，我们预期能找到维数为 $2$
的线性子空间。
\end{enumerate}
\end{exercise}

\begin{exercise}[高度]
\label{exercises-exercise-heights}
设 $K$ 为域。令 $h_n : \mathbf{P}^n(K) \to \mathbf{R}$（$n \geq 0$）
为满足课堂中讨论的 $2$ 条公理的一族函数。令 $X$ 为
$K$ 上的射影代数簇。令 $\mathcal{L}$ 为可逆的
$\mathcal{O}_X$-模，并回顾课堂中构造的关联高度函数
$h_\mathcal{L} : X(K) \to \mathbf{R}$。令 $\alpha : X \to X$
为 $X$ 上定义在 $K$ 上的自同构。
\begin{enumerate}
\item 证明 $P \mapsto h_\mathcal{L}(\alpha(P))$
与函数 $h_{\alpha^*\mathcal{L}}$ 相差有界量。（提示：回顾若存在态射
$\varphi : X \to \mathbf{P}^n$，使得
$\mathcal{L} = \varphi^*\mathcal{O}_{\mathbf{P}^n}(1)$，则按构造
$h_\mathcal{L}(P) = h_n(\varphi(P))$，并利用这一点。
一般地，将 $\mathcal{L}$ 写成两个此类模之差。）
\item 假设 $h_\mathcal{L}(P) - h_\mathcal{L}(\alpha(P))$
在 $X(K)$ 上无界。证明对于
$\mathcal{N} = \mathcal{L} \otimes \alpha^*\mathcal{L}^{\otimes -1}$，
$h_\mathcal{N}$ 在 $X(K)$ 上无界。
\item 假设 $X$ 是椭圆曲线，且 $\mathcal{L}$ 是 $X$ 上对称的充足可逆模，
使得 $h_\mathcal{L}$ 在 $X(K)$ 上无界。证明存在一个次数为 $0$ 的可逆模
$\mathcal{N}$，使 $h_\mathcal{N}$ 无界。
（提示：回顾 $X$ 是维数为 $1$ 的阿贝尔簇。因此由课堂结果，
$h_\mathcal{L}$ 除一个常数外是二次的。选取适当的点
$P_0 \in X(K)$。令 $\alpha : X \to X$ 为按 $P_0$ 平移。
考虑 $P \mapsto h_\mathcal{L}(P) - h_\mathcal{L}(P + P_0)$。
应用你在上面证明的结果。）
\end{enumerate}
\end{exercise}

\begin{exercise}[单态射]
\label{exercises-exercise-monomorphisms}
设 $f : X \to Y$ 是概形范畴中的单态射：对于任意一对概形态射
$a, b : T \to X$，若 $f \circ a = f \circ b$，则 $a = b$。
证明 $f$ 在点上是单射。% P12-ERR-0504
你的论证还说明了其他什么吗？
\end{exercise}

\begin{exercise}[不动点]
\label{exercises-exercise-fix-points}
设 $k$ 为代数闭域。
\begin{enumerate}
\item 若 $G = \mathbf{G}_{m, k}$，证明当 $G$ 作用于 $k$ 上的射影代数簇
$X$ 时，该作用有不动点；即证明存在 $x \in X(k)$，使得对所有
$g \in G(k)$ 都有 $a(g, x) = x$。
\item 对 $G = (\mathbf{G}_{m, k})^n$ 作同样证明，其中它等于
$n \geq 1$ 个乘法群的直积。
\item 给出一个连通群概形 $G$ 作用于光滑射影代数簇 $X$ 的例子，
使该作用没有不动点。
\end{enumerate}
\end{exercise}

\section{代数几何，2025 年春季期末考试}
\label{exercises-section-final-exam-spring-2025}

\noindent
以下是哥伦比亚大学 2025 年春季代数几何课程期末考试的试题。

\begin{exercise}[定义]
\label{exercises-exercise-definitions-spring-2025}
简要定义下列斜体概念。
\begin{enumerate}
\item 一个 {\it 概形} $X$，
\item 概形 $X$ 上的一个 {\it 拟凝聚模}，
\item 概形 $X$ 的 {\it Picard 群}，
% P12-ERR-0505
\item 给定概形态射 $f : X \to Y$ 与 $\mathcal{O}_X$-模
$\mathcal{F}$，其 {\it 推前} $f_*\mathcal{F}$，
\item 概形的一个 {\it 闭浸入}，以及
\item 给定域 $k$ 上的一个 {\it 代数簇}。
\end{enumerate}
\end{exercise}

\begin{exercise}[定理]
\label{exercises-exercise-results-spring-2025}
% P12-ERR-0506
准确而简洁地陈述一个课堂上讨论过的、与下列每项相关的非平凡事实
（若有多个，只需任选一个）。
\begin{enumerate}
\item 概形 $X$ 的拓扑空间，
\item 概形范畴上的函子的一个可表性判据，
\item 代数闭域 $k$ 上光滑射影曲线 $X$ 的 Picard 函子，
\item 代数闭域 $k$ 上光滑射影曲线 $X$ 的 Picard 群中的挠元，
\item 一个关于光滑射影代数簇 $X$、$Y$、$Z$ 的乘积
$X \times Y \times Z$ 的 Picard 群的定理，其中这些代数簇定义在
代数闭域 $k$ 上。
\end{enumerate}
\end{exercise}

\begin{exercise}
\label{exercises-exercise-affine-line-infinite-nr-points}
设 $k$ 为域。解释为什么 $\Spec(k[x])$ 有无穷多个点。
\end{exercise}

\begin{exercise}
\label{exercises-exercise-hypersurface-variety}
设 $k$ 为域，$x_1, \ldots, x_n$ 为变量。
令 $f \in k[x_1, \ldots, x_n]$ 为非常数多项式。
令 $X = \Spec(k[x_1, \ldots, x_n]/(f))$。为了使 $X$ 成为
$k$ 上的代数簇，$f$ 应具有什么性质？
\end{exercise}

\begin{exercise}
\label{exercises-exercise-find-singular-point}
找出 $\mathbf{C}[x, y]/(x^5 - 5xy^4 + 20y - 16)$ 的谱上
一个奇点（无需找出全部奇点），并将其视为 $\mathbf{C}$ 上的代数簇。
\end{exercise}

\begin{exercise}
\label{exercises-exercise-kernel-restriction-pic-curve}
设 $X$ 为代数闭域 $k$ 上的光滑射影曲线。令 $x \in X(k)$
为 $X$ 上的闭点。令 $U = X \setminus \{x\}$。
关于限制映射 $\Pic(X) \to \Pic(U)$ 的核，你能说些什么？
\end{exercise}

\begin{exercise}
\label{exercises-exercise-sections-degree-2g-minus-4}
令 $X$ 为代数闭域 $k$ 上亏格 $g \geq 100$ 的光滑射影曲线。
% P12-ERR-0507
对于次数为 $2g - 4$ 的 $X$ 上可逆 $\mathcal{O}_X$-模
$\mathcal{L}$，$h^0(\mathcal{L}) = \dim_k H^0(X, \mathcal{L})$
可以取哪些值？请尽可能完整地回答。
\end{exercise}
